% Produit vectoriel — Mathématiques Tle C — Cours signé OnBuch+
% Programme officiel MINESEC. Compiler avec : tectonic M17-produit-vectoriel.tex
\newcommand{\DOCMATIERE}{Mathématiques}
\newcommand{\DOCNIVEAU}{Tle C}
\newcommand{\DOCMODULE}{Géométrie dans l'espace}
\newcommand{\DOCLECON}{17}
\newcommand{\DOCTITRE}{Produit vectoriel}
% OnBuch+ — préambule « cours signés » ENRICHI (compilé avec tectonic / XeLaTeX)
% Système visuel « Pop » d'OnBuch+ : fond crème (jamais blanc), bordures encre
% épaisses, ombres « dures » décalées, titres Archivo Black en capitales.
% Version MATHÉMATIQUES : graphes (pgfplots/tikz), boîtes pédagogiques
% (définition, propriété, méthode, attention, le savais-tu, exercice résolu,
% exercice, TP, bilan), page de garde et signature OnBuch+ ancrée en bas.
%
% Macros attendues dans main.tex AVANT \input{preamble} :
%   \DOCMATIERE  \DOCNIVEAU  \DOCMODULE  \DOCLECON (numéro)  \DOCTITRE
\documentclass[11pt]{article}

\usepackage{fontspec}
\usepackage[french]{babel}
\usepackage[a4paper,margin=2cm,top=3.4cm,bottom=2.8cm,headheight=1.4cm,footskip=1.3cm]{geometry}
\usepackage{amsmath,amssymb,amsfonts}
\usepackage{mathtools} % \xmapsto, \coloneqq... souvent utilisés par les modèles
\usepackage{cancel} % simplifications barrées
% \abs{z} (module, valeur absolue) et \norm{u} (norme), utilisés par les modèles
\providecommand{\abs}{}\let\abs\relax\DeclarePairedDelimiter{\abs}{\lvert}{\rvert}
\providecommand{\norm}{}\let\norm\relax\DeclarePairedDelimiter{\norm}{\lVert}{\rVert}
\usepackage[table]{xcolor} % [table] : fournit aussi \rowcolors (alternance de lignes)
\usepackage{pagecolor}
\usepackage{tikz}
\usetikzlibrary{arrows.meta,positioning,calc,shapes.geometric,shapes.misc,shapes.arrows,decorations.pathmorphing,decorations.markings,patterns,shadows,mindmap,trees,fit,backgrounds,matrix,intersections}
\usepackage{pgfplots}
\usepackage{pgf-pie} % diagrammes circulaires \pie (statistiques)
\pgfplotsset{compat=1.18}
\usepgfplotslibrary{fillbetween,groupplots} % aires entre courbes (calcul intégral)
\usepackage{enumitem}
\usepackage{fancyhdr}
% [most] charge toutes les bibliothèques tcolorbox (listings, minted, raster,
% documentation...) et gonfle inutilement la mémoire TeX sur un document long
% (~300 boîtes) : on ne charge que ce qui est réellement utilisé.
\usepackage[skins,breakable]{tcolorbox}
\usepackage{graphicx}
\usepackage{colortbl}
\usepackage{booktabs}
\usepackage{tabularx}
\usepackage{diagbox} % case d'en-tête diagonale des tableaux à double entrée
% Colonnes extensibles centrée / gauche / droite (C, L, R), souvent utilisées par les modèles
\newcolumntype{C}{>{\centering\arraybackslash}X}
\newcolumntype{L}{>{\raggedright\arraybackslash}X}
\newcolumntype{R}{>{\raggedleft\arraybackslash}X}
\usepackage{array}
\usepackage{multicol}
\usepackage{siunitx}
% parse-numbers=false : accepte aussi \SI{2,5 \times 10^{-3}}{...}
\sisetup{locale=FR,output-decimal-marker={,},group-minimum-digits=4,parse-numbers=false}
\DeclareSIUnit\ohm{\ensuremath{\symup{\Omega}}} % U+2126 absent de la police
\usepackage{qrcode}
% \degree : commande fréquemment utilisée par les modèles pour un angle ou une
% température en dehors de \SI{}{} (non fournie par les paquets chargés).
\providecommand{\degree}{^{\circ}}
% \qed (fin de démonstration), utilisé par les modèles sans amsthm
\providecommand{\qed}{\hfill$\square$}
% \barre : double barre des valeurs interdites dans un tableau de variations (tkz-tab)
\providecommand{\barre}{\ensuremath{\|}}
% environnement proof (amsthm non chargé) utilisé par les modèles
\ifdefined\proof\else\newenvironment{proof}[1][Démonstration]{\par\noindent\textit{#1.}\ }{\qed\par}\fi
\usepackage{lastpage}
\usepackage{hyperref}
\hypersetup{hidelinks}

% ============================================================
% Palette « Pop » OnBuch+ — mêmes hex que la classe P (plus_ui.dart)
% ============================================================
\definecolor{poporange}{HTML}{F59321}
\definecolor{popdark}{HTML}{E07A0C}
\definecolor{popcream}{HTML}{FFF6E8}
\definecolor{popink}{HTML}{1C1714}
\definecolor{poppurple}{HTML}{7A5AE0}
\definecolor{popgreen}{HTML}{2E9E6A}
\definecolor{poppink}{HTML}{E0457B}
\definecolor{popblue}{HTML}{2F7DE1}
\definecolor{popgold}{HTML}{C98A12}
\definecolor{popmuted}{HTML}{8A7F76}
\definecolor{popline}{HTML}{EADFD2}
\definecolor{popgray}{HTML}{8A7F76}
\colorlet{popgrey}{popgray}
% Alias tolérants (le modèle nomme parfois les couleurs différemment)
\colorlet{popwhite}{white}
\colorlet{popblack}{popink}
\colorlet{popred}{poppink}
\colorlet{popyellow}{popgold}
% Teintes claires pour fonds de tableaux / zones
\colorlet{popblueL}{popblue!10!white}
\colorlet{poporangeL}{poporange!14!white}
\colorlet{popgreenL}{popgreen!12!white}
\colorlet{poppurpleL}{poppurple!12!white}
\colorlet{poppinkL}{poppink!12!white}
\colorlet{popgoldL}{popgold!14!white}

\pagecolor{popcream}

% ============================================================
% Typographie
% ============================================================
\defaultfontfeatures{Ligatures=TeX}
\setmainfont{PlusJakartaSans-Medium.ttf}[Path=../../../fonts/,BoldFont=PlusJakartaSans-Bold.ttf]
% Maths en sans-serif (Fira Math) avec chiffres et lettres droites de Plus
% Jakarta, pour que les formules se fondent dans le texte.
\usepackage[math-style=ISO,bold-style=ISO]{unicode-math}
\setmathfont{FiraMath-Regular.otf}[Path=../../../fonts/]
\setmathfont{PlusJakartaSans-Medium.ttf}[Path=../../../fonts/,range={up/num,up/{Latin,latin},\mathup}]
\usepackage{icomma} % virgule décimale sans espace en mode math
% Caractères mathématiques Unicode absents des polices de texte (Plus Jakarta,
% Archivo Black) : rendus par leur équivalent mathématique, en texte comme en
% formule — sinon ils s'affichent en carré vide (ex. « Arithmétique dans ℤ »).
\usepackage{newunicodechar}
\newunicodechar{ℝ}{\ensuremath{\mathbb{R}}}
\newunicodechar{ℤ}{\ensuremath{\mathbb{Z}}}
\newunicodechar{ℂ}{\ensuremath{\mathbb{C}}}
\newunicodechar{ℕ}{\ensuremath{\mathbb{N}}}
\newunicodechar{ℚ}{\ensuremath{\mathbb{Q}}}
\newunicodechar{𝔻}{\ensuremath{\mathbb{D}}}
\newunicodechar{α}{\ensuremath{\alpha}}
\newunicodechar{β}{\ensuremath{\beta}}
\newunicodechar{γ}{\ensuremath{\gamma}}
\newunicodechar{δ}{\ensuremath{\delta}}
\newunicodechar{ε}{\ensuremath{\varepsilon}}
\newunicodechar{θ}{\ensuremath{\theta}}
\newunicodechar{λ}{\ensuremath{\lambda}}
\newunicodechar{μ}{\ensuremath{\mu}}
\newunicodechar{σ}{\ensuremath{\sigma}}
\newunicodechar{τ}{\ensuremath{\tau}}
\newunicodechar{φ}{\ensuremath{\varphi}}
\newunicodechar{ψ}{\ensuremath{\psi}}
\newunicodechar{ω}{\ensuremath{\omega}}
\newunicodechar{Σ}{\ensuremath{\Sigma}}
% majuscules produites par \MakeUppercase dans les titres (α-aminés -> Α)
\newunicodechar{Α}{\ensuremath{\alpha}}
\newunicodechar{Β}{\ensuremath{\beta}}
\newunicodechar{Γ}{\ensuremath{\Gamma}}
\newunicodechar{Δ}{\ensuremath{\Delta}}
\newunicodechar{Ε}{\ensuremath{\varepsilon}}
\newunicodechar{Θ}{\ensuremath{\Theta}}
\newunicodechar{Λ}{\ensuremath{\Lambda}}
\newunicodechar{Μ}{\ensuremath{\mu}}
\newunicodechar{Τ}{\ensuremath{\tau}}
\newunicodechar{Φ}{\ensuremath{\Phi}}
\newunicodechar{Ψ}{\ensuremath{\Psi}}
\newunicodechar{Ω}{\ensuremath{\Omega}}
\newunicodechar{↦}{\ensuremath{\mapsto}}
\newunicodechar{∈}{\ensuremath{\in}}
\newunicodechar{∉}{\ensuremath{\notin}}
\newunicodechar{∧}{\ensuremath{\wedge}}
\newunicodechar{⇔}{\ensuremath{\Leftrightarrow}}
\newunicodechar{⟺}{\ensuremath{\Longleftrightarrow}}
\newunicodechar{⇒}{\ensuremath{\Rightarrow}}
\newunicodechar{⟹}{\ensuremath{\Longrightarrow}}
\newunicodechar{∘}{\ensuremath{\circ}}
\newunicodechar{∪}{\ensuremath{\cup}}
\newunicodechar{∩}{\ensuremath{\cap}}
\newunicodechar{≡}{\ensuremath{\equiv}}
\newunicodechar{⊂}{\ensuremath{\subset}}
\newunicodechar{⊥}{\ensuremath{\perp}}
\newunicodechar{∃}{\ensuremath{\exists}}
\newunicodechar{∀}{\ensuremath{\forall}}
\newunicodechar{∛}{\ensuremath{\sqrt[3]{\,}}}
\newunicodechar{∜}{\ensuremath{\sqrt[4]{\,}}}
\newunicodechar{⃗}{\ensuremath{{}^{\to}}}
\newunicodechar{ⁿ}{\ifmmode^{n}\else\textsuperscript{\ensuremath{n}}\fi}
\newunicodechar{ˣ}{\ifmmode^{x}\else\textsuperscript{\ensuremath{x}}\fi}
\newunicodechar{ᵇ}{\ifmmode^{b}\else\textsuperscript{\ensuremath{b}}\fi}
\newunicodechar{ʳ}{\ifmmode^{r}\else\textsuperscript{\ensuremath{r}}\fi}
\newunicodechar{ᵘ}{\ifmmode^{u}\else\textsuperscript{\ensuremath{u}}\fi}
\newunicodechar{ʸ}{\ifmmode^{y}\else\textsuperscript{\ensuremath{y}}\fi}
\newunicodechar{ᵃ}{\ifmmode^{a}\else\textsuperscript{\ensuremath{a}}\fi}
\newunicodechar{ᵉ}{\ifmmode^{e}\else\textsuperscript{\ensuremath{e}}\fi}
\newunicodechar{ᵏ}{\ifmmode^{k}\else\textsuperscript{\ensuremath{k}}\fi}
\newunicodechar{ᵖ}{\ifmmode^{p}\else\textsuperscript{\ensuremath{p}}\fi}
\newunicodechar{ᴺ}{\ifmmode^{N}\else\textsuperscript{\ensuremath{N}}\fi}
\newunicodechar{ᶜ}{\ifmmode^{c}\else\textsuperscript{\ensuremath{c}}\fi}
\newunicodechar{ʰ}{\ifmmode^{h}\else\textsuperscript{\ensuremath{h}}\fi}
\newunicodechar{ᵠ}{\ifmmode^{q}\else\textsuperscript{\ensuremath{q}}\fi}
\newunicodechar{ₙ}{\ifmmode_{n}\else\textsubscript{\ensuremath{n}}\fi}
\newunicodechar{ₐ}{\ifmmode_{a}\else\textsubscript{\ensuremath{a}}\fi}
\newunicodechar{ᵢ}{\ifmmode_{i}\else\textsubscript{\ensuremath{i}}\fi}
\newunicodechar{ᵣ}{\ifmmode_{r}\else\textsubscript{\ensuremath{r}}\fi}
\newunicodechar{ₖ}{\ifmmode_{k}\else\textsubscript{\ensuremath{k}}\fi}
\newunicodechar{ᵦ}{\ifmmode_{\beta}\else\textsubscript{\ensuremath{\beta}}\fi}
\newunicodechar{ₓ}{\ifmmode_{x}\else\textsubscript{\ensuremath{x}}\fi}
\newfontfamily\popdisplay{ArchivoBlack-Regular.ttf}[Path=../../../fonts/]
\newfontfamily\poplogo{SpaceGrotesk-Bold.ttf}[Path=../../../fonts/]
\newfontfamily\popsemi{PlusJakartaSans-SemiBold.ttf}[Path=../../../fonts/]
\newfontfamily\popxbold{PlusJakartaSans-ExtraBold.ttf}[Path=../../../fonts/]
\newfontfamily\popxboldls{PlusJakartaSans-ExtraBold.ttf}[Path=../../../fonts/,LetterSpace=3.0]

\setlist[enumerate]{leftmargin=1.7em,itemsep=5pt,topsep=4pt}
\setlist[itemize]{leftmargin=1.7em,itemsep=4pt,topsep=4pt,label={\color{poporange}\textbullet}}
\setlength{\parindent}{0pt}
\setlength{\parskip}{5pt}
\renewcommand{\arraystretch}{1.25}

% pgfplots : style Pop par défaut
\pgfplotsset{
  every axis/.append style={axis line style={popink,line width=1pt},tick style={popink},
    grid style={popline,line width=0.5pt},label style={font=\small},tick label style={font=\footnotesize}},
  popcurve/.style={poporange,line width=1.8pt,no markers,smooth},
}
% Même style disponible dans un \draw TikZ simple (hors pgfplots)
\tikzset{popcurve/.style={poporange,line width=1.8pt}}

% ============================================================
% Pill / squiggle
% ============================================================
\newcommand{\poppill}[3]{%
  \tikz[baseline=-0.55ex]\node[draw=popink,line width=1.2pt,rounded corners=2.6pt,%
    fill=#2,inner xsep=6.5pt,inner ysep=3pt]{%
    \popxboldls\fontsize{7}{7}\selectfont\color{#3}\MakeUppercase{#1}};%
}
\newcommand{\popsquiggle}[1]{%
  \tikz[baseline=-0.4ex]{
    \draw[#1,line width=2.6pt,line cap=round]
      plot [smooth] coordinates {(0,0.09) (0.34,0.01) (0.68,0.21) (1.02,0.01) (1.36,0.10)};
  }%
}
% Mot-clé surligné (marqueur orange sous le texte)
\newcommand{\cle}[1]{{\popxbold\color{popdark}#1}}

% ============================================================
% En-tête / pied
% ============================================================
\pagestyle{fancy}
\fancyhf{}
\renewcommand{\headrulewidth}{0pt}
\renewcommand{\footrulewidth}{0pt}
\lhead{\raisebox{-0.15ex}{\poplogo\fontsize{13}{13}\selectfont\color{popink}OnBuch\color{poporange}+}}
\rhead{\poppill{\DOCMATIERE\ \textbullet\ \DOCNIVEAU\ \textbullet\ Leçon \DOCLECON}{white}{popink}}
\lfoot{\popxbold\fontsize{7.6}{7.6}\selectfont\color{popmuted}\MakeUppercase{Cours signé OnBuch+}\ \textbullet\ {\popsemi\DOCTITRE}}
\rfoot{\tikz[baseline=-0.6ex]\node[rounded corners=8.5pt,fill=popink,minimum height=17pt,minimum width=17pt,inner xsep=5pt,inner sep=0]{\popxbold\fontsize{8}{8}\selectfont\color{popcream}\thepage\,/\,\pageref*{LastPage}};}

% ============================================================
% Titres
% ============================================================
\newcommand{\chaptitre}[2]{%
  \begin{tcolorbox}[enhanced,colback=poporange,colframe=popink,boxrule=2.6pt,arc=11pt,
      drop shadow=popink,left=15pt,right=15pt,top=12pt,bottom=11pt,before skip=3pt,after skip=18pt]
    \poppill{#2}{popink}{popcream}\par\vspace{9pt}
    {\raggedright\hyphenpenalty=10000\exhyphenpenalty=10000\popdisplay\fontsize{22}{24}\selectfont\color{popcream}\MakeUppercase{#1}\par}
  \end{tcolorbox}
}
% Section numérotée : pastille orange avec le numéro + titre Archivo
\newcounter{popsec}
\newcommand{\coursec}[1]{%
  \stepcounter{popsec}\setcounter{popsub}{0}%
  \par\vspace{16pt}\noindent
  \begin{minipage}{\linewidth}
  \tikz[baseline=-0.7ex]\node[draw=popink,line width=1.6pt,fill=poporange,rounded corners=4pt,minimum size=22pt,inner sep=0]{\popdisplay\fontsize{12}{12}\selectfont\color{popcream}\thepopsec};%
  \hspace{8pt}{\popdisplay\fontsize{15}{17}\selectfont\color{popink}\MakeUppercase{#1}}\par
  \vspace{4pt}\noindent{\color{poporange}\rule{36pt}{3.6pt}}
  \end{minipage}\par\nopagebreak\vspace{8pt}
}
\newcounter{popsub}
\newcommand{\courssub}[1]{%
  \stepcounter{popsub}%
  \par\vspace{9pt}\noindent{\popxbold\fontsize{11.5}{13}\selectfont\color{popink}{\color{poporange}\thepopsec.\thepopsub}\hspace{6pt}#1}\par\nopagebreak\vspace{4pt}
}

% ============================================================
% Boîtes pédagogiques — même recette : fond blanc, bordure épaisse colorée,
% ombre dure encre, badge plein. Titre optionnel [#1].
% ============================================================
\tcbset{popbase/.style={breakable,enhanced,colback=white,boxrule=2.3pt,arc=9pt,
  shadow={3pt}{-3pt}{0pt}{popink},left=14pt,right=14pt,top=11pt,bottom=11pt,before skip=11pt,after skip=15pt,
  fontupper=\popsemi\fontsize{10.6}{14.5}\selectfont\color{popink},
  fontlower=\popsemi\fontsize{10.6}{14.5}\selectfont\color{popink}}}
% Coche dessinée (le glyphe ✓ n'existe pas dans les polices du document)
\newcommand{\popcheck}[1]{\tikz[baseline=-0.5ex]\draw[#1,line width=1.6pt,line cap=round,line join=round](-0.13,0.0)--(-0.04,-0.09)--(0.13,0.1);}
\newcommand{\popboxhead}[3]{\poppill{#1}{#2}{white}\ifx\relax#3\relax\else\hspace{6pt}{\popxbold\fontsize{10.6}{12}\selectfont\color{popink}#3}\fi\par\vspace{7pt}}

\newtcolorbox{aretenir}[1][]{popbase,colframe=popgreen,before upper={\popboxhead{À retenir}{popgreen}{#1}}}
\newtcolorbox{exemplebox}[1][]{popbase,colframe=poporange,before upper={\popboxhead{Exemple}{poporange}{#1}}}
\newtcolorbox{definition}[1][]{popbase,colframe=popblue,before upper={\popboxhead{Définition}{popblue}{#1}}}
\newtcolorbox{propriete}[1][]{popbase,colframe=poppurple,before upper={\popboxhead{Propriété}{poppurple}{#1}}}
\newtcolorbox{methode}[1][]{popbase,colframe=popgold,before upper={\popboxhead{Méthode}{popgold}{#1}}}
\newtcolorbox{attention}[1][]{popbase,colframe=poppink,before upper={\popboxhead{Attention}{poppink}{#1}}}
\newtcolorbox{experience}[1][]{popbase,colframe=popblue,colback=popblueL,before upper={\popboxhead{Expérience}{popblue}{#1}}}
% « Le savais-tu ? » : carte violette pleine (contexte, culture, Cameroun)
\newtcolorbox{savaistu}[1][]{popbase,colback=poppurple,colframe=popink,
  fontupper=\popsemi\fontsize{10.4}{14.2}\selectfont\color{white},
  before upper={\poppill{Le savais-tu\,?}{white}{poppurple}\ifx\relax#1\relax\else\hspace{6pt}{\popxbold\color{white}#1}\fi\par\vspace{7pt}}}
% Exercice résolu : énoncé en haut, \tcblower puis la solution sur fond crème
\newcounter{popexo}\newcounter{popexr}
\newtcolorbox{exoresolu}[1][]{popbase,colframe=poporange,colbacklower=popcream,
  segmentation style={popink,line width=1.2pt,dashed},
  before upper={\stepcounter{popexr}\popboxhead{Exercice résolu \thepopexr}{poporange}{#1}},
  before lower={\poppill{Solution}{popink}{popcream}\par\vspace{6pt}}}
% Exercice (non résolu en place) — la correction est dans la partie Corrigés
\newtcolorbox{exercice}[1][]{popbase,colframe=popink,
  before upper={\stepcounter{popexo}\popboxhead{Exercice \thepopexo}{popink}{#1}}}
% Corrigé
\newtcolorbox{corrige}[1][]{popbase,colframe=popgreen,colback=popgreenL,
  before upper={\popboxhead{Corrigé}{popgreen}{#1}}}
% Objectifs / prérequis / bilan
\newtcolorbox{objectifs}{popbase,colframe=popink,colback=poporangeL,
  before upper={\popboxhead{Objectifs de la leçon}{popink}{}}}
\newtcolorbox{prerequis}{popbase,colframe=popink,colback=popblueL,
  before upper={\popboxhead{Prérequis}{popblue}{}}}
\newtcolorbox{bilan}{popbase,colframe=popink,colback=white,boxrule=2.8pt,
  before upper={\popboxhead{Fiche bilan}{poporange}{L'essentiel à connaître pour l'examen}}}
% Figure encadrée avec légende
\newtcolorbox{popfigure}[1][]{enhanced,breakable=false,colback=white,colframe=popline,boxrule=1.4pt,arc=8pt,
  left=10pt,right=10pt,top=10pt,bottom=8pt,before skip=10pt,after skip=12pt,halign=center,
  lower separated=false,halign lower=center,
  fontlower=\popsemi\fontsize{9}{11.5}\selectfont\color{popmuted},
  #1}
\newcommand{\legende}[1]{\par\vspace{6pt}{\popsemi\fontsize{9}{11.5}\selectfont\color{popmuted}#1}}

% ============================================================
% Signature OnBuch+
% ============================================================
\newcommand{\poplogoinline}[1]{{\poplogo\fontsize{#1}{#1}\selectfont\color{popink}OnBuch\color{poporange}+}}
\newcommand{\signatureonbuch}{%
  \begin{tcolorbox}[enhanced,colback=popink,colframe=popink,boxrule=2.2pt,arc=10pt,
      drop shadow=poporange,left=14pt,right=14pt,top=10pt,bottom=10pt,before skip=0pt,after skip=0pt]
    \tikz[baseline=-0.7ex]\node[circle,fill=popgreen,draw=popcream,line width=1.3pt,minimum size=22pt,inner sep=0]{\popcheck{white}};%
    \hspace{9pt}\begin{minipage}[c]{0.62\linewidth}
      {\popxbold\fontsize{12}{13}\selectfont\color{popcream}Signé\ \textbullet\ {\poplogo OnBuch\color{poporange}+}}\par\vspace{2pt}
      {\raggedright\popsemi\fontsize{8}{10.5}\selectfont\color{popline}\MakeUppercase{Équipe pédagogique OnBuch+}\\\MakeUppercase{Conforme au programme officiel MINESEC}\par}
    \end{minipage}\hfill
    {\poplogo\fontsize{22}{22}\selectfont\color{popcream}OnBuch\color{poporange}+}
  \end{tcolorbox}%
}

% ============================================================
% Page de garde
% #1 titre · #2 matière · #3 niveau · #4 module · #5 n° leçon · #6 durée ·
% #7 description / objectifs (texte ou itemize)
% ============================================================
\newcommand{\pagegarde}[7]{%
  \thispagestyle{empty}
  \newgeometry{a4paper,margin=1.7cm,top=1.6cm,bottom=1.4cm}%
  \begin{tikzpicture}[remember picture,overlay]
    \fill[poporange!18] ([xshift=-3.2cm,yshift=-3.6cm]current page.north east) circle (4.6cm);
    \fill[poppurple!14] ([xshift=2.4cm,yshift=7.6cm]current page.south west) circle (3.8cm);
  \end{tikzpicture}%
  {\poplogo\fontsize{30}{30}\selectfont\color{popink}OnBuch\color{poporange}+}\hfill
  \raisebox{0.6ex}{\poppill{Cours signé \textbullet\ Premium}{popink}{popcream}}\par
  \vspace{1pt}\popsquiggle{poppurple}\par
  \vspace{8pt}
  \begin{tcolorbox}[enhanced,colback=poporange,colframe=popink,boxrule=2.8pt,arc=13pt,
      drop shadow=popink,left=18pt,right=18pt,top=12pt,bottom=13pt,after skip=10pt]
    \poppill{#2 \textbullet\ #3}{popink}{popcream}\hspace{5pt}\poppill{#4}{white}{popink}\par\vspace{8pt}
    {\popdisplay\fontsize{14}{14}\selectfont\color{popink}LEÇON #5}\par\vspace{4pt}
    {\raggedright\hyphenpenalty=10000\exhyphenpenalty=10000\popdisplay\fontsize{25}{27}\selectfont\color{popcream}\MakeUppercase{#1}\par}
    \vspace{8pt}
    {\popxbold\fontsize{9.5}{9.5}\selectfont\color{popink}\MakeUppercase{Durée conseillée : #6}}
  \end{tcolorbox}
  \begin{tcolorbox}[enhanced,colback=white,colframe=popink,boxrule=2.2pt,arc=10pt,
      drop shadow=popink,left=16pt,right=16pt,top=10pt,bottom=10pt,after skip=8pt]
    \poppill{Dans cette leçon}{poppurple}{white}\par\vspace{5pt}
    \popsemi\fontsize{9.3}{12.6}\selectfont\color{popink}#7
  \end{tcolorbox}
  \noindent\begin{minipage}[t]{0.487\linewidth}
    \begin{tcolorbox}[enhanced,colback=popgreen,colframe=popink,boxrule=2.2pt,arc=10pt,
        drop shadow=popink,left=11pt,right=11pt,top=10pt,bottom=10pt,height=3.1cm,valign=center]
      \tcbox[colback=white,colframe=popink,boxrule=1.3pt,arc=5pt,boxsep=0pt,nobeforeafter,
        left=3pt,right=3pt,top=3pt,bottom=3pt,valign=center]{\qrcode[height=1.9cm]{https://play.google.com/store/apps/details?id=com.luvvix.onbuch&hl=fr}}%
      \hspace{8pt}\begin{minipage}[c]{0.5\linewidth}
        {\popdisplay\fontsize{11}{12.5}\selectfont\color{white}\MakeUppercase{Télécharge OnBuch}}\par\vspace{4pt}
        {\raggedright\popsemi\fontsize{8.4}{11}\selectfont\color{white}Gratuit \textbullet\ Google Play\\Tuteur IA, annales, concours, résultats\par}
      \end{minipage}
    \end{tcolorbox}
  \end{minipage}\hfill
  \begin{minipage}[t]{0.487\linewidth}
    \begin{tcolorbox}[enhanced,colback=poppurple,colframe=popink,boxrule=2.2pt,arc=10pt,
        drop shadow=popink,left=11pt,right=11pt,top=10pt,bottom=10pt,height=3.1cm,valign=center]
      \tikz[baseline=-0.7ex]\node[circle,fill=white,draw=popink,line width=1.3pt,minimum size=1.6cm,inner sep=0]{\popdisplay\fontsize{24}{24}\selectfont\color{poppurple}+};%
      \hspace{8pt}\begin{minipage}[c]{0.6\linewidth}
        {\popdisplay\fontsize{11}{12.5}\selectfont\color{white}\MakeUppercase{Passe à OnBuch+}}\par\vspace{4pt}
        {\raggedright\popsemi\fontsize{8.4}{11}\selectfont\color{white}Tous les cours signés, Tuteur IA illimité, fiches PDF — onglet OnBuch+ de l'app\par}
      \end{minipage}
    \end{tcolorbox}
  \end{minipage}
  \vspace{8pt}
  \popcontenu
  \vfill
  \signatureonbuch
  \restoregeometry
  \newpage
}

% Rangée « Ce cours contient » (page de garde)
\newcommand{\poptile}[3]{% #1 couleur #2 gros texte #3 légende
  \begin{tcolorbox}[enhanced,colback=#1,colframe=popink,boxrule=2pt,arc=9pt,shadow={2.5pt}{-2.5pt}{0pt}{popink},
      left=6pt,right=6pt,top=9pt,bottom=9pt,halign=center,valign=center,height=1.75cm,width=\linewidth,nobeforeafter]
    {\popdisplay\fontsize{12.5}{13}\selectfont\color{white}\MakeUppercase{#2}}\par\vspace{3pt}
    {\centering\popsemi\fontsize{7.8}{9.5}\selectfont\color{white}#3\par}
  \end{tcolorbox}}
\newcommand{\popcontenu}{%
  \par\noindent{\popxboldls\fontsize{7.5}{7.5}\selectfont\color{popmuted}CE COURS SIGNÉ CONTIENT}\par\vspace{7pt}
  \noindent\begin{minipage}[t]{0.235\linewidth}\poptile{popblue}{Cours}{complet, illustré, pas à pas}\end{minipage}\hfill
  \begin{minipage}[t]{0.235\linewidth}\poptile{poporange}{Méthodes}{et exercices résolus}\end{minipage}\hfill
  \begin{minipage}[t]{0.235\linewidth}\poptile{poppink}{Exercices}{type examen + corrigés}\end{minipage}\hfill
  \begin{minipage}[t]{0.235\linewidth}\poptile{popgreen}{Fiche bilan}{l'essentiel à retenir}\end{minipage}\par}

% Sommaire
\newtcolorbox{sommaire}{enhanced,colback=white,colframe=popink,boxrule=2.2pt,arc=10pt,
  drop shadow=popink,left=18pt,right=18pt,top=13pt,bottom=13pt,before skip=6pt,after skip=20pt,
  before upper={\poppill{Sommaire}{poppurple}{white}\par\vspace{9pt}\popsemi\fontsize{10.6}{14.5}\selectfont\color{popink}}}

% Signature de fin de cours — poussée tout en bas de la dernière page
\newcommand{\findecours}{%
  \par\vfill\nopagebreak
  \begin{center}\popsquiggle{poporange}\end{center}\vspace{4pt}
  \signatureonbuch
}
\AtBeginDocument{\def\llbracket{\mathopen{[\mkern-3.5mu[}}\def\rrbracket{\mathclose{]\mkern-3.5mu]}}} % intervalles d'entiers (glyphe ⟦ absent de la police)
% Glyphes absents de Fira Math : redéfinis à partir de glyphes disponibles
\AtBeginDocument{%
  \def\setminus{\mathbin{\backslash}}\let\smallsetminus\setminus
  \def\vdots{\mathord{\vbox{\baselineskip3.2pt\lineskiplimit0pt\kern4pt\hbox{.}\hbox{.}\hbox{.}}}}%
  \def\ddots{\mathinner{\mkern1mu\raise6.5pt\hbox{.}\mkern2mu\raise3.6pt\hbox{.}\mkern2mu\raise0.7pt\hbox{.}\mkern1mu}}%
  \def\triangle{\mathord{\scalebox{0.9}{$\symup{\Delta}$}}}%
  \def\nabla{\mathord{\raisebox{\depth}{\scalebox{1}[-1]{$\symup{\Delta}$}}}}%
  \def\checkmark{\popcheck{popdark}}%
  \def\star{\mathbin{\ast}}\def\bigstar{\mathord{\ast}}%
  \def\diamond{\mathbin{\scalebox{0.6}{$\Diamond$}}}%
}

%% FIGFIX-BEGIN v4 — correctifs globaux des figures (docs/onbuch-generation/tools/figfix)
\usepackage{adjustbox}\usepackage{etoolbox}
% toute figure TikZ/pgfplots plus large que la ligne est réduite à la largeur disponible
\BeforeBeginEnvironment{tikzpicture}{\begin{adjustbox}{max width=\linewidth}}
\AfterEndEnvironment{tikzpicture}{\end{adjustbox}}
% légendes pgfplots lisibles par-dessus les courbes
\pgfplotsset{every axis/.append style={legend style={fill=white,fill opacity=0.92,text opacity=1,draw=black!15,inner sep=3pt}}}
% étiquettes d'arêtes (syntaxe edge["..."]) avec fond blanc
\tikzset{every edge quotes/.append style={fill=white,inner sep=1.2pt}}
%% FIGFIX-END
\begin{document}

\pagegarde{Produit vectoriel}{Mathématiques}{Tle C}{Géométrie dans l'espace}{17}{12 h}{Cette leçon introduit le produit vectoriel de deux vecteurs de l'espace, outil fondamental pour construire des vecteurs orthogonaux, calculer des aires, des volumes et des distances. Elle consolide la géométrie analytique de l'espace et prépare efficacement aux épreuves du Baccalauréat C et E.
\vspace{4pt}
\begin{multicols}{2}\small
\begin{itemize}
\item À la fin de cette leçon, je sais orienter l'espace à l'aide de la règle du bonhomme d'Ampère ou de la main droite et reconnaître une base orthonormée directe
\item À la fin de cette leçon, je sais définir le produit vectoriel de deux vecteurs de l'espace et interpréter géométriquement sa norme
\item À la fin de cette leçon, je sais calculer les coordonnées d'un produit vectoriel dans une base orthonormée directe
\item À la fin de cette leçon, je sais utiliser les propriétés (antisymétrie, bilinéarité) pour simplifier des calculs vectoriels
\item À la fin de cette leçon, je sais caractériser la colinéarité de deux vecteurs par la nullité de leur produit vectoriel
\item À la fin de cette leçon, je sais déterminer un vecteur normal à un plan à l'aide du produit vectoriel
\end{itemize}\end{multicols}}

\begin{sommaire}
\begin{multicols}{2}
\begin{enumerate}
\item Orientation de l'espace
\item Définition et interprétation géométrique du produit vectoriel
\item Expression analytique dans une base orthonormée directe
\item Propriétés algébriques du produit vectoriel
\item Distances : point-droite et point-plan
\item Aires, produit mixte et volume du tétraèdre
\item Activité d'intégration
\item Méthodes et exercices résolus
\item Exercices
\item Corrigés des exercices
\item Fiche bilan
\end{enumerate}
\end{multicols}
\end{sommaire}

\begin{objectifs}
\begin{itemize}
\item À la fin de cette leçon, je sais orienter l'espace à l'aide de la règle du bonhomme d'Ampère ou de la main droite et reconnaître une base orthonormée directe
\item À la fin de cette leçon, je sais définir le produit vectoriel de deux vecteurs de l'espace et interpréter géométriquement sa norme
\item À la fin de cette leçon, je sais calculer les coordonnées d'un produit vectoriel dans une base orthonormée directe
\item À la fin de cette leçon, je sais utiliser les propriétés (antisymétrie, bilinéarité) pour simplifier des calculs vectoriels
\item À la fin de cette leçon, je sais caractériser la colinéarité de deux vecteurs par la nullité de leur produit vectoriel
\item À la fin de cette leçon, je sais déterminer un vecteur normal à un plan à l'aide du produit vectoriel
\item À la fin de cette leçon, je sais calculer la distance d'un point à une droite et la distance d'un point à un plan
\item À la fin de cette leçon, je sais calculer l'aire d'un triangle (ou d'un parallélogramme) de l'espace
\item À la fin de cette leçon, je sais calculer le volume d'un tétraèdre à l'aide du produit mixte
\end{itemize}
\end{objectifs}

\begin{prerequis}
\begin{itemize}
\item Vecteurs de l'espace : opérations, colinéarité, repérage (Première C/E)
\item Produit scalaire dans l'espace et orthogonalité de deux vecteurs (leçon précédente de Terminale)
\item Repères orthonormés de l'espace et coordonnées d'un point, d'un vecteur
\item Équation cartésienne d'un plan et vecteur normal (leçon sur le produit scalaire)
\item Notions d'aire d'un triangle et de volume d'un tétraèdre (formule V = (1/3) × aire de base × hauteur)
\end{itemize}
\end{prerequis}

\newpage

\coursec{Orientation de l'espace}

\courssub{Une situation concrète : le panneau solaire de Yaoundé}

À Yaoundé, un technicien doit installer un panneau solaire sur le toit incliné d'une maison. Pour que le panneau soit bien fixé, le support doit être \emph{perpendiculaire aux deux barres métalliques} du cadre. Deux questions se posent alors immédiatement :

\begin{itemize}
\item Comment trouver la direction exacte de cette perpendiculaire commune aux deux barres ?
\item Une fois la direction trouvée, de quel côté faut-il la pointer : vers le haut (pour que le panneau regarde le ciel) ou vers le bas (vers l'intérieur du toit) ?
\end{itemize}

Le menuisier du quartier, lui, veut connaître le volume d'un pilier décoratif en forme de tétraèdre, et le couvreur doit calculer la surface de tôle nécessaire pour une face triangulaire du toit.

Toutes ces questions — perpendiculaire commune, aire d'un triangle, volume d'un tétraèdre — trouvent une réponse élégante et rapide grâce à un nouvel outil : le \cle{produit vectoriel}. Mais attention : le produit vectoriel de deux vecteurs est un \emph{vecteur}, et un vecteur possède un \emph{sens}. Or, dans l'espace, décider ce qui est « vers le haut » ou « vers l'avant » n'est pas automatique : il faut d'abord \cle{orienter l'espace}, c'est-à-dire choisir une convention qui distingue les « bonnes » bases des « mauvaises ». C'est l'objet de cette première section, fondation indispensable de toute la leçon.

\textbf{Problématique.} Comment choisir, de manière cohérente et universelle, le sens du troisième vecteur d'une base de l'espace, de sorte que deux mathématiciens (ou deux techniciens !) obtiennent toujours le même résultat ?

\courssub{Rappel : l'orientation du plan}

Commençons par ce que tu connais déjà depuis la classe de Seconde : l'orientation du plan.

\begin{aretenir}[Sens direct dans le plan]
Dans le plan muni d'un repère orthonormé $(O\,;\vec{i},\vec{j})$, le \cle{sens direct} (ou sens trigonométrique) est le \textbf{sens inverse des aiguilles d'une montre}. On dit que la base $(\vec{i},\vec{j})$ est \emph{directe} lorsque l'angle orienté $(\vec{i},\vec{j})$ mesure $+\dfrac{\pi}{2}$ radians, c'est-à-dire lorsqu'on passe de $\vec{i}$ à $\vec{j}$ par un quart de tour dans le sens inverse des aiguilles d'une montre.
\end{aretenir}

Concrètement : si $\vec{i}$ pointe vers la droite (l'« Est ») et $\vec{j}$ vers le haut (le « Nord »), alors $(\vec{i},\vec{j})$ est directe. Si tu échanges les deux vecteurs, la base $(\vec{j},\vec{i})$ devient \emph{indirecte} : l'angle orienté vaut $-\dfrac{\pi}{2}$.

Dans le plan, tout est simple : il n'y a que deux sens de rotation possibles, et l'expression « sens inverse des aiguilles d'une montre » suffit à tout dire, parce que nous regardons tous le plan \emph{du même côté} (depuis le dessus de la feuille).

\courssub{Le problème nouveau : dans l'espace, la montre ne suffit plus}

Dans l'espace, la situation se complique radicalement. Imaginons deux bases orthonormées $(\vec{i},\vec{j},\vec{k})$ et $(\vec{i},\vec{j},\vec{k}')$ qui partagent les deux mêmes premiers vecteurs, mais avec $\vec{k}' = -\vec{k}$ : l'une pointe « vers le haut », l'autre « vers le bas ».

\begin{itemize}
\item Ces deux bases sont toutes les deux orthonormées : les vecteurs sont unitaires et deux à deux orthogonaux.
\item Pourtant, \textbf{aucune rotation de l'espace ne permet de superposer l'une sur l'autre}. Essaye avec tes doigts : si le pouce et l'index sont fixés, le majeur ne peut pointer que d'un seul côté sans casser le poignet !
\end{itemize}

\begin{attention}[Deux « moitiés » de bases irréconciliables]
L'ensemble des bases orthonormées de l'espace se partage en \textbf{deux familles} : deux bases d'une même famille peuvent être déformées continûment (par rotation) l'une en l'autre ; deux bases de familles différentes ne le peuvent jamais. C'est exactement comme ta main droite et ta main gauche : ce sont des images l'une de l'autre dans un miroir, mais aucun mouvement dans l'espace ne permet de superposer un gant droit sur une main gauche. On dit que l'espace possède \cle{deux orientations possibles}, et il faut en \emph{choisir une}.
\end{attention}

Voilà pourquoi l'expression « sens inverse des aiguilles d'une montre » devient insuffisante : une rotation qui semble directe vue d'un côté du plan semble indirecte vue de l'autre côté ! Il faut une convention qui fixe un point de vue. Deux conventions classiques, rigoureusement équivalentes, sont au programme : la règle du \cle{bonhomme d'Ampère} et la règle de la \cle{main droite}.

\courssub{La règle du bonhomme d'Ampère}

\begin{definition}[Base orthonormée directe — règle du bonhomme d'Ampère]
Soit $(\vec{i},\vec{j},\vec{k})$ une base orthonormée de l'espace (trois vecteurs unitaires, deux à deux orthogonaux). On dit que cette base est \cle{directe} lorsque la condition suivante est satisfaite :

\emph{Un bonhomme debout sur le plan $(O\,;\vec{i},\vec{j})$, les pieds en $O$, la tête du côté de $\vec{k}$, et qui regarde dans la direction de $\vec{i}$, voit le vecteur $\vec{j}$ \textbf{à sa gauche}.}

De manière équivalente : le bonhomme, tête vers $\vec{k}$, voit le plan $(O\,;\vec{i},\vec{j})$ orienté dans le sens direct habituel (le passage de $\vec{i}$ à $\vec{j}$ se fait dans le sens inverse des aiguilles d'une montre \emph{vu de son point de vue}).

Si la condition n'est pas satisfaite (le bonhomme voit $\vec{j}$ à sa droite), la base est dite \cle{indirecte}.
\end{definition}

Le point essentiel est la position du bonhomme : c'est lui qui fixe « de quel côté » on regarde le plan. Si $\vec{k}$ pointe vers le haut, le bonhomme regarde le plan depuis le dessus, et tout se passe comme en Seconde.

\begin{popfigure}
\begin{center}
\begin{tikzpicture}[scale=1.1, line cap=round, line join=round]
% Plan (O;i,j)
\fill[popblueL, opacity=0.55] (-2.2,-0.9) -- (2.6,-0.9) -- (3.4,0.9) -- (-1.4,0.9) -- cycle;
\draw[popmuted] (-2.2,-0.9) -- (2.6,-0.9) -- (3.4,0.9) -- (-1.4,0.9) -- cycle;
% Axes
\draw[-{Stealth[length=3mm]}, very thick, popblue] (0,0) -- (2.4,0) node[below right] {$\vec{i}$};
\draw[-{Stealth[length=3mm]}, very thick, poporange] (0,0) -- (1.1,0.75) node[right] {$\vec{j}$};
\draw[-{Stealth[length=3mm]}, very thick, popgreen] (0,0) -- (0,2.6) node[above] {$\vec{k}$};
\fill (0,0) circle (1.6pt) node[below left] {$O$};
% Angle direct i -> j
\draw[-{Stealth[length=2mm]}, poppurple, thick] (0.85,0) arc (0:34:0.85) node[midway, right, xshift=2pt] {$+\frac{\pi}{2}$};
% Bonhomme d'Ampère stylisé (debout, tête vers k, regarde vers i)
\begin{scope}[shift={(-1.15,0)}]
\draw[popdark, thick] (0,0) -- (0,0.75); % corps
\draw[popdark, thick] (0,0.75) -- (0,1.15); % cou
\fill[popgold] (0,1.32) circle (0.17); % tête
% Regarde vers i (droite dans le repère local du bonhomme, mais i est à droite sur le schéma)
% Le bonhomme est à x=-1.15, il regarde vers l'origine (0,0) donc vers la droite (direction i)
\draw[-{Stealth[length=1.6mm]}, popdark, thick] (0,1.32) -- (0.5,1.32) node[midway, above, font=\tiny] {regard $\vec{i}$};
\draw[popdark, thick] (0,0.55) -- (-0.35,0.85); % bras gauche (vers j, à sa gauche)
\draw[popdark, thick] (0,0.55) -- (0.35,0.4); % bras droit
\draw[popdark, thick] (0,0) -- (-0.22,-0.45); % jambe
\draw[popdark, thick] (0,0) -- (0.22,-0.45); % jambe
\draw[-{Stealth[length=1.6mm]}, popdark] (0,1.5) -- (0,1.95) node[above, font=\tiny] {$\vec{k}$};
\end{scope}
\node[popdark, align=center, font=\small] at (-1.15,-0.95) {bonhomme d'Ampère\\ tête vers $\vec{k}$, regarde $\vec{i}$};
\end{tikzpicture}
\end{center}
\legende{Base orthonormée directe $(\vec{i},\vec{j},\vec{k})$ : le bonhomme d'Ampère, debout sur le plan $(O\,;\vec{i},\vec{j})$, la tête du côté de $\vec{k}$, regarde vers $\vec{i}$ et voit $\vec{j}$ à sa gauche.}
\end{popfigure}

\begin{savaistu}[D'où vient ce bonhomme ?]
André-Marie Ampère (1775--1836) était un physicien et mathématicien français, père de l'électromagnétisme. Sa règle du « bonhomme » servait à l'origine à déterminer le sens du \emph{champ magnétique} autour d'un fil parcouru par un courant électrique : le bonhomme est couché le long du fil, le courant entrant par les pieds et sortant par la tête ; il regarde le point où l'on mesure le champ, et le champ magnétique pointe alors vers sa \emph{gauche}. C'est cette même règle, transposée, que nous utilisons pour orienter l'espace. Sans elle, pas de moteurs électriques, pas de centrales hydroélectriques comme celle de Song Loulou, et pas de réseau ENEO !
\end{savaistu}

\courssub{La règle de la main droite}

La règle du bonhomme d'Ampère demande un peu d'imagination. Heureusement, tu possèdes un outil toujours disponible, même le jour du Baccalauréat : ta \textbf{main droite}.

\begin{definition}[Règle de la main droite]
Soit $(\vec{i},\vec{j},\vec{k})$ une base orthonormée de l'espace. Elle est \cle{directe} si et seulement si l'on peut associer :
\begin{itemize}
\item le \textbf{pouce} de la main droite à $\vec{i}$ ;
\item l'\textbf{index} à $\vec{j}$ ;
\item le \textbf{majeur} à $\vec{k}$ ;
\end{itemize}
les trois doigts étant écartés de manière à former un trièdre (pouce et index à angle droit, majeur dressé perpendiculairement à la paume).
\end{definition}

\begin{attention}[Toujours la main DROITE]
C'est le piège classique : la règle fonctionne \textbf{uniquement avec la main droite}. Avec la main gauche, tu obtiens une base indirecte ! Pendant l'épreuve, n'hésite pas à poser ton stylo une seconde et à vérifier physiquement avec ta main. Les trois doigts doivent former un vrai trièdre : pouce horizontal vers l'avant, index horizontal vers la droite (ou l'inverse selon la convention choisie), majeur vertical vers le haut.
\end{attention}

\begin{popfigure}
\begin{center}
\begin{tikzpicture}[scale=1.0, line cap=round, line join=round]
% Paume stylisée
\fill[popgoldL] (0,0) ellipse (1.05 and 0.8);
\draw[popdark, thick] (0,0) ellipse (1.05 and 0.8);
% Pouce -> i (vers la gauche-haut)
\draw[popdark, thick, fill=popgoldL] (-0.85,0.25) .. controls (-1.5,0.55) and (-1.9,0.75) .. (-2.15,0.9)
   .. controls (-2.25,0.98) and (-2.2,1.1) .. (-2.05,1.08)
   .. controls (-1.7,1.0) and (-1.25,0.85) .. (-0.75,0.62) -- cycle;
% Index -> j (vers la droite)
\draw[popdark, thick, fill=popgoldL] (0.55,0.55) .. controls (1.3,0.85) and (2.1,1.0) .. (2.5,1.05)
   .. controls (2.65,1.06) and (2.68,1.2) .. (2.55,1.25)
   .. controls (2.05,1.35) and (1.25,1.25) .. (0.5,0.95) -- cycle;
% Majeur -> k (vers le haut)
\draw[popdark, thick, fill=popgoldL] (-0.12,0.72) .. controls (-0.18,1.5) and (-0.18,2.2) .. (-0.12,2.6)
   .. controls (-0.1,2.75) and (0.12,2.75) .. (0.14,2.6)
   .. controls (0.2,2.2) and (0.2,1.5) .. (0.16,0.72) -- cycle;
% Flèches et labels
\draw[-{Stealth[length=3mm]}, very thick, popblue] (-2.15,1.0) -- (-2.9,1.35) node[above left] {$\vec{i}$ (pouce)};
\draw[-{Stealth[length=3mm]}, very thick, poporange] (2.55,1.15) -- (3.3,1.3) node[right] {$\vec{j}$ (index)};
\draw[-{Stealth[length=3mm]}, very thick, popgreen] (0,2.7) -- (0,3.4) node[above] {$\vec{k}$ (majeur)};
\node[popdark, font=\small] at (0,-0.35) {main droite};
\end{tikzpicture}
\end{center}
\legende{La règle de la main droite : pouce $=\vec{i}$, index $=\vec{j}$, majeur $=\vec{k}$. Si les trois vecteurs de la base coïncident avec ces trois doigts, la base est directe.}
\end{popfigure}

\begin{methode}[Vérifier rapidement qu'une base est directe]
Pour vérifier qu'une base orthonormée $(\vec{u},\vec{v},\vec{w})$ est directe :
\begin{enumerate}
\item \textbf{Identifie} le plan des deux premiers vecteurs $(\vec{u},\vec{v})$ et le sens de rotation de $\vec{u}$ vers $\vec{v}$ (le plus court chemin angulaire).
\item \textbf{Place-toi} mentalement (ou avec le bonhomme d'Ampère) du côté vers lequel pointe $\vec{w}$, et regarde le plan depuis ce côté.
\item \textbf{Observe} : si le passage de $\vec{u}$ à $\vec{v}$ se fait dans le sens inverse des aiguilles d'une montre (sens direct), la base est \emph{directe} ; sinon elle est \emph{indirecte}.
\item \textbf{Variante main droite} : associe pouce $=\vec{u}$, index $=\vec{v}$, majeur $=\vec{w}$. Si c'est possible sans tordre le poignet, la base est directe.
\end{enumerate}
\end{methode}

\courssub{Bases directes et bases indirectes : définition et premiers exemples}

\begin{definition}[Base orthonormée directe (BOND) et base orthonormée indirecte]
L'espace étant orienté par l'une des deux règles équivalentes précédentes :
\begin{itemize}
\item une base orthonormée $(\vec{i},\vec{j},\vec{k})$ est dite \cle{directe} (on note parfois BOND : base orthonormée directe) si elle satisfait la règle du bonhomme d'Ampère, ou de façon équivalente la règle de la main droite ;
\item elle est dite \cle{indirecte} dans le cas contraire.
\end{itemize}
Un repère $(O\,;\vec{i},\vec{j},\vec{k})$ est dit \emph{orthonormé direct} lorsque sa base est une BOND.
\end{definition}

\begin{exemplebox}[Le repère « classique » est direct]
Considérons le repère usuel de l'espace : $\vec{i}$ vers l'Est, $\vec{j}$ vers le Nord, $\vec{k}$ vers le haut (vers le zénith).
\begin{itemize}
\item \textbf{Main droite :} pouce $\to$ Est ($\vec{i}$), index $\to$ Nord ($\vec{j}$), majeur $\to$ Haut ($\vec{k}$). C'est parfaitement naturel : la base est \textbf{directe}.
\item \textbf{Bonhomme d'Ampère :} debout sur le sol (plan horizontal), la tête vers le ciel ($\vec{k}$), il regarde vers l'Est ($\vec{i}$). Le Nord ($\vec{j}$) est bien \emph{à sa gauche}.
\end{itemize}
Les deux règles donnent le même verdict. Dans les figures de ce cours, on prendra toujours $\vec{i}$ vers la droite, $\vec{j}$ en profondeur vers l'arrière-droit, et $\vec{k}$ vers le haut : c'est une base directe (vérifie avec ta main droite).
\end{exemplebox}

\begin{popfigure}
\begin{center}
\begin{tikzpicture}[scale=0.92, line cap=round, line join=round]
% ---- Trièdre direct ----
\begin{scope}
\fill[popgreenL, opacity=0.5] (-1.9,-0.7) -- (2.1,-0.7) -- (2.8,0.7) -- (-1.2,0.7) -- cycle;
\draw[-{Stealth[length=2.6mm]}, very thick, popblue] (0,0) -- (1.9,0) node[below right] {$\vec{i}$};
\draw[-{Stealth[length=2.6mm]}, very thick, poporange] (0,0) -- (0.95,0.65) node[right] {$\vec{j}$};
\draw[-{Stealth[length=2.6mm]}, very thick, popgreen] (0,0) -- (0,2.1) node[above] {$\vec{k}$};
\fill (0,0) circle (1.4pt) node[below left] {$O$};
\draw[-{Stealth[length=1.8mm]}, poppurple, thick] (0.7,0) arc (0:34:0.7);
\node[font=\small\bfseries, popdark] at (0.4,-1.15) {Base DIRECTE};
\end{scope}
% ---- Trièdre indirect ----
\begin{scope}[xshift=7.2cm]
\fill[poppinkL, opacity=0.5] (-1.9,-0.7) -- (2.1,-0.7) -- (2.8,0.7) -- (-1.2,0.7) -- cycle;
\draw[-{Stealth[length=2.6mm]}, very thick, poporange] (0,0) -- (1.9,0) node[below right] {$\vec{j}$};
\draw[-{Stealth[length=2.6mm]}, very thick, popblue] (0,0) -- (0.95,0.65) node[right] {$\vec{i}$};
\draw[-{Stealth[length=2.6mm]}, very thick, popgreen] (0,0) -- (0,2.1) node[above] {$\vec{k}$};
\fill (0,0) circle (1.4pt) node[below left] {$O$};
\draw[-{Stealth[length=1.8mm]}, poppurple, thick] (0.7,0) arc (0:34:0.7);
\node[font=\small\bfseries, popdark] at (0.4,-1.15) {Base INDIRECTE ($\vec{i}$ et $\vec{j}$ échangés)};
\end{scope}
\end{tikzpicture}
\end{center}
\legende{À gauche, une base directe : vu du dessus (côté $\vec{k}$), on passe de $\vec{i}$ à $\vec{j}$ dans le sens trigonométrique. À droite, on a simplement permuté $\vec{i}$ et $\vec{j}$ : la base devient indirecte.}
\end{popfigure}

\courssub{Effet d'une permutation ou d'un changement de signe}

Voici le point le plus délicat — et le plus souvent testé — de cette section. Que devient l'orientation d'une base lorsqu'on en modifie l'ordre ou le sens des vecteurs ?

\begin{propriete}[Permutations et orientation]
Soit $(\vec{i},\vec{j},\vec{k})$ une base orthonormée directe. Alors :
\begin{enumerate}
\item \textbf{Permutation de deux vecteurs} : les bases $(\vec{j},\vec{i},\vec{k})$, $(\vec{i},\vec{k},\vec{j})$ et $(\vec{k},\vec{j},\vec{i})$ sont \emph{indirectes}. Échanger deux vecteurs \textbf{change toujours} l'orientation.
\item \textbf{Permutation circulaire} : les bases $(\vec{j},\vec{k},\vec{i})$ et $(\vec{k},\vec{i},\vec{j})$ sont \emph{directes}. Faire « tourner » les trois vecteurs conserve l'orientation.
\item \textbf{Changement de signe} : remplacer un seul vecteur par son opposé (par exemple $(\vec{i},\vec{j},-\vec{k})$) rend la base \emph{indirecte}. Changer le signe de \textbf{deux} vecteurs (par exemple $(-\vec{i},-\vec{j},\vec{k})$) conserve l'orientation.
\end{enumerate}
\end{propriete}

\begin{experience}[Justification par la main droite]
Vérifions ces affirmations avec la main droite, sans aucun calcul.
\begin{enumerate}
\item \emph{Permutation de $\vec{i}$ et $\vec{j}$.} Place pouce $=\vec{i}$, index $=\vec{j}$, majeur $=\vec{k}$ (base directe). Pour la base $(\vec{j},\vec{i},\vec{k})$, il faudrait pouce $=\vec{j}$ et index $=\vec{i}$ : c'est impossible avec la main droite sans retourner la main paume vers le haut — auquel cas le majeur pointe vers $-\vec{k}$. La base est donc indirecte.
\item \emph{Permutation circulaire.} Pour $(\vec{j},\vec{k},\vec{i})$ : tourne simplement ton poignet pour que le pouce suive $\vec{j}$, l'index $\vec{k}$ ; le majeur pointe alors naturellement vers $\vec{i}$. Aucune contorsion : la base reste directe. C'est logique : une permutation circulaire, c'est \emph{deux} échanges successifs ($\vec{i}\leftrightarrow\vec{j}$ puis $\vec{i}\leftrightarrow\vec{k}$), et deux changements d'orientation se compensent.
\item \emph{Un signe moins.} Pour $(\vec{i},\vec{j},-\vec{k})$, le majeur devrait pointer à l'opposé de sa position naturelle : impossible sans passer à la main gauche. La base est indirecte. Avec \emph{deux} signes moins, par exemple $(-\vec{i},-\vec{j},\vec{k})$, il suffit de pivoter la main d'un demi-tour autour de l'axe vertical : pouce et index pointent vers les opposés, le majeur reste vers $\vec{k}$. Base directe.
\end{enumerate}
Retiens la règle de comptage : \textbf{chaque échange de deux vecteurs et chaque changement de signe d'un vecteur « coûte » un changement d'orientation}. Un nombre pair de telles opérations conserve l'orientation ; un nombre impair la renverse.
\end{experience}

\begin{attention}[Le piège de la permutation]
L'erreur la plus fréquente au Baccalauréat : on te donne une base directe $(\vec{i},\vec{j},\vec{k})$ et on te demande de travailler dans $(\vec{j},\vec{i},\vec{k})$ ou $(\vec{k},\vec{j},\vec{i})$ « par commodité ». Ces bases sont \textbf{indirectes} : tous les résultats de produit vectoriel que tu calculeras dedans seront \emph{faux en signe} ! Avant tout calcul, pose-toi la question : « ma base est-elle directe ? ». Si elle ne l'est pas, réordonne-la par permutation circulaire (qui conserve l'orientation) plutôt que par simple échange.
\end{attention}

\begin{exoresolu}[Directe ou indirecte ?]
L'espace est muni d'une base orthonormée directe $(\vec{i},\vec{j},\vec{k})$. Déterminer, en justifiant, l'orientation de chacune des bases suivantes :
\begin{enumerate}
\item $(\vec{k},\vec{i},\vec{j})$ ;
\item $(\vec{i},\vec{k},\vec{j})$ ;
\item $(-\vec{i},\vec{j},-\vec{k})$.
\end{enumerate}
\tcblower
\textbf{Solution détaillée.}

\textbf{1. Base $(\vec{k},\vec{i},\vec{j})$.} C'est une permutation circulaire de $(\vec{i},\vec{j},\vec{k})$ : chaque vecteur prend la place du précédent ($\vec{k}$ passe en tête, $\vec{i}$ et $\vec{j}$ reculent). On peut aussi compter : pour passer de $(\vec{i},\vec{j},\vec{k})$ à $(\vec{k},\vec{i},\vec{j})$, on échange $\vec{j}$ et $\vec{k}$ (1 changement), puis $\vec{j}$ et $\vec{i}$ (2 changements) : nombre \emph{pair} de permutations. La base $(\vec{k},\vec{i},\vec{j})$ est donc \textbf{directe}. Vérification main droite : pouce $=\vec{k}$ (vers le haut), index $=\vec{i}$, majeur $=\vec{j}$ : faisable en inclinant la main. Confirmé.

\textbf{2. Base $(\vec{i},\vec{k},\vec{j})$.} On a simplement échangé les deux derniers vecteurs $\vec{j}$ et $\vec{k}$ : un \emph{seul} échange, donc un nombre impair de permutations. La base $(\vec{i},\vec{k},\vec{j})$ est \textbf{indirecte}.

\textbf{3. Base $(-\vec{i},\vec{j},-\vec{k})$.} Deux vecteurs ont changé de signe : $\vec{i}$ et $\vec{k}$. Chaque changement de signe renverse l'orientation ; deux changements la restaurent. La base $(-\vec{i},\vec{j},-\vec{k})$ est donc \textbf{directe}. Géométriquement, on a effectué un demi-tour autour de l'axe $(O\,;\vec{j})$, ce qui est une rotation : l'orientation est conservée.
\end{exoresolu}

\courssub{Convention pour toute la suite du cours}

\begin{aretenir}[Convention définitive]
Dans toute la suite de cette leçon — et dans tout le programme de Terminale C et E — \textbf{l'espace est supposé orienté} et muni d'un repère orthonormé \cle{direct} $(O\,;\vec{i},\vec{j},\vec{k})$. Toutes les formules du produit vectoriel que nous établirons (expression analytique, calcul d'aires, de distances et de volumes) ne sont valables que dans un tel repère. Si un exercice fournit un repère sans préciser son orientation, vérifie-la avec la main droite avant de commencer tout calcul.
\end{aretenir}

\begin{exercice}[À toi de jouer]
L'espace est muni d'une base orthonormée directe $(\vec{i},\vec{j},\vec{k})$. Pour chacune des bases suivantes, dire si elle est directe ou indirecte, en comptant les permutations et changements de signe :
\begin{enumerate}
\item $(\vec{j},\vec{k},\vec{i})$ ;
\item $(-\vec{i},\vec{j},\vec{k})$ ;
\item $(-\vec{i},-\vec{j},-\vec{k})$ ;
\item $(\vec{j},\vec{i},-\vec{k})$.
\end{enumerate}
\end{exercice}

\begin{corrige}[Exercice 1]
\textbf{1.} $(\vec{j},\vec{k},\vec{i})$ : permutation circulaire, équivalente à deux échanges (nombre pair) : base \textbf{directe}.

\textbf{2.} $(-\vec{i},\vec{j},\vec{k})$ : un seul changement de signe (nombre impair) : base \textbf{indirecte}.

\textbf{3.} $(-\vec{i},-\vec{j},-\vec{k})$ : trois changements de signe (nombre impair) : base \textbf{indirecte}. C'est la base « symétrique » par rapport à $O$, image dans un miroir : l'orientation est renversée.

\textbf{4.} $(\vec{j},\vec{i},-\vec{k})$ : un échange ($\vec{i}\leftrightarrow\vec{j}$) et un changement de signe ($\vec{k}\to-\vec{k}$), soit deux opérations (nombre pair) : base \textbf{directe}. Vérification main droite : pouce $=\vec{j}$, index $=\vec{i}$, le majeur pointe alors vers $-\vec{k}$ : confirmé.
\end{corrige}

Tu sais désormais ce qu'orienter l'espace signifie, et tu maîtrises les deux règles — bonhomme d'Ampère et main droite — qui permettent de reconnaître une base directe. Cette convention, qui peut sembler anodine, est la clé de voûte de toute la leçon : c'est elle qui donnera un \emph{sens} bien déterminé au produit vectoriel, et donc au support perpendiculaire du panneau solaire de notre technicien. Passons maintenant à la définition proprement dite de cet outil remarquable.

\coursec{Définition et interprétation géométrique du produit vectoriel}

Dans la section précédente, nous avons appris à \cle{orienter} l'espace~: grâce à la règle du bonhomme d'Ampère (ou de la main droite), nous savons désormais dire si un triplet de vecteurs $(\vec{u},\vec{v},\vec{w})$ forme une base \emph{directe} ou \emph{indirecte}. Cet outil va maintenant nous servir à construire un objet nouveau et remarquable~: à partir de deux vecteurs $\vec{u}$ et $\vec{v}$, nous allons fabriquer un \textbf{troisième vecteur}, noté $\vec{u}\wedge\vec{v}$, qui sera perpendiculaire aux deux premiers et dont la longueur mesurera une \emph{aire}. C'est le \cle{produit vectoriel}.

Pourquoi un tel objet~? Imagine un mécanicien de quartier à Douala qui serre un boulon avec une clé~: l'efficacité du serrage (le \emph{moment} de la force) dépend de la force appliquée, de la longueur du bras de levier, et de l'angle entre les deux — et son effet est une rotation autour d'un axe \emph{perpendiculaire} au plan formé par la clé et la force. Ce « vecteur de rotation » est exactement un produit vectoriel. La physique regorge de telles situations, et la géométrie nous en donne la définition précise.

\courssub{Le cas particulier~: deux vecteurs colinéaires}

Commençons par le cas le plus simple. Si $\vec{u}$ et $\vec{v}$ sont colinéaires (c'est-à-dire portés par des droites parallèles, ou si l'un des deux est nul), il n'existe pas de direction privilégiée perpendiculaire « au plan » qu'ils définiraient... puisqu'ils ne définissent aucun plan~! Par convention, on pose alors~:

\begin{definition}[Produit vectoriel de deux vecteurs colinéaires]
Si $\vec{u}$ et $\vec{v}$ sont deux vecteurs \textbf{colinéaires} de l'espace (en particulier si l'un d'eux est nul), alors
\[ \vec{u}\wedge\vec{v} = \vec{0}. \]
\end{definition}

\begin{attention}[Le résultat est le vecteur nul]
On écrit $\vec{u}\wedge\vec{v} = \vec{0}$ (avec la flèche~!) et non $\vec{u}\wedge\vec{v} = 0$. Le produit vectoriel de deux vecteurs est un \textbf{vecteur}, jamais un nombre. C'est une différence fondamentale avec le produit scalaire, qui, lui, renvoie un réel.
\end{attention}

\courssub{Définition du produit vectoriel de deux vecteurs non colinéaires}

Supposons maintenant $\vec{u}$ et $\vec{v}$ \textbf{non colinéaires}. Ils définissent alors un plan (tout plan contenant des représentants de $\vec{u}$ et $\vec{v}$ de même origine). Le produit vectoriel $\vec{u}\wedge\vec{v}$ sera entièrement déterminé par trois données~: sa \emph{direction}, son \emph{sens} et sa \emph{norme}.

\begin{definition}[Produit vectoriel de deux vecteurs non colinéaires]
Soit $\vec{u}$ et $\vec{v}$ deux vecteurs non colinéaires de l'espace orienté. Le \cle{produit vectoriel} de $\vec{u}$ par $\vec{v}$, noté $\vec{u}\wedge\vec{v}$ (lire « $\vec{u}$ vectoriel $\vec{v}$ »), est le vecteur caractérisé par~:
\begin{enumerate}
\item \textbf{Direction~:} $\vec{u}\wedge\vec{v}$ est orthogonal à $\vec{u}$ \emph{et} à $\vec{v}$~;
\item \textbf{Sens~:} le triplet $\bigl(\vec{u},\,\vec{v},\,\vec{u}\wedge\vec{v}\bigr)$ est une base \textbf{directe} de l'espace (règle du bonhomme d'Ampère~: l'observateur, les pieds sur $\vec{u}$, regarde vers $\vec{v}$~; $\vec{u}\wedge\vec{v}$ pointe vers le haut si $\vec{v}$ est à sa gauche)~;
\item \textbf{Norme~:} si $\theta$ désigne la mesure dans $[0\,;\pi]$ de l'angle géométrique formé par $\vec{u}$ et $\vec{v}$, alors
\[ \bigl\|\vec{u}\wedge\vec{v}\bigr\| = \|\vec{u}\|\times\|\vec{v}\|\times\sin\theta. \]
\end{enumerate}
\end{definition}

\begin{aretenir}[Les trois ingrédients à retenir]
Pour décrire $\vec{u}\wedge\vec{v}$, retiens le triptyque~:
\begin{itemize}
\item \textbf{perpendiculaire} au plan $(\vec{u},\vec{v})$~;
\item orienté de sorte que $(\vec{u},\vec{v},\vec{u}\wedge\vec{v})$ soit \textbf{directe}~;
\item de longueur $\|\vec{u}\|\,\|\vec{v}\|\,\sin\theta$.
\end{itemize}
Comme $\theta\in[0\,;\pi]$, on a $\sin\theta\geq 0$~: la norme est bien un réel positif, ce qui est rassurant~!
\end{aretenir}

\begin{popfigure}
\begin{center}
\begin{tikzpicture}[scale=1.2, line cap=round, line join=round]
% Coordonnées dans le plan (x,y) pour le parallélogramme, z vertical pour le produit vectoriel
\coordinate (O) at (0,0);
\coordinate (A) at (3.5,0);
\coordinate (B) at (1.2,2);
\coordinate (C) at (4.7,2);
% Parallélogramme (plan horizontal projeté)
\fill[popblueL, opacity=0.4] (O) -- (A) -- (C) -- (B) -- cycle;
\draw[popmuted, thick] (O) -- (A) -- (C) -- (B) -- cycle;
\draw[popmuted, dashed] (O) -- (B) (A) -- (C);
% Vecteurs u et v
\draw[-{Stealth[length=3mm]}, very thick, popdark] (O) -- (A) node[midway, below] {$\vec{u}$};
\draw[-{Stealth[length=3mm]}, very thick, popdark] (O) -- (B) node[midway, left] {$\vec{v}$};
% Produit vectoriel (vertical)
\draw[-{Stealth[length=4mm]}, very thick, popink] (O) -- (0,3.2) node[above] {$\vec{u}\wedge\vec{v}$};
% Angle theta
\draw[popgreen, thick] (0.8,0) arc[start angle=0, end angle=59, radius=0.8];
\node[popgreen] at (1.1,0.4) {$\theta$};
% Hauteur h (projection de v sur la perpendiculaire à u)
\coordinate (H) at (1.2,0);
\draw[dashed, poppurple, thick] (B) -- (H) node[midway, right, font=\small] {$h$};
\draw[poppurple] (1.0,0) -- (1.0,0.15) -- (1.15,0.15);
% Point O
\fill (O) circle (1.5pt) node[below left] {$O$};
% Base orthonormée discrète pour repère
\draw[popmuted, thin, ->] (0,0) -- (-0.5,-0.3) node[left, font=\tiny] {$x$};
\draw[popmuted, thin, ->] (0,0) -- (0.5,-0.3) node[below, font=\tiny] {$y$};
\draw[popmuted, thin, ->] (0,0) -- (0,0.5) node[above, font=\tiny] {$z$};
\end{tikzpicture}
\end{center}
\legende{Le produit vectoriel $\vec{u}\wedge\vec{v}$ est orthogonal au plan contenant $\vec{u}$ et $\vec{v}$~; sa norme est l'aire du parallélogramme bleu construit sur $\vec{u}$ et $\vec{v}$.}
\end{popfigure}

\begin{exemplebox}[Lecture immédiate de la définition]
Soit $\vec{u}$ et $\vec{v}$ tels que $\|\vec{u}\|=4$, $\|\vec{v}\|=3$ et l'angle $(\vec{u},\vec{v})$ mesure $30^\circ$. Alors
\[ \bigl\|\vec{u}\wedge\vec{v}\bigr\| = 4\times 3\times \sin 30^\circ = 12\times \tfrac{1}{2} = 6. \]
Le vecteur $\vec{u}\wedge\vec{v}$ est perpendiculaire au plan de $\vec{u}$ et $\vec{v}$, de longueur $6$, dirigé de sorte que $(\vec{u},\vec{v},\vec{u}\wedge\vec{v})$ soit directe.
\end{exemplebox}

\begin{attention}[Deux pièges classiques]
\begin{itemize}
\item \textbf{Ne pas confondre avec le produit scalaire~!} Le produit scalaire $\vec{u}\cdot\vec{v} = \|\vec{u}\|\,\|\vec{v}\|\cos\theta$ est un \emph{nombre}~; le produit vectoriel $\vec{u}\wedge\vec{v}$ est un \emph{vecteur}. On ne peut pas écrire « $\vec{u}\wedge\vec{v} = 6$ »~: on écrit $\|\vec{u}\wedge\vec{v}\| = 6$.
\item \textbf{Le sinus, pas le cosinus.} Si $\vec{u}$ et $\vec{v}$ sont orthogonaux, $\sin\theta = 1$ et $\|\vec{u}\wedge\vec{v}\| = \|\vec{u}\|\,\|\vec{v}\|$~: le produit vectoriel est alors de norme \emph{maximale}. C'est l'inverse du produit scalaire, qui est nul dans ce cas.
\end{itemize}
\end{attention}

\courssub{Interprétation géométrique de la norme~: l'aire du parallélogramme}

Voici l'une des plus belles idées de cette leçon~: la norme du produit vectoriel n'est pas un nombre abstrait, c'est une \textbf{aire} que l'on peut voir et mesurer.

\begin{propriete}[Norme et aire du parallélogramme]
Soit $\vec{u}$ et $\vec{v}$ deux vecteurs non colinéaires, et $O$ un point de l'espace. Posons $\overrightarrow{OA}=\vec{u}$ et $\overrightarrow{OB}=\vec{v}$. Alors
\[ \bigl\|\vec{u}\wedge\vec{v}\bigr\| = \mathcal{A}_{\text{parallélogramme}}, \]
où $\mathcal{A}_{\text{parallélogramme}}$ est l'aire du parallélogramme construit sur les côtés $[OA]$ et $[OB]$.
\end{propriete}

\begin{experience}[Justification~: base fois hauteur]
Notons $\theta$ l'angle géométrique entre $\vec{u}$ et $\vec{v}$, et $h$ la hauteur du parallélogramme relative à la base $[OA]$ (voir la figure ci-dessus~: $h$ est la distance de $B$ à la droite $(OA)$).
\begin{enumerate}
\item Dans le triangle rectangle formé par $B$ et son projeté orthogonal sur $(OA)$, on lit~: $h = OB\,\sin\theta = \|\vec{v}\|\sin\theta$.
\item L'aire d'un parallélogramme est « base $\times$ hauteur »~:
\[ \mathcal{A} = OA \times h = \|\vec{u}\|\times\|\vec{v}\|\times\sin\theta. \]
\item Or, par définition, $\|\vec{u}\wedge\vec{v}\| = \|\vec{u}\|\,\|\vec{v}\|\,\sin\theta$. On retrouve exactement l'aire du parallélogramme. \hfill$\blacksquare$
\end{enumerate}
\end{experience}

\begin{propriete}[Aire du triangle]
Avec les mêmes notations, l'aire du triangle $OAB$ est la moitié de celle du parallélogramme~:
\[ \mathcal{A}_{OAB} = \frac{1}{2}\,\bigl\|\overrightarrow{OA}\wedge\overrightarrow{OB}\bigr\|. \]
\end{propriete}

En effet, la diagonale $[AB]$ partage le parallélogramme en deux triangles de même aire. Cette formule deviendra un outil de calcul redoutable dès que nous disposerons de l'expression analytique du produit vectoriel (section suivante).

\begin{exoresolu}[Aire d'un triangle par le produit vectoriel]
Dans un plan, on donne trois points $M$, $N$, $P$ tels que $MN = 5$, $MP = 8$ et $\widehat{NMP} = 60^\circ$. Calculer l'aire du triangle $MNP$.
\tcblower
On applique la formule avec $\vec{u}=\overrightarrow{MN}$ et $\vec{v}=\overrightarrow{MP}$~:
\begin{align*}
\mathcal{A}_{MNP} &= \frac{1}{2}\,\bigl\|\overrightarrow{MN}\wedge\overrightarrow{MP}\bigr\| \\
&= \frac{1}{2}\times MN\times MP\times \sin\widehat{NMP} \\
&= \frac{1}{2}\times 5\times 8\times \sin 60^\circ = 20\times\frac{\sqrt{3}}{2} = 10\sqrt{3}.
\end{align*}
L'aire du triangle $MNP$ vaut $10\sqrt{3}\approx 17{,}3$ unités d'aire.
\end{exoresolu}

\courssub{Produits vectoriels des vecteurs d'une base orthonormée directe}

Plaçons-nous dans une base orthonormée \textbf{directe} $(\vec{i},\vec{j},\vec{k})$ de l'espace. Appliquons la définition à $\vec{i}\wedge\vec{j}$~:
\begin{itemize}
\item \textbf{Direction~:} $\vec{i}\wedge\vec{j}$ est orthogonal à $\vec{i}$ et à $\vec{j}$, donc colinéaire à $\vec{k}$~;
\item \textbf{Sens~:} la base étant directe, $\vec{i}\wedge\vec{j}$ a le même sens que $\vec{k}$~;
\item \textbf{Norme~:} $\|\vec{i}\wedge\vec{j}\| = \|\vec{i}\|\,\|\vec{j}\|\,\sin 90^\circ = 1\times 1\times 1 = 1$.
\end{itemize}
Donc $\vec{i}\wedge\vec{j} = \vec{k}$. En permutant les rôles, on obtient le cercle des relations fondamentales~:

\begin{aretenir}[Produits vectoriels dans une BOND]
Dans une base orthonormée directe $(\vec{i},\vec{j},\vec{k})$~:
\[ \vec{i}\wedge\vec{j} = \vec{k}, \qquad \vec{j}\wedge\vec{k} = \vec{i}, \qquad \vec{k}\wedge\vec{i} = \vec{j}. \]
En inversant l'ordre des facteurs, on obtient les opposés (nous justifierons cette antisymétrie dans la section 4)~:
\[ \vec{j}\wedge\vec{i} = -\vec{k}, \qquad \vec{k}\wedge\vec{j} = -\vec{i}, \qquad \vec{i}\wedge\vec{k} = -\vec{j}. \]
\textbf{Astuce mnémotechnique~:} dans l'ordre cyclique $\vec{i}\to\vec{j}\to\vec{k}\to\vec{i}$, le produit de deux vecteurs consécutifs donne le suivant~; dans l'ordre inverse, on obtient son opposé.
\end{aretenir}

\courssub{Premières lectures géométriques dans un cube}

Le cube est le terrain de jeu idéal pour s'entraîner à « voir » un produit vectoriel sans aucun calcul de coordonnées. Considérons le cube $ABCDEFGH$ d'arête $1$, avec $ABCD$ la face du bas et $EFGH$ la face du haut ($E$ au-dessus de $A$, etc.). Le triplet $\bigl(\overrightarrow{AB},\overrightarrow{AD},\overrightarrow{AE}\bigr)$ forme une base orthonormée que nous supposons directe.

\begin{exemplebox}[Dans le cube~: $\overrightarrow{AB}\wedge\overrightarrow{AD}$]
Déterminons $\overrightarrow{AB}\wedge\overrightarrow{AD}$ par la définition~:
\begin{enumerate}
\item \textbf{Direction~:} $\overrightarrow{AB}$ et $\overrightarrow{AD}$ engendrent le plan de la face $ABCD$~; un vecteur orthogonal aux deux est porté par la verticale $(AE)$.
\item \textbf{Sens~:} la base $\bigl(\overrightarrow{AB},\overrightarrow{AD},\overrightarrow{AE}\bigr)$ est directe, donc $\overrightarrow{AB}\wedge\overrightarrow{AD}$ a le sens de $\overrightarrow{AE}$.
\item \textbf{Norme~:} $\|\overrightarrow{AB}\wedge\overrightarrow{AD}\| = 1\times 1\times \sin 90^\circ = 1$, qui est bien l'aire de la face carrée $ABCD$.
\end{enumerate}
Conclusion~: $\overrightarrow{AB}\wedge\overrightarrow{AD} = \overrightarrow{AE}$.
\end{exemplebox}

\begin{popfigure}
\begin{center}
\begin{tikzpicture}[scale=1.5, line join=round, line cap=round]
% Sommets du cube (perspective cavalière)
\coordinate (A) at (0,0);
\coordinate (B) at (2.4,0);
\coordinate (C) at (3.1,0.9);
\coordinate (D) at (0.7,0.9);
\coordinate (E) at (0,2.4);
\coordinate (F) at (2.4,2.4);
\coordinate (G) at (3.1,3.3);
\coordinate (H) at (0.7,3.3);
% Face du bas coloriée
\fill[poporangeL, opacity=0.6] (A) -- (B) -- (C) -- (D) -- cycle;
% Arêtes visibles
\draw[thick, popdark] (A) -- (B) -- (C) -- (G) -- (F) -- (B);
\draw[thick, popdark] (F) -- (E) -- (A);
\draw[thick, popdark] (E) -- (H) -- (G);
\draw[thick, popdark] (H) -- (D);
% Arêtes cachées
\draw[dashed, popmuted] (A) -- (D) -- (C);
% Vecteurs
\draw[-{Stealth[length=3mm]}, very thick, popblue] (A) -- (B) node[midway, below] {$\overrightarrow{AB}$};
\draw[-{Stealth[length=3mm]}, very thick, popgreen] (A) -- (D) node[midway, below right, font=\small] {$\overrightarrow{AD}$};
\draw[-{Stealth[length=3.5mm]}, very thick, popink] (A) -- (E) node[midway, left] {$\overrightarrow{AE}$};
% Sommets
\foreach \p/\pos in {A/below left, B/below, C/right, D/above left, E/above left, F/above, G/above right, H/above}
  \fill (\p) circle (1.2pt) node[\pos, font=\small] {$\p$};
\end{tikzpicture}
\end{center}
\legende{Dans le cube $ABCDEFGH$ d'arête $1$~: $\overrightarrow{AB}\wedge\overrightarrow{AD} = \overrightarrow{AE}$. La norme du produit vectoriel est l'aire de la face orange $ABCD$, soit $1$.}
\end{popfigure}

\begin{exoresolu}[Lecture dans le cube]
Dans le cube $ABCDEFGH$ d'arête $1$ ci-dessus, déterminer les vecteurs $\overrightarrow{AB}\wedge\overrightarrow{AE}$ et $\overrightarrow{AD}\wedge\overrightarrow{AB}$, puis l'aire du triangle $ABE$.
\tcblower
\textbf{1. Calcul de $\overrightarrow{AB}\wedge\overrightarrow{AE}$.} Ces deux vecteurs sont orthogonaux et de norme $1$, donc leur produit vectoriel est de norme $1$ et orthogonal au plan de la face $ABFE$~: il est colinéaire à $\overrightarrow{AD}$. Pour le sens, on vérifie avec la règle de la main droite que $\bigl(\overrightarrow{AB},\overrightarrow{AE},\overrightarrow{AD}\bigr)$ est \emph{indirecte} (on a échangé deux vecteurs par rapport à la base directe). Donc
\[ \overrightarrow{AB}\wedge\overrightarrow{AE} = -\overrightarrow{AD}. \]
\textbf{2. Calcul de $\overrightarrow{AD}\wedge\overrightarrow{AB}$.} On inverse l'ordre par rapport à l'exemple précédent~: le sens change, la norme reste $1$~:
\[ \overrightarrow{AD}\wedge\overrightarrow{AB} = -\overrightarrow{AE}. \]
\textbf{3. Aire du triangle $ABE$.} C'est la moitié de l'aire de la face carrée $ABFE$~:
\[ \mathcal{A}_{ABE} = \frac{1}{2}\bigl\|\overrightarrow{AB}\wedge\overrightarrow{AE}\bigr\| = \frac{1}{2}\times 1 = \frac{1}{2}. \]
\end{exoresolu}

\begin{attention}[Bien vérifier l'orientation de la base]
Dans le cube, tout repose sur le fait que $\bigl(\overrightarrow{AB},\overrightarrow{AD},\overrightarrow{AE}\bigr)$ est \emph{directe}. Si l'énoncé ne le précise pas, c'est à toi de le vérifier avec le bonhomme d'Ampère avant de conclure sur le sens du produit vectoriel. Un sens faux, et tout le résultat est opposé~!
\end{attention}

\begin{savaistu}[Le produit vectoriel est partout en physique]
Le produit vectoriel n'est pas qu'un outil de géomètre~: il gouverne des phénomènes que tu observes chaque jour.
\begin{itemize}
\item \textbf{Moment d'une force~:} quand le mécanicien serre un boulon, le moment de la force $\vec{F}$ par rapport à l'axe est $\vec{\mathcal{M}} = \overrightarrow{OA}\wedge\vec{F}$. C'est pourquoi une clé plus longue (bras de levier plus grand) serre plus fort à force égale~: la norme $\|\vec{\mathcal{M}}\| = OA\times\|\vec{F}\|\times\sin\theta$ augmente.
\item \textbf{Force de Lorentz~:} une particule de charge $q$ animée d'une vitesse $\vec{v}$ dans un champ magnétique $\vec{B}$ subit la force $\vec{F} = q\,\vec{v}\wedge\vec{B}$. C'est ce principe qui fait tourner les moteurs électriques... y compris ceux des groupes électrogènes qui dépannent lors des délestages~!
\item \textbf{Rotation de la Terre~:} la force de Coriolis, responsable du sens de rotation des cyclones, s'écrit aussi avec un produit vectoriel.
\end{itemize}
Le mathématicien irlandais William Rowan Hamilton et le physicien américain Josiah Willard Gibbs ont contribué, au XIX\textsuperscript{e} siècle, à faire du produit vectoriel le langage universel de la physique de l'espace.
\end{savaistu}

\begin{methode}[Déterminer un produit vectoriel par la définition]
Pour déterminer géométriquement $\vec{u}\wedge\vec{v}$ (sans coordonnées) ~:
\begin{enumerate}
\item Vérifier si $\vec{u}$ et $\vec{v}$ sont colinéaires~: si oui, $\vec{u}\wedge\vec{v}=\vec{0}$, c'est terminé.
\item Sinon, identifier le plan $(\vec{u},\vec{v})$ et chercher un vecteur connu orthogonal à ce plan (une arête de cube, une hauteur...)~: c'est la \textbf{direction}.
\item Appliquer la règle du bonhomme d'Ampère au triplet $(\vec{u},\vec{v},\vec{w})$ pour fixer le \textbf{sens}.
\item Calculer la \textbf{norme} $\|\vec{u}\|\,\|\vec{v}\|\,\sin\theta$, en l'interprétant si possible comme une aire.
\end{enumerate}
\end{methode}

Nous savons désormais ce qu'\emph{est} le produit vectoriel et ce que sa norme \emph{mesure}. Mais le calcul géométrique direct atteint vite ses limites~: comment trouver $\vec{u}\wedge\vec{v}$ lorsque $\vec{u}$ et $\vec{v}$ sont donnés par leurs coordonnées dans un repère quelconque~? C'est toute l'affaire de la section suivante, où nous établirons l'expression analytique du produit vectoriel dans une base orthonormée directe.

\coursec{Expression analytique dans une base orthonormée directe}

Dans la section précédente, nous avons défini le produit vectoriel $\vec{u} \wedge \vec{v}$ de manière \emph{géométrique} : direction (perpendiculaire au plan des deux vecteurs), sens (règle de la main droite), norme ($\|\vec{u}\|\,\|\vec{v}\|\,|\sin\theta|$). Cette définition est belle, mais peu pratique pour calculer : dès que les vecteurs sont donnés par leurs coordonnées, mesurer un angle et orienter « à la main » devient acrobatique. Toute cette section a donc un seul objectif : transformer la géométrie en \cle{calcul}. Nous allons établir une formule qui, à partir des coordonnées de $\vec{u}$ et $\vec{v}$ dans une base orthonormée directe, donne \emph{directement} les coordonnées de $\vec{u} \wedge \vec{v}$. C'est cette formule que tu utiliseras dans presque tous les exercices du Baccalauréat.

\courssub{Le cadre : une base orthonormée directe}

\begin{definition}[Base orthonormée directe]
Une base $(\vec{i}, \vec{j}, \vec{k})$ de l'espace est dite \cle{orthonormée directe} (en abrégé \cle{BOND}) lorsque :
\begin{itemize}
\item les trois vecteurs sont unitaires et deux à deux orthogonaux : $\|\vec{i}\| = \|\vec{j}\| = \|\vec{k}\| = 1$ et $\vec{i} \cdot \vec{j} = \vec{j} \cdot \vec{k} = \vec{k} \cdot \vec{i} = 0$ ;
\item le trièdre $(\vec{i}, \vec{j}, \vec{k})$ est orienté dans le sens direct (règle de la main droite : pouce $= \vec{i}$, index $= \vec{j}$, majeur $= \vec{k}$).
\end{itemize}
\end{definition}

Dans une telle base, la définition géométrique du produit vectoriel donne immédiatement les produits des vecteurs de base entre eux. En effet, $\vec{i}$ et $\vec{j}$ sont orthogonaux et unitaires, donc $\|\vec{i} \wedge \vec{j}\| = 1 \times 1 \times \sin 90\degres = 1$, et le sens direct impose :

\begin{aretenir}[Produits vectoriels des vecteurs de base]
Dans une base orthonormée \textbf{directe} $(\vec{i}, \vec{j}, \vec{k})$ :
\[ \vec{i} \wedge \vec{j} = \vec{k}, \qquad \vec{j} \wedge \vec{k} = \vec{i}, \qquad \vec{k} \wedge \vec{i} = \vec{j}. \]
Par antisymétrie : $\vec{j} \wedge \vec{i} = -\vec{k}$, $\vec{k} \wedge \vec{j} = -\vec{i}$, $\vec{i} \wedge \vec{k} = -\vec{j}$. Et bien sûr $\vec{i} \wedge \vec{i} = \vec{j} \wedge \vec{j} = \vec{k} \wedge \vec{k} = \vec{0}$.
\end{aretenir}

\begin{methode}[Retenir le cycle $\vec{i} \to \vec{j} \to \vec{k}$]
Les trois égalités $\vec{i} \wedge \vec{j} = \vec{k}$, $\vec{j} \wedge \vec{k} = \vec{i}$, $\vec{k} \wedge \vec{i} = \vec{j}$ suivent un \emph{cycle} : le produit de deux vecteurs consécutifs du cycle donne le troisième, avec le signe $+$ si l'on suit le sens du cycle, et le signe $-$ sinon. Dessine un cercle avec $\vec{i}$, $\vec{j}$, $\vec{k}$ placés dans cet ordre : tu ne pourras plus te tromper.
\end{methode}

\courssub{Le théorème fondamental}

\begin{propriete}[Expression du produit vectoriel dans une BOND]
Soit $(\vec{i}, \vec{j}, \vec{k})$ une base orthonormée directe de l'espace, et soient $\vec{u}(x\,; y\,; z)$ et $\vec{v}(x'\,; y'\,; z')$ deux vecteurs. Alors :
\[ \vec{u} \wedge \vec{v} = \begin{pmatrix} yz' - zy' \\ zx' - xz' \\ xy' - yx' \end{pmatrix} \]
Autrement dit : $\vec{u} \wedge \vec{v} = (yz' - zy')\,\vec{i} + (zx' - xz')\,\vec{j} + (xy' - yx')\,\vec{k}$.
\emph{(Démonstration exigible au Bac via la bilinéarité, voir l'activité guidée ci-dessous.)}
\end{propriete}

\begin{experience}[Démonstration guidée : développer par bilinéarité]
Nous admettons provisoirement la \emph{bilinéarité} du produit vectoriel (elle sera établie dans la section suivante) : on peut développer un produit vectoriel de sommes comme un produit ordinaire, en respectant l'ordre des facteurs (car $\vec{a} \wedge \vec{b} = -\vec{b} \wedge \vec{a}$).

\textbf{Étape 1.} Écris les décompositions :
\[ \vec{u} = x\vec{i} + y\vec{j} + z\vec{k}, \qquad \vec{v} = x'\vec{i} + y'\vec{j} + z'\vec{k}. \]

\textbf{Étape 2.} Développe le produit en neuf termes :
\begin{align*}
\vec{u} \wedge \vec{v} =\ & xx'\,(\vec{i} \wedge \vec{i}) + xy'\,(\vec{i} \wedge \vec{j}) + xz'\,(\vec{i} \wedge \vec{k}) \\
& + yx'\,(\vec{j} \wedge \vec{i}) + yy'\,(\vec{j} \wedge \vec{j}) + yz'\,(\vec{j} \wedge \vec{k}) \\
& + zx'\,(\vec{k} \wedge \vec{i}) + zy'\,(\vec{k} \wedge \vec{j}) + zz'\,(\vec{k} \wedge \vec{k}).
\end{align*}

\textbf{Étape 3.} Les produits d'un vecteur par lui-même sont nuls : $\vec{i} \wedge \vec{i} = \vec{j} \wedge \vec{j} = \vec{k} \wedge \vec{k} = \vec{0}$. Il reste six termes.

\textbf{Étape 4.} Remplace par les produits de la base directe :
$\vec{i} \wedge \vec{j} = \vec{k}$, $\vec{j} \wedge \vec{k} = \vec{i}$, $\vec{k} \wedge \vec{i} = \vec{j}$, et $\vec{i} \wedge \vec{k} = -\vec{j}$, $\vec{j} \wedge \vec{i} = -\vec{k}$, $\vec{k} \wedge \vec{j} = -\vec{i}$.

\textbf{Étape 5.} Regroupe selon $\vec{i}$, $\vec{j}$, $\vec{k}$ :
\begin{align*}
\vec{u} \wedge \vec{v} &= yz'\,\vec{i} - zy'\,\vec{i} + zx'\,\vec{j} - xz'\,\vec{j} + xy'\,\vec{k} - yx'\,\vec{k} \\
&= (yz' - zy')\,\vec{i} + (zx' - xz')\,\vec{j} + (xy' - yx')\,\vec{k}.
\end{align*}
C'est exactement la formule annoncée. $\blacksquare$
\end{experience}

\courssub{Comment retenir la formule : la règle en croix (règle du gamma)}

La formule semble indigeste, mais il existe une disposition visuelle qui la rend mécanique. Écris les coordonnées des deux vecteurs \emph{l'une sous l'autre}, en colonnes :

\begin{popfigure}
\begin{center}
\begin{tikzpicture}[scale=1.05, every node/.style={font=\normalsize}]
% En-têtes
\node[popblue, font=\bfseries] at (-2.2, 0.6) {$\vec{u}$};
\node[poporange, font=\bfseries] at (-2.2, -0.6) {$\vec{v}$};
% Coordonnées
\node (x)  at (0, 0.6)  {$x$};
\node (y)  at (1.6, 0.6)  {$y$};
\node (z)  at (3.2, 0.6)  {$z$};
\node (xp) at (0, -0.6) {$x'$};
\node (yp) at (1.6, -0.6) {$y'$};
\node (zp) at (3.2, -0.6) {$z'$};
% Croix pour la 1re coordonnée (y z')
\draw[->, popgreen, very thick] (y.south) -- (zp.north);
\draw[->, popgreen, very thick, dashed] (z.south) -- (yp.north);
\node[popgreen, font=\bfseries] at (5.3, 0) {$1^{\text{re}}$ coord. : $yz' - zy'$};
% Croix pour la 2e coordonnée (z x')
\draw[->, poppurple, very thick] (z.south west) -- (xp.north east);
\draw[->, poppurple, very thick, dashed] (x.south east) -- (zp.north west);
\node[poppurple, font=\bfseries] at (5.3, -1.2) {$2^{\text{e}}$ coord. : $zx' - xz'$};
% Croix pour la 3e coordonnée (x y')
\draw[->, popink, very thick] (x.south) -- (yp.north);
\draw[->, popink, very thick, dashed] (y.south) -- (xp.north);
\node[popink, font=\bfseries] at (5.3, -2.4) {$3^{\text{e}}$ coord. : $xy' - yx'$};
\end{tikzpicture}
\end{center}
\legende{Règle en croix : pour chaque coordonnée du résultat, on « masque » la colonne correspondante et on croise les deux autres (trait plein moins trait pointillé).}
\end{popfigure}

\begin{methode}[Calculer $\vec{u} \wedge \vec{v}$ par la règle en croix]
\begin{enumerate}
\item Écris les coordonnées en deux lignes superposées : $\vec{u}$ au-dessus, $\vec{v}$ en dessous.
\item \textbf{Première coordonnée} : masque la colonne des $x$, croise les colonnes $y$ et $z$ : calcule $yz' - zy'$ (diagonale descendante moins diagonale montante).
\item \textbf{Deuxième coordonnée} : masque la colonne des $y$, croise $z$ et $x$ : calcule $zx' - xz'$. Attention à l'ordre : on commence par $z$ !
\item \textbf{Troisième coordonnée} : masque la colonne des $z$, croise $x$ et $y$ : calcule $xy' - yx'$.
\item \textbf{Vérification systématique} : le résultat $\vec{n} = \vec{u} \wedge \vec{v}$ doit être orthogonal à $\vec{u}$ \emph{et} à $\vec{v}$. Calcule les deux produits scalaires $\vec{n} \cdot \vec{u}$ et $\vec{n} \cdot \vec{v}$ : ils doivent être nuls. Si l'un des deux ne l'est pas, il y a une erreur de calcul.
\end{enumerate}
\end{methode}

\begin{attention}[Pièges de la formule]
\begin{itemize}
\item Pour la deuxième coordonnée, c'est $zx' - xz'$ et \textbf{pas} $xz' - zx'$ : c'est l'erreur la plus fréquente. Retiens le cycle $x \to y \to z \to x$ : chaque coordonnée commence par la lettre qui \emph{suit} celle de la coordonnée calculée.
\item Dans chaque différence, le premier terme est le produit « coordonnée de $\vec{u}$ $\times$ coordonnée de $\vec{v}$ ». Inverser $\vec{u}$ et $\vec{v}$ change le signe du résultat (antisymétrie) : $\vec{v} \wedge \vec{u} = -\vec{u} \wedge \vec{v}$.
\item Chaque coordonnée est un \emph{déterminant} $2 \times 2$ : $yz' - zy' = \begin{vmatrix} y & z \\ y' & z' \end{vmatrix}$, etc. Si tu connais les déterminants, c'est un excellent moyen de contrôle.
\end{itemize}
\end{attention}

\courssub{Exemples chiffrés progressifs}

\begin{exemplebox}[Niveau 1 : coordonnées simples]
Soient $\vec{u}(1\,; 2\,; 0)$ et $\vec{v}(0\,; 1\,; 3)$ dans une BOND. Calculons $\vec{u} \wedge \vec{v}$.
\begin{align*}
\vec{u} \wedge \vec{v} = \begin{pmatrix} 2 \times 3 - 0 \times 1 \\ 0 \times 0 - 1 \times 3 \\ 1 \times 1 - 2 \times 0 \end{pmatrix} = \begin{pmatrix} 6 \\ -3 \\ 1 \end{pmatrix}.
\end{align*}
\emph{Vérification} : $\vec{n} \cdot \vec{u} = 6 \times 1 + (-3) \times 2 + 1 \times 0 = 0$ \checkmark{} et $\vec{n} \cdot \vec{v} = 6 \times 0 + (-3) \times 1 + 1 \times 3 = 0$ \checkmark. Le résultat est bien orthogonal aux deux vecteurs.
\end{exemplebox}

\begin{exemplebox}[Niveau 2 : coordonnées quelconques, avec des signes]
Soient $\vec{u}(2\,; -1\,; 3)$ et $\vec{v}(1\,; 4\,; -2)$. Alors :
\begin{align*}
\vec{u} \wedge \vec{v} = \begin{pmatrix} (-1)(-2) - 3 \times 4 \\ 3 \times 1 - 2 \times (-2) \\ 2 \times 4 - (-1) \times 1 \end{pmatrix} = \begin{pmatrix} 2 - 12 \\ 3 + 4 \\ 8 + 1 \end{pmatrix} = \begin{pmatrix} -10 \\ 7 \\ 9 \end{pmatrix}.
\end{align*}
\emph{Vérification} : $(-10) \times 2 + 7 \times (-1) + 9 \times 3 = -20 - 7 + 27 = 0$ \checkmark{} ; $(-10) \times 1 + 7 \times 4 + 9 \times (-2) = -10 + 28 - 18 = 0$ \checkmark.
\end{exemplebox}

\begin{exoresolu}[Un calcul complet avec vérification]
\textbf{Énoncé.} Dans une base orthonormée directe, on donne $\vec{u}(3\,; 0\,; -1)$ et $\vec{v}(-2\,; 5\,; 1)$. Calculer $\vec{u} \wedge \vec{v}$, puis $\vec{v} \wedge \vec{u}$, et vérifier l'orthogonalité.
\tcblower
\textbf{Solution.} Appliquons la règle en croix.
\begin{align*}
\vec{u} \wedge \vec{v} = \begin{pmatrix} 0 \times 1 - (-1) \times 5 \\ (-1) \times (-2) - 3 \times 1 \\ 3 \times 5 - 0 \times (-2) \end{pmatrix} = \begin{pmatrix} 0 + 5 \\ 2 - 3 \\ 15 - 0 \end{pmatrix} = \begin{pmatrix} 5 \\ -1 \\ 15 \end{pmatrix}.
\end{align*}
Par antisymétrie, $\vec{v} \wedge \vec{u} = -\vec{u} \wedge \vec{v} = (-5\,; 1\,; -15)$, ce que l'on retrouve en refaisant le calcul avec les lignes échangées.

\emph{Vérification d'orthogonalité} :
\begin{align*}
(\vec{u} \wedge \vec{v}) \cdot \vec{u} &= 5 \times 3 + (-1) \times 0 + 15 \times (-1) = 15 - 15 = 0, \\
(\vec{u} \wedge \vec{v}) \cdot \vec{v} &= 5 \times (-2) + (-1) \times 5 + 15 \times 1 = -10 - 5 + 15 = 0.
\end{align*}
Les deux produits scalaires sont nuls : le calcul est validé.
\end{exoresolu}

\courssub{Application immédiate : vecteur normal à un plan}

Voici la première grande utilité pratique de la formule. Un plan $(P)$ est déterminé par un point $A$ et deux vecteurs directeurs $\vec{u}$ et $\vec{v}$ non colinéaires. Le produit vectoriel $\vec{n} = \vec{u} \wedge \vec{v}$ est, par définition, orthogonal à $\vec{u}$ et à $\vec{v}$ : il est donc orthogonal à \emph{tout} vecteur du plan. C'est un \cle{vecteur normal} au plan $(P)$.

\begin{popfigure}
\begin{center}
\begin{tikzpicture}[scale=1.1, >=Stealth]
% Plan (parallélogramme)
\fill[popblueL, opacity=0.55] (-2.6,-0.4) -- (0.4,-1.4) -- (3.2,-0.2) -- (0.2,0.9) -- cycle;
\draw[popblue, thick] (-2.6,-0.4) -- (0.4,-1.4) -- (3.2,-0.2) -- (0.2,0.9) -- cycle;
\node[popblue] at (2.6, 0.55) {$(P)$};
% Point A
\coordinate (A) at (0.1,-0.25);
\fill[popdark] (A) circle (1.6pt) node[below left, popdark] {$A$};
% Vecteurs u et v dans le plan
\draw[->, poporange, very thick] (A) -- ++(1.9, 0.35) node[midway, above, poporange] {$\vec{u}$};
\draw[->, popgreen, very thick] (A) -- ++(1.0, -0.75) node[midway, below right, popgreen] {$\vec{v}$};
% Vecteur normal
\draw[->, popink, very thick] (A) -- ++(0.35, 1.9) node[right, popink] {$\vec{n} = \vec{u} \wedge \vec{v}$};
% Angle droit symbolique
\draw[popink] (0.28, 0.12) -- (0.42, 0.05) -- (0.35, -0.09);
\end{tikzpicture}
\end{center}
\legende{Le plan $(P)$ passe par $A$ et est dirigé par $\vec{u}$ et $\vec{v}$ ; le vecteur $\vec{n} = \vec{u} \wedge \vec{v}$ est normal à $(P)$.}
\end{popfigure}

\begin{methode}[Déterminer un vecteur normal au plan $(ABC)$]
\begin{enumerate}
\item Calcule les coordonnées de $\overrightarrow{AB}$ et $\overrightarrow{AC}$.
\item Vérifie qu'ils ne sont pas colinéaires (sinon $A$, $B$, $C$ sont alignés et ne définissent pas un plan).
\item Calcule $\vec{n} = \overrightarrow{AB} \wedge \overrightarrow{AC}$ par la règle en croix : c'est un vecteur normal au plan $(ABC)$.
\item Contrôle : $\vec{n} \cdot \overrightarrow{AB} = 0$ et $\vec{n} \cdot \overrightarrow{AC} = 0$.
\end{enumerate}
\end{methode}

\begin{exemplebox}[Vecteur normal et équation cartésienne du plan $(ABC)$]
Dans un repère orthonormé direct, on donne $A(1\,; 0\,; 2)$, $B(2\,; 1\,; 0)$ et $C(0\,; 1\,; 1)$.

\textbf{Étape 1.} $\overrightarrow{AB}(1\,; 1\,; -2)$ et $\overrightarrow{AC}(-1\,; 1\,; -1)$. Ils ne sont pas colinéaires (les coordonnées ne sont pas proportionnelles), donc $A$, $B$, $C$ définissent bien un plan.

\textbf{Étape 2.} Calculons le vecteur normal :
\begin{align*}
\vec{n} = \overrightarrow{AB} \wedge \overrightarrow{AC} = \begin{pmatrix} 1 \times (-1) - (-2) \times 1 \\ (-2) \times (-1) - 1 \times (-1) \\ 1 \times 1 - 1 \times (-1) \end{pmatrix} = \begin{pmatrix} -1 + 2 \\ 2 + 1 \\ 1 + 1 \end{pmatrix} = \begin{pmatrix} 1 \\ 3 \\ 2 \end{pmatrix}.
\end{align*}
\emph{Contrôle} : $\vec{n} \cdot \overrightarrow{AB} = 1 + 3 - 4 = 0$ \checkmark{} ; $\vec{n} \cdot \overrightarrow{AC} = -1 + 3 - 2 = 0$ \checkmark.

\textbf{Étape 3.} Une équation cartésienne du plan $(ABC)$ est de la forme $x + 3y + 2z + d = 0$ (les coefficients de $x$, $y$, $z$ sont les coordonnées de $\vec{n}$). Comme $A(1\,; 0\,; 2)$ appartient au plan : $1 + 0 + 4 + d = 0$, donc $d = -5$. Finalement :
\[ (ABC) : \quad x + 3y + 2z - 5 = 0. \]
\end{exemplebox}

\begin{aretenir}[Le lien essentiel avec l'équation d'un plan]
Si $\vec{n}(a\,; b\,; c)$ est un vecteur normal au plan $(P)$, alors $(P)$ admet une équation cartésienne de la forme
\[ ax + by + cz + d = 0, \]
où $d$ se détermine en écrivant qu'un point connu du plan vérifie l'équation. Le produit vectoriel fournit donc \emph{gratuitement} les trois premiers coefficients de l'équation : c'est l'outil standard pour déterminer l'équation d'un plan passant par trois points.
\end{aretenir}

\begin{attention}[La formule exige une base orthonormée DIRECTE]
La formule $\vec{u} \wedge \vec{v} = (yz' - zy'\,;\, zx' - xz'\,;\, xy' - yx')$ n'est valable que si la base est à la fois \textbf{orthonormée} et \textbf{directe}.
\begin{itemize}
\item Si la base est orthonormée mais \emph{indirecte}, le résultat obtenu est l'opposé du vrai produit vectoriel (le sens est inversé).
\item Si la base n'est pas orthonormée, la formule est tout simplement fausse.
\end{itemize}
Au Baccalauréat, l'énoncé précise presque toujours « repère orthonormé direct $(O\,; \vec{i}, \vec{j}, \vec{k})$ ». Si ce n'est pas le cas, vérifie l'orientation avant d'appliquer la formule : c'est une condition d'application, pas un détail.
\end{attention}

\begin{savaistu}[Pourquoi « directe » ? Une histoire de vis]
Le choix de l'orientation directe est une convention, adoptée internationalement : elle correspond au sens de vissage d'une vis ordinaire (ou d'un tire-bouchon) quand on tourne de $\vec{u}$ vers $\vec{v}$. Les physiciens utilisent la même convention pour le champ magnétique créé par un courant — c'est la fameuse règle du bonhomme d'Ampère vue en première partie de la leçon. En pratique, c'est ce produit vectoriel qui se cache derrière le calcul des moments de forces dans les engins : du tournevis du mécanicien de Douala au rotor des générateurs du barrage de Song Loulou.
\end{savaistu}

\begin{exercice}[À toi de jouer]
Dans un repère orthonormé direct, on donne les points $A(1\,; -1\,; 2)$, $B(3\,; 0\,; 1)$ et $C(2\,; 2\,; -1)$.
\begin{enumerate}
\item Calculer $\overrightarrow{AB} \wedge \overrightarrow{AC}$ et vérifier l'orthogonalité par deux produits scalaires.
\item En déduire une équation cartésienne du plan $(ABC)$.
\item Le point $D(0\,; 1\,; 3)$ appartient-il au plan $(ABC)$ ?
\end{enumerate}
\end{exercice}

\begin{corrige}[Exercice 1]
\textbf{1.} $\overrightarrow{AB}(2\,; 1\,; -1)$ et $\overrightarrow{AC}(1\,; 3\,; -3)$.
\begin{align*}
\overrightarrow{AB} \wedge \overrightarrow{AC} = \begin{pmatrix} 1 \times (-3) - (-1) \times 3 \\ (-1) \times 1 - 2 \times (-3) \\ 2 \times 3 - 1 \times 1 \end{pmatrix} = \begin{pmatrix} -3 + 3 \\ -1 + 6 \\ 6 - 1 \end{pmatrix} = \begin{pmatrix} 0 \\ 5 \\ 5 \end{pmatrix}.
\end{align*}
Contrôle : $0 \times 2 + 5 \times 1 + 5 \times (-1) = 0$ \checkmark{} ; $0 \times 1 + 5 \times 3 + 5 \times (-3) = 0$ \checkmark. On peut simplifier et prendre $\vec{n}(0\,; 1\,; 1)$ comme vecteur normal.

\textbf{2.} Équation : $y + z + d = 0$. Avec $A(1\,; -1\,; 2)$ : $-1 + 2 + d = 0$ donne $d = -1$. Donc $(ABC) : y + z - 1 = 0$.

\textbf{3.} Pour $D(0\,; 1\,; 3)$ : $1 + 3 - 1 = 3 \neq 0$. Donc $D \notin (ABC)$.
\end{corrige}

Tu sais désormais calculer un produit vectoriel de manière purement algébrique, le vérifier, et l'exploiter pour trouver un vecteur normal et l'équation d'un plan. Dans la section suivante, nous étudierons les \emph{propriétés algébriques} du produit vectoriel — antisymétrie, bilinéarité, caractérisation de la colinéarité — qui justifieront a posteriori tous les calculs que nous venons de pratiquer.

\coursec{Propriétés algébriques du produit vectoriel}

Dans la section précédente, nous avons établi l'expression analytique du produit vectoriel dans une base orthonormée directe $(\vec{i},\vec{j},\vec{k})$ : pour $\vec{u}(x\,;y\,;z)$ et $\vec{v}(x'\,;y'\,;z')$,
\[
\vec{u}\wedge\vec{v}\;=\;\begin{pmatrix} yz'-zy' \\ zx'-xz' \\ xy'-yx' \end{pmatrix}.
\]
Cette formule est un outil de calcul formidable. Mais pour l'utiliser avec efficacité et sans erreur, il faut connaître les \emph{règles du jeu} : comment le produit vectoriel se comporte-t-il vis-à-vis de l'addition, de la multiplication par un scalaire, de l'ordre des vecteurs ? C'est l'objet de cette section. Tu vas découvrir que le produit vectoriel possède des propriétés algébriques riches, mais aussi quelques comportements \emph{déconcertants} par rapport à la multiplication ordinaire : il change de signe quand on permute les facteurs, il n'est pas associatif, et on ne peut pas « diviser » par un vecteur. Chacune de ces particularités est une source classique d'erreurs au Baccalauréat : prenons le temps de les comprendre en profondeur.

\courssub{Antisymétrie du produit vectoriel}

Commençons par une expérience de pensée. Prends deux vecteurs $\vec{u}$ et $\vec{v}$ non colinéaires. Le vecteur $\vec{u}\wedge\vec{v}$ est orthogonal au plan $(\vec{u},\vec{v})$, et son sens est donné par la règle de la main droite : pouce selon $\vec{u}$, index selon $\vec{v}$, le majeur indique $\vec{u}\wedge\vec{v}$. Maintenant, permute les rôles : pouce selon $\vec{v}$, index selon $\vec{u}$. Ta main est obligée de se retourner, et ton majeur pointe \emph{en sens opposé} ! La norme, elle, ne change pas : $\|\vec{u}\|\,\|\vec{v}\|\,|\sin(\vec{u},\vec{v})|$ est symétrique en $\vec{u}$ et $\vec{v}$. Conclusion : $\vec{v}\wedge\vec{u}$ et $\vec{u}\wedge\vec{v}$ ont même direction, même norme, mais des sens opposés. C'est exactement dire qu'ils sont opposés.

\begin{propriete}[Antisymétrie du produit vectoriel]
Pour tous vecteurs $\vec{u}$ et $\vec{v}$ de l'espace :
\[
\vec{v}\wedge\vec{u}=-\bigl(\vec{u}\wedge\vec{v}\bigr).
\]
On dit que le produit vectoriel est \cle{antisymétrique}.
\end{propriete}

\begin{experience}[Démonstration de l'antisymétrie]
Démontrons cette propriété à partir de la définition géométrique.

\textbf{Cas 1 : $\vec{u}$ et $\vec{v}$ sont colinéaires.} Alors $\vec{u}\wedge\vec{v}=\vec{0}$ et $\vec{v}\wedge\vec{u}=\vec{0}$, donc $\vec{v}\wedge\vec{u}=-(\vec{u}\wedge\vec{v})=\vec{0}$ : l'égalité est vraie.

\textbf{Cas 2 : $\vec{u}$ et $\vec{v}$ ne sont pas colinéaires.} Comparons les trois caractéristiques des vecteurs $\vec{u}\wedge\vec{v}$ et $\vec{v}\wedge\vec{u}$.

\emph{Direction.} Par définition, chacun des deux est orthogonal à la fois à $\vec{u}$ et à $\vec{v}$, donc orthogonal au plan $(\vec{u},\vec{v})$. Ils ont donc la \textbf{même direction} (la normale à ce plan).

\emph{Norme.} On a
\[
\|\vec{u}\wedge\vec{v}\|=\|\vec{u}\|\,\|\vec{v}\|\,\bigl|\sin(\vec{u},\vec{v})\bigr|
=\|\vec{v}\|\,\|\vec{u}\|\,\bigl|\sin(\vec{v},\vec{u})\bigr|=\|\vec{v}\wedge\vec{u}\|,
\]
car $|\sin|$ prend la même valeur pour les angles $(\vec{u},\vec{v})$ et $(\vec{v},\vec{u})$. Les normes sont \textbf{égales}.

\emph{Sens.} Le trièdre $(\vec{u},\vec{v},\vec{u}\wedge\vec{v})$ est direct par définition. Or permuter deux vecteurs d'un trièdre direct le transforme en trièdre \emph{indirect} (règle du bonhomme d'Ampère : en échangeant $\vec{u}$ et $\vec{v}$, l'observateur doit se placer « tête en bas »). Donc le trièdre $(\vec{v},\vec{u},\vec{u}\wedge\vec{v})$ est indirect, ce qui signifie que $\vec{v}\wedge\vec{u}$ pointe dans le sens \textbf{opposé} à $\vec{u}\wedge\vec{v}$.

\emph{Conclusion.} Les deux vecteurs ont même direction, même norme et des sens opposés : $\vec{v}\wedge\vec{u}=-(\vec{u}\wedge\vec{v})$. $\blacksquare$
\end{experience}

\begin{popfigure}
\begin{tikzpicture}[scale=1.05,>=Stealth]
% Plan
\fill[popblueL,opacity=0.55] (-2.6,-1.3) -- (2.2,-1.3) -- (3.1,0.9) -- (-1.7,0.9) -- cycle;
\draw[popmuted] (-2.6,-1.3) -- (2.2,-1.3) -- (3.1,0.9) -- (-1.7,0.9) -- cycle;
\node[popmuted] at (2.6,0.6) {\small plan $(\vec{u},\vec{v})$};
% Vecteurs u et v
\draw[->,very thick,popink] (0,0) -- (2.1,0.35) node[right] {$\vec{u}$};
\draw[->,very thick,poporange] (0,0) -- (0.9,0.75) node[above right] {$\vec{v}$};
% u wedge v (vers le haut)
\draw[->,ultra thick,popblue] (0,0) -- (0,2.3) node[above] {$\vec{u}\wedge\vec{v}$};
% v wedge u (vers le bas)
\draw[->,ultra thick,poppurple] (0,0) -- (0,-2.3) node[below] {$\vec{v}\wedge\vec{u}=-(\vec{u}\wedge\vec{v})$};
% Angle droit
\draw[popblue] (0,0.35) -- (-0.28,0.35) -- (-0.28,0.02);
\end{tikzpicture}
\legende{Permuter $\vec{u}$ et $\vec{v}$ renverse le produit vectoriel : $\vec{u}\wedge\vec{v}$ et $\vec{v}\wedge\vec{u}$ sont deux vecteurs opposés, de part et d'autre du plan $(\vec{u},\vec{v})$.}
\end{popfigure}

Une conséquence immédiate mais capitale :

\begin{propriete}[Produit vectoriel d'un vecteur par lui-même]
Pour tout vecteur $\vec{u}$ de l'espace :
\[
\vec{u}\wedge\vec{u}=\vec{0}.
\]
\end{propriete}

En effet, par antisymétrie, $\vec{u}\wedge\vec{u}=-(\vec{u}\wedge\vec{u})$, donc $2(\vec{u}\wedge\vec{u})=\vec{0}$, d'où $\vec{u}\wedge\vec{u}=\vec{0}$. On retrouve bien sûr la définition : un vecteur est colinéaire à lui-même, donc le produit vectoriel est nul. Mais la démonstration par l'antisymétrie est plus élégante : le vecteur nul est le seul vecteur égal à son opposé.

\begin{attention}[L'ordre des facteurs compte !]
Contrairement au produit scalaire qui est \emph{symétrique} ($\vec{u}\cdot\vec{v}=\vec{v}\cdot\vec{u}$), le produit vectoriel est \emph{antisymétrique}. Quand tu permutes les deux vecteurs dans un calcul, n'oublie \textbf{jamais} le signe moins. C'est l'erreur la plus fréquente dans les calculs d'aires et de vecteurs normaux : un signe perdu en route fausse tout le résultat.
\end{attention}

\courssub{Bilinéarité du produit vectoriel}

Le produit vectoriel se comporte « gentiment » vis-à-vis de l'addition et de la multiplication par un scalaire, \emph{dans chacune de ses deux variables}. C'est ce qu'on appelle la \cle{bilinéarité}. Concrètement, cela signifie que tu peux développer un produit vectoriel comme tu développerais un produit ordinaire $(a+b)(c+d)$, à condition de respecter scrupuleusement l'ordre des facteurs.

\begin{propriete}[Bilinéarité du produit vectoriel]
Pour tous vecteurs $\vec{u}$, $\vec{u}'$, $\vec{v}$, $\vec{v}'$ de l'espace et tout réel $\lambda$ :

\textbf{Linéarité à gauche :}
\[
(\vec{u}+\vec{u}')\wedge\vec{v}=\vec{u}\wedge\vec{v}+\vec{u}'\wedge\vec{v}
\qquad\text{et}\qquad
(\lambda\vec{u})\wedge\vec{v}=\lambda\,(\vec{u}\wedge\vec{v}).
\]

\textbf{Linéarité à droite :}
\[
\vec{u}\wedge(\vec{v}+\vec{v}')=\vec{u}\wedge\vec{v}+\vec{u}\wedge\vec{v}'
\qquad\text{et}\qquad
\vec{u}\wedge(\lambda\vec{v})=\lambda\,(\vec{u}\wedge\vec{v}).
\]
\end{propriete}

\begin{experience}[Démonstration de la bilinéarité par les coordonnées]
Plaçons-nous dans une base orthonormée directe et notons $\vec{u}(x\,;y\,;z)$, $\vec{u}'(x'\,;y'\,;z')$ et $\vec{v}(a\,;b\,;c)$. Démontrons la distributivité à gauche ; les autres égalités se démontrent de la même manière.

Le vecteur $\vec{u}+\vec{u}'$ a pour coordonnées $(x+x'\,;y+y'\,;z+z')$. La première coordonnée de $(\vec{u}+\vec{u}')\wedge\vec{v}$ est donc :
\begin{align*}
(y+y')c-(z+z')b
&= (yc-zb)+(y'c-z'b)\\
&= \text{(1\textsuperscript{re} coordonnée de }\vec{u}\wedge\vec{v})+\text{(1\textsuperscript{re} coordonnée de }\vec{u}'\wedge\vec{v}).
\end{align*}
Le même calcul sur les deuxième et troisième coordonnées montre que $(\vec{u}+\vec{u}')\wedge\vec{v}$ et $\vec{u}\wedge\vec{v}+\vec{u}'\wedge\vec{v}$ ont les mêmes coordonnées : ils sont égaux.

Pour l'homogénéité, $\lambda\vec{u}$ a pour coordonnées $(\lambda x\,;\lambda y\,;\lambda z)$, et la première coordonnée de $(\lambda\vec{u})\wedge\vec{v}$ est
\[
(\lambda y)c-(\lambda z)b=\lambda(yc-zb),
\]
c'est-à-dire $\lambda$ fois la première coordonnée de $\vec{u}\wedge\vec{v}$ ; idem pour les autres. Donc $(\lambda\vec{u})\wedge\vec{v}=\lambda(\vec{u}\wedge\vec{v})$.

Enfin, la linéarité à droite se déduit de la linéarité à gauche grâce à l'antisymétrie :
\[
\vec{u}\wedge(\vec{v}+\vec{v}')=-\bigl((\vec{v}+\vec{v}')\wedge\vec{u}\bigr)
=-\bigl(\vec{v}\wedge\vec{u}+\vec{v}'\wedge\vec{u}\bigr)
=\vec{u}\wedge\vec{v}+\vec{u}\wedge\vec{v}'.
\quad\blacksquare
\]
\end{experience}

\begin{exemplebox}[Un développement remarquable : $(\vec{u}+\vec{v})\wedge(\vec{u}-\vec{v})=-2(\vec{u}\wedge\vec{v})$]
Développons pas à pas, en veillant à l'ordre des facteurs :
\begin{align*}
(\vec{u}+\vec{v})\wedge(\vec{u}-\vec{v})
&= \vec{u}\wedge\vec{u}-\vec{u}\wedge\vec{v}+\vec{v}\wedge\vec{u}-\vec{v}\wedge\vec{v}\\
&= \vec{0}-\vec{u}\wedge\vec{v}-\bigl(\vec{u}\wedge\vec{v}\bigr)-\vec{0}
\quad\text{(car }\vec{v}\wedge\vec{u}=-\vec{u}\wedge\vec{v}\text{)}\\
&= -2\,(\vec{u}\wedge\vec{v}).
\end{align*}
\textbf{Application numérique rapide.} Soit $\vec{u}(1\,;2\,;-1)$ et $\vec{v}(3\,;0\,;2)$. On demande $(\vec{u}+\vec{v})\wedge(\vec{u}-\vec{v})$. Plutôt que de calculer $\vec{u}+\vec{v}(4\,;2\,;1)$ et $\vec{u}-\vec{v}(-2\,;2\,;-3)$ puis leur produit vectoriel, on calcule une seule fois :
\[
\vec{u}\wedge\vec{v}
=\begin{pmatrix} 2\times2-(-1)\times0 \\ (-1)\times3-1\times2 \\ 1\times0-2\times3 \end{pmatrix}
=\begin{pmatrix} 4 \\ -5 \\ -6 \end{pmatrix},
\qquad\text{donc}\qquad
(\vec{u}+\vec{v})\wedge(\vec{u}-\vec{v})=\begin{pmatrix} -8 \\ 10 \\ 12 \end{pmatrix}.
\]
Deux fois moins de calculs, deux fois moins de risques d'erreur !
\end{exemplebox}

\courssub{Caractérisation de la colinéarité}

Voici le théorème central de cette section, celui qui justifie à lui seul l'existence du produit vectoriel dans le programme : il fournit un \emph{critère calculatoire} de colinéarité, et donc un moyen de tester l'alignement de trois points ou le parallélisme de deux droites.

\begin{propriete}[Théorème : caractérisation de la colinéarité]
Soit $\vec{u}$ et $\vec{v}$ deux vecteurs de l'espace. Alors :
\[
\vec{u}\ \text{et}\ \vec{v}\ \text{sont colinéaires}
\iff
\vec{u}\wedge\vec{v}=\vec{0}.
\]
Cette démonstration est \textbf{exigible} au Baccalauréat.
\end{propriete}

\begin{experience}[Démonstration complète du théorème]
Démontrons les deux implications.

\textbf{Sens direct ($\Rightarrow$).} Supposons $\vec{u}$ et $\vec{v}$ colinéaires. C'est immédiat par la définition géométrique : le produit vectoriel de deux vecteurs colinéaires est le vecteur nul. Retrouvons-le aussi par la bilinéarité, ce qui sera utile pour la suite. Si $\vec{u}=\vec{0}$, alors $\vec{u}\wedge\vec{v}=\vec{0}$. Sinon, il existe un réel $\lambda$ tel que $\vec{v}=\lambda\vec{u}$, et :
\[
\vec{u}\wedge\vec{v}=\vec{u}\wedge(\lambda\vec{u})=\lambda\,(\vec{u}\wedge\vec{u})=\lambda\,\vec{0}=\vec{0}.
\]

\textbf{Réciproque ($\Leftarrow$).} Supposons $\vec{u}\wedge\vec{v}=\vec{0}$. Si $\vec{u}=\vec{0}$, alors $\vec{u}$ est colinéaire à tout vecteur, donc à $\vec{v}$ : c'est terminé. Supposons donc $\vec{u}\neq\vec{0}$ et notons $\vec{u}(x\,;y\,;z)$ et $\vec{v}(x'\,;y'\,;z')$ dans une base orthonormée directe. L'hypothèse s'écrit :
\[
yz'-zy'=0,\qquad zx'-xz'=0,\qquad xy'-yx'=0. \tag{$\star$}
\]
Comme $\vec{u}\neq\vec{0}$, l'une de ses coordonnées est non nulle ; supposons par exemple $x\neq0$ (le raisonnement est identique avec $y$ ou $z$). Posons $\lambda=\dfrac{x'}{x}$, de sorte que $x'=\lambda x$. Reportons dans les deux dernières relations de $(\star)$ :
\[
zx'-xz'=0 \implies z(\lambda x)-xz'=0 \implies xz'=\lambda xz \implies z'=\lambda z
\quad(\text{car }x\neq0),
\]
\[
xy'-yx'=0 \implies xy'-y(\lambda x)=0 \implies xy'=\lambda xy \implies y'=\lambda y.
\]
Ainsi $(x'\,;y'\,;z')=(\lambda x\,;\lambda y\,;\lambda z)$, c'est-à-dire $\vec{v}=\lambda\vec{u}$ : les vecteurs sont colinéaires. $\blacksquare$
\end{experience}

\begin{methode}[Tester la colinéarité de deux vecteurs de l'espace]
\begin{enumerate}
\item Écris les coordonnées de $\vec{u}$ et $\vec{v}$ dans une base orthonormée directe.
\item Calcule $\vec{u}\wedge\vec{v}$ par la formule analytique.
\item Conclus : si les trois coordonnées sont nulles, les vecteurs sont colinéaires ; si l'une au moins est non nulle, ils ne le sont pas.
\end{enumerate}
\emph{Application à l'alignement :} trois points $A$, $B$, $C$ sont alignés si et seulement si $\overrightarrow{AB}\wedge\overrightarrow{AC}=\vec{0}$.
\end{methode}

\begin{exoresolu}[Alignement de trois points]
Dans un repère orthonormé direct $(O\,;\vec{i},\vec{j},\vec{k})$, on donne $A(1\,;0\,;2)$, $B(3\,;2\,;0)$ et $C(7\,;6\,;-4)$. Les points $A$, $B$ et $C$ sont-ils alignés ?
\tcblower
On calcule $\overrightarrow{AB}(2\,;2\,;-2)$ et $\overrightarrow{AC}(6\,;6\,;-6)$. Puis :
\[
\overrightarrow{AB}\wedge\overrightarrow{AC}
=\begin{pmatrix}
2\times(-6)-(-2)\times6\\
(-2)\times6-2\times(-6)\\
2\times6-2\times6
\end{pmatrix}
=\begin{pmatrix}
-12+12\\
-12+12\\
12-12
\end{pmatrix}
=\begin{pmatrix}0\\0\\0\end{pmatrix}.
\]
D'après le théorème de caractérisation, $\overrightarrow{AB}$ et $\overrightarrow{AC}$ sont colinéaires (on observe d'ailleurs que $\overrightarrow{AC}=3\,\overrightarrow{AB}$). Les points $A$, $B$ et $C$ sont donc \textbf{alignés}.
\end{exoresolu}

\courssub{Ce que le produit vectoriel n'est PAS : non-associativité et impossibilité de simplifier}

Tu connais bien la multiplication des nombres réels : elle est commutative, associative, possède un élément neutre ($1$), et on peut « simplifier » ($ab=ac$ avec $a\neq0$ implique $b=c$). Le produit vectoriel ne vérifie \textbf{aucune} de ces propriétés. Nous avons déjà vu qu'il n'est pas commutatif (il est antisymétrique). Voyons le reste.

\begin{attention}[Le produit vectoriel n'est pas associatif et on ne peut pas « diviser » par un vecteur]
\textbf{1. Non-associativité.} En général,
\[
(\vec{u}\wedge\vec{v})\wedge\vec{w}\ \neq\ \vec{u}\wedge(\vec{v}\wedge\vec{w}).
\]
Il est donc \textbf{interdit} d'écrire $\vec{u}\wedge\vec{v}\wedge\vec{w}$ sans parenthèses : l'expression n'a pas de sens !

\textbf{2. Pas d'élément neutre.} Il n'existe aucun vecteur $\vec{e}$ tel que $\vec{u}\wedge\vec{e}=\vec{u}$ pour tout $\vec{u}$. En effet, $\vec{u}\wedge\vec{e}$ est toujours orthogonal à $\vec{u}$ (ou nul) : il ne peut être égal à $\vec{u}$ que si $\vec{u}=\vec{0}$.

\textbf{3. Pas de simplification.} De $\vec{u}\wedge\vec{v}=\vec{u}\wedge\vec{w}$, on ne peut \textbf{pas} conclure $\vec{v}=\vec{w}$. En effet, par bilinéarité :
\[
\vec{u}\wedge\vec{v}=\vec{u}\wedge\vec{w}
\iff \vec{u}\wedge(\vec{v}-\vec{w})=\vec{0}
\iff \vec{u}\ \text{et}\ \vec{v}-\vec{w}\ \text{sont colinéaires},
\]
ce qui est beaucoup plus faible que $\vec{v}=\vec{w}$. Contre-exemple éclatant : pour tout vecteur $\vec{u}$ non nul, $\vec{u}\wedge\vec{u}=\vec{u}\wedge\vec{0}=\vec{0}$, et pourtant $\vec{u}\neq\vec{0}$.
\end{attention}

\begin{exemplebox}[Contre-exemple à l'associativité, avec les vecteurs de base]
Prenons $\vec{u}=\vec{i}$, $\vec{v}=\vec{i}$, $\vec{w}=\vec{j}$ dans la base orthonormée directe $(\vec{i},\vec{j},\vec{k})$. Rappelons les produits de base : $\vec{i}\wedge\vec{j}=\vec{k}$, $\vec{j}\wedge\vec{k}=\vec{i}$, $\vec{k}\wedge\vec{i}=\vec{j}$, et $\vec{i}\wedge\vec{i}=\vec{0}$.

\textbf{Calcul à gauche :}
\[
(\vec{i}\wedge\vec{i})\wedge\vec{j}=\vec{0}\wedge\vec{j}=\vec{0}.
\]

\textbf{Calcul à droite :}
\[
\vec{i}\wedge(\vec{i}\wedge\vec{j})=\vec{i}\wedge\vec{k}=-(\vec{k}\wedge\vec{i})=-\vec{j}.
\]

Comme $\vec{0}\neq-\vec{j}$, on a bien $(\vec{i}\wedge\vec{i})\wedge\vec{j}\neq\vec{i}\wedge(\vec{i}\wedge\vec{j})$. Le produit vectoriel n'est \textbf{pas associatif}.
\end{exemplebox}

\begin{popfigure}
\begin{tikzpicture}[scale=1.15,>=Stealth]
% Repère 3D
\draw[->,thick,popdark] (0,0) -- (2.6,0) node[right] {$\vec{i}$};
\draw[->,thick,popdark] (0,0) -- (-1.1,-1.1) node[below left] {$\vec{j}$};
\draw[->,thick,popdark] (0,0) -- (0,2.6) node[above] {$\vec{k}$};
\fill (0,0) circle (1.5pt) node[below right=-2pt] {$O$};
% Membre de gauche : résultat nul
\node[align=center,popmuted] at (0,-2.1) {\small $(\vec{i}\wedge\vec{i})\wedge\vec{j}=\vec{0}\wedge\vec{j}=\vec{0}$};
% Membre de droite
\begin{scope}[xshift=7.2cm]
\draw[->,thick,popdark] (0,0) -- (2.6,0) node[right] {$\vec{i}$};
\draw[->,thick,popdark] (0,0) -- (-1.1,-1.1) node[below left] {$\vec{j}$};
\draw[->,thick,popdark] (0,0) -- (0,2.6) node[above] {$\vec{k}$};
\fill (0,0) circle (1.5pt);
% i wedge j = k
\draw[->,very thick,popblue] (0.35,0) -- (0.35,2.1) node[midway,right] {\small $\vec{i}\wedge\vec{j}=\vec{k}$};
% i wedge k = -j
\draw[->,very thick,poporange] (0,0) -- (1.5,1.5) node[above right] {\small $\vec{i}\wedge\vec{k}=-\vec{j}$};
\node[align=center,popmuted] at (0.4,-2.1) {\small $\vec{i}\wedge(\vec{i}\wedge\vec{j})=\vec{i}\wedge\vec{k}=-\vec{j}\neq\vec{0}$};
\end{scope}
\end{tikzpicture}
\legende{Contre-exemple à l'associativité : à gauche, $(\vec{i}\wedge\vec{i})\wedge\vec{j}=\vec{0}$ ; à droite, $\vec{i}\wedge(\vec{i}\wedge\vec{j})=-\vec{j}$. Le placement des parenthèses change le résultat.}
\end{popfigure}

\begin{savaistu}[Pourquoi une opération aussi « mal élevée » ?]
Le produit vectoriel est antisymétrique et non associatif : à première vue, c'est un mauvais élève comparé à la multiplication des nombres ! Pourtant, ces « défauts » sont exactement ce qu'il faut pour décrire la physique de l'espace. Le moment d'une force (qui fait tourner une porte ou une clé à molette), la force de Laplace sur un câble électrique, le champ magnétique créé par un courant : tous ces phénomènes changent de sens quand on inverse le sens du courant ou de la force — exactement comme $\vec{u}\wedge\vec{v}$ change de signe quand on permute $\vec{u}$ et $\vec{v}$. Les mathématiciens disent que l'espace $\mathbb{R}^3$ muni du produit vectoriel forme une \emph{algèbre de Lie} ; cette structure est au cœur de la mécanique du solide et de la robotique moderne.
\end{savaistu}

\courssub{Entraînement : simplifier des expressions mêlant produit vectoriel et produit scalaire}

Les exercices de Baccalauréat adorent mélanger produit vectoriel et produit scalaire dans une même expression. Les outils à mobiliser sont : la bilinéarité (développer), l'antisymétrie (permuter avec un signe moins), le fait que $\vec{u}\wedge\vec{u}=\vec{0}$, et l'orthogonalité fondamentale : $\vec{u}\wedge\vec{v}$ est orthogonal à $\vec{u}$ et à $\vec{v}$, donc
\[
\vec{u}\cdot(\vec{u}\wedge\vec{v})=0
\qquad\text{et}\qquad
\vec{v}\cdot(\vec{u}\wedge\vec{v})=0.
\]

\begin{exoresolu}[Simplification d'expressions vectorielles]
Soit $\vec{u}$, $\vec{v}$, $\vec{w}$ trois vecteurs de l'espace.
\begin{enumerate}
\item Simplifier $\vec{A}=(2\vec{u}-\vec{v})\wedge(\vec{u}+3\vec{v})$.
\item Montrer que $B=\bigl(\vec{u}\wedge\vec{v}\bigr)\cdot\bigl(\vec{u}+\vec{v}\bigr)=0$.
\item Simplifier $\vec{C}=(\vec{u}\wedge\vec{v})\wedge\vec{v}+\vec{v}\wedge(\vec{u}\wedge\vec{v})$.
\end{enumerate}
\tcblower
\textbf{1.} On développe par bilinéarité :
\begin{align*}
\vec{A}
&= 2\,\vec{u}\wedge\vec{u}+6\,\vec{u}\wedge\vec{v}-\vec{v}\wedge\vec{u}-3\,\vec{v}\wedge\vec{v}\\
&= \vec{0}+6\,\vec{u}\wedge\vec{v}+\vec{u}\wedge\vec{v}-\vec{0}
\quad\text{(car }\vec{v}\wedge\vec{u}=-\vec{u}\wedge\vec{v}\text{)}\\
&= 7\,\vec{u}\wedge\vec{v}.
\end{align*}

\textbf{2.} Le vecteur $\vec{u}\wedge\vec{v}$ est orthogonal à $\vec{u}$ et à $\vec{v}$, donc $(\vec{u}\wedge\vec{v})\cdot\vec{u}=0$ et $(\vec{u}\wedge\vec{v})\cdot\vec{v}=0$. Par linéarité du produit scalaire :
\[
B=(\vec{u}\wedge\vec{v})\cdot\vec{u}+(\vec{u}\wedge\vec{v})\cdot\vec{v}=0+0=0.
\]

\textbf{3.} Par antisymétrie, $\vec{v}\wedge(\vec{u}\wedge\vec{v})=-(\vec{u}\wedge\vec{v})\wedge\vec{v}$. Donc :
\[
\vec{C}=(\vec{u}\wedge\vec{v})\wedge\vec{v}-(\vec{u}\wedge\vec{v})\wedge\vec{v}=\vec{0}.
\]
Les deux termes sont opposés : leur somme est le vecteur nul.
\end{exoresolu}

\begin{exercice}[À toi de jouer]
Soit $\vec{u}$ et $\vec{v}$ deux vecteurs non colinéaires de l'espace.
\begin{enumerate}
\item Développer et simplifier $(\vec{u}+2\vec{v})\wedge(3\vec{u}-\vec{v})$.
\item On donne $\vec{u}(1\,;-1\,;2)$ et $\vec{v}(2\,;1\,;0)$. Calculer $(\vec{u}+\vec{v})\wedge(\vec{u}-2\vec{v})$ en utilisant la question 1 comme modèle.
\item Montrer que $\bigl(\vec{u}\wedge\vec{v}\bigr)\cdot\bigl(2\vec{u}-5\vec{v}\bigr)=0$.
\end{enumerate}
\end{exercice}

\begin{corrige}[Exercice : À toi de jouer]
\textbf{1.} Par bilinéarité :
\begin{align*}
(\vec{u}+2\vec{v})\wedge(3\vec{u}-\vec{v})
&= 3\,\vec{u}\wedge\vec{u}-\vec{u}\wedge\vec{v}+6\,\vec{v}\wedge\vec{u}-2\,\vec{v}\wedge\vec{v}\\
&= \vec{0}-\vec{u}\wedge\vec{v}-6\,\vec{u}\wedge\vec{v}-\vec{0}
= -7\,\vec{u}\wedge\vec{v}.
\end{align*}

\textbf{2.} D'après le même schéma de développement :
\[
(\vec{u}+\vec{v})\wedge(\vec{u}-2\vec{v})=\vec{u}\wedge\vec{u}-2\,\vec{u}\wedge\vec{v}+\vec{v}\wedge\vec{u}-2\,\vec{v}\wedge\vec{v}=-3\,\vec{u}\wedge\vec{v}.
\]
Or
\[
\vec{u}\wedge\vec{v}
=\begin{pmatrix} (-1)\times0-2\times1 \\ 2\times2-1\times0 \\ 1\times1-(-1)\times2 \end{pmatrix}
=\begin{pmatrix} -2 \\ 4 \\ 3 \end{pmatrix},
\qquad\text{donc}\qquad
(\vec{u}+\vec{v})\wedge(\vec{u}-2\vec{v})=\begin{pmatrix} 6 \\ -12 \\ -9 \end{pmatrix}.
\]

\textbf{3.} Le vecteur $\vec{u}\wedge\vec{v}$ est orthogonal à $\vec{u}$ et à $\vec{v}$, donc à toute combinaison linéaire de $\vec{u}$ et $\vec{v}$, en particulier à $2\vec{u}-5\vec{v}$. Le produit scalaire est donc nul.
\end{corrige}

\begin{aretenir}[L'essentiel des propriétés algébriques]
\begin{itemize}
\item \textbf{Antisymétrie :} $\vec{v}\wedge\vec{u}=-(\vec{u}\wedge\vec{v})$ ; en particulier $\vec{u}\wedge\vec{u}=\vec{0}$.
\item \textbf{Bilinéarité :} on développe comme un produit ordinaire, en respectant l'ordre des facteurs : $(\vec{u}+\vec{u}')\wedge\vec{v}=\vec{u}\wedge\vec{v}+\vec{u}'\wedge\vec{v}$ et $(\lambda\vec{u})\wedge\vec{v}=\lambda(\vec{u}\wedge\vec{v})$.
\item \textbf{Colinéarité :} $\vec{u}$ et $\vec{v}$ colinéaires $\iff \vec{u}\wedge\vec{v}=\vec{0}$ (démonstration exigible).
\item \textbf{Pièges :} pas d'associativité (parenthèses obligatoires), pas d'élément neutre, pas de simplification : $\vec{u}\wedge\vec{v}=\vec{u}\wedge\vec{w}$ n'implique pas $\vec{v}=\vec{w}$.
\item \textbf{Orthogonalité :} $\vec{u}\cdot(\vec{u}\wedge\vec{v})=\vec{v}\cdot(\vec{u}\wedge\vec{v})=0$, outil précieux pour simplifier les expressions mixtes.
\end{itemize}
\end{aretenir}

Ces règles algébriques maîtrisées, nous sommes prêts pour les applications géométriques spectaculaires du produit vectoriel : le calcul des distances d'un point à une droite ou à un plan, que nous aborderons dans la section suivante.

\coursec{Distances : point-droite et point-plan}

Dans la section précédente, nous avons établi les propriétés algébriques du produit vectoriel : antisymétrie, bilinéarité, et ce critère fondamental selon lequel $\vec{u} \wedge \vec{v} = \vec{0}$ si et seulement si $\vec{u}$ et $\vec{v}$ sont colinéaires. Nous avons aussi vu que la norme $\|\vec{u} \wedge \vec{v}\|$ possède une interprétation géométrique remarquable : c'est l'\cle{aire du parallélogramme} construit sur $\vec{u}$ et $\vec{v}$. C'est précisément cette interprétation que nous allons maintenant exploiter pour résoudre deux problèmes classiques de l'espace, très fréquents au Baccalauréat : \emph{à quelle distance un point se trouve-t-il d'une droite ? d'un plan ?}

Imagine un technicien de CAMTEL qui doit raccorder une maison (un point $M$) à une ligne électrique existante (une droite $D$) : le câble le plus court part de $M$ perpendiculairement à la ligne. De même, un drone de livraison qui survole un terrain incliné (un plan $P$) doit connaître son altitude, c'est-à-dire sa distance au plan. Ces deux situations concrètes se ramènent à des calculs que le produit vectoriel rend remarquablement simples.

\courssub{Distance d'un point à une droite}

\textbf{Le problème et l'intuition.}\par
Soit $D$ une droite de l'espace passant par un point $A$ et dirigée par un vecteur non nul $\vec{u}$, et soit $M$ un point quelconque. La distance de $M$ à $D$, notée $d(M, D)$, est la \emph{plus courte} distance entre $M$ et un point de $D$ : c'est la longueur $MH$, où $H$ est le \cle{projeté orthogonal} de $M$ sur $D$.

L'idée géniale consiste à ne pas chercher explicitement le point $H$ (ce qui serait fastidieux), mais à utiliser l'aire du parallélogramme construit sur les vecteurs $\overrightarrow{AM}$ et $\vec{u}$. Cette aire peut s'écrire de deux façons :
\begin{itemize}
\item avec le produit vectoriel : $\mathcal{A} = \|\overrightarrow{AM} \wedge \vec{u}\|$ ;
\item avec la formule classique « base $\times$ hauteur », en prenant $\vec{u}$ comme base : la hauteur correspondante est exactement $MH$, donc $\mathcal{A} = \|\vec{u}\| \times MH$.
\end{itemize}
En égalisant les deux expressions, on obtient immédiatement $MH = \dfrac{\|\overrightarrow{AM} \wedge \vec{u}\|}{\|\vec{u}\|}$. C'est la formule attendue !

\begin{popfigure}
\begin{center}
\begin{tikzpicture}[scale=1.6, >=Stealth]
% Droite D
\draw[popline, thick] (-0.6,0) -- (5.6,0);
\node[popdark, below] at (5.4,0) {$D$};
% Points A et H
\fill[popblue] (0.8,0) circle (1.5pt) node[below left, popdark] {$A$};
\fill[popblue] (3.2,0) circle (1.5pt) node[below, popdark] {$H$};
% Point M
\fill[popink] (3.2,2.1) circle (1.5pt) node[above, popdark] {$M$};
% Vecteur u
\draw[->, poporange, very thick] (0.8,0) -- (2.3,0) node[midway, below, poporange] {$\vec{u}$};
% Vecteur AM
\draw[->, popblue, very thick] (0.8,0) -- (3.2,2.1) node[midway, above left, popblue] {$\overrightarrow{AM}$};
% Hauteur MH
\draw[popink, very thick, dashed] (3.2,2.1) -- (3.2,0) node[midway, right, popink] {$MH$};
% Angle droit
\draw[popink] (3.2,0.18) -- (3.02,0.18) -- (3.02,0);
% Parallélogramme
\draw[popgreen, thick, dashed] (3.2,2.1) -- (4.7,2.1);
\draw[popgreen, thick, dashed] (2.3,0) -- (4.7,2.1);
\fill[popgreenL, opacity=0.35] (0.8,0) -- (2.3,0) -- (4.7,2.1) -- (3.2,2.1) -- cycle;
\node[popgreen!60!black] at (3.0,1.15) {$\mathcal{A} = \|\overrightarrow{AM} \wedge \vec{u}\|$};
\end{tikzpicture}
\end{center}
\legende{Le parallélogramme construit sur $\overrightarrow{AM}$ et $\vec{u}$ : son aire vaut à la fois $\|\overrightarrow{AM} \wedge \vec{u}\|$ et $\|\vec{u}\| \times MH$, où $MH$ est la distance cherchée.}
\end{popfigure}

\begin{propriete}[Distance d'un point à une droite (théorème)]
Soit $D$ une droite passant par un point $A$ et dirigée par un vecteur non nul $\vec{u}$, et $M$ un point de l'espace. Alors :
\[ d(M, D) = \frac{\left\|\overrightarrow{AM} \wedge \vec{u}\,\right\|}{\|\vec{u}\|} \]
Cette distance est nulle si et seulement si $M \in D$.
\end{propriete}

\begin{experience}[Démonstration guidée du théorème]
Soit $H$ le projeté orthogonal de $M$ sur $D$. Nous allons établir la formule en comparant deux calculs de l'aire d'un même parallélogramme.
\begin{enumerate}
\item \textbf{Construis le parallélogramme.} À partir du point $A$, porte les vecteurs $\overrightarrow{AM}$ et $\vec{u}$ ; complète en un parallélogramme ayant pour côtés adjacents $\overrightarrow{AM}$ et $\vec{u}$.
\item \textbf{Première expression de l'aire.} Par l'interprétation géométrique du produit vectoriel, l'aire de ce parallélogramme est $\mathcal{A} = \|\overrightarrow{AM} \wedge \vec{u}\|$.
\item \textbf{Deuxième expression de l'aire.} Choisissons le côté porté par $\vec{u}$ comme base : sa longueur est $\|\vec{u}\|$. La hauteur relative à cette base est la distance perpendiculaire de $M$ à la droite $D$, c'est-à-dire $MH$. Donc $\mathcal{A} = \|\vec{u}\| \times MH$.
\item \textbf{Conclus.} L'égalité $\|\vec{u}\| \times MH = \|\overrightarrow{AM} \wedge \vec{u}\|$ donne, puisque $\vec{u} \neq \vec{0}$ :
\[ MH = \frac{\|\overrightarrow{AM} \wedge \vec{u}\|}{\|\vec{u}\|}. \]
\item \textbf{Cas particulier.} Si $M \in D$, alors $\overrightarrow{AM}$ et $\vec{u}$ sont colinéaires, donc $\overrightarrow{AM} \wedge \vec{u} = \vec{0}$ et $d(M, D) = 0$. Réciproquement, si la distance est nulle, le produit vectoriel est nul, donc $\overrightarrow{AM}$ est colinéaire à $\vec{u}$, donc $M \in D$.
\end{enumerate}
\end{experience}

\begin{methode}[Calculer la distance d'un point à une droite en trois étapes]
Pour calculer la distance de $M$ à la droite $D(A, \vec{u})$ :
\begin{enumerate}
\item \textbf{Former le vecteur $\overrightarrow{AM}$} : soustraire les coordonnées de $A$ à celles de $M$.
\item \textbf{Calculer le produit vectoriel $\overrightarrow{AM} \wedge \vec{u}$} puis sa norme (racine carrée de la somme des carrés des coordonnées).
\item \textbf{Diviser par $\|\vec{u}\|$} : le quotient obtenu est la distance cherchée. Simplifier la racine si possible.
\end{enumerate}
\end{methode}

\begin{exemplebox}[Distance de $M(1\,;\,2\,;\,3)$ à la droite $D(A, \vec{u})$]
L'espace est muni d'un repère orthonormé direct $(O\,;\vec{i},\vec{j},\vec{k})$. Soit $D$ la droite passant par $A(0\,;\,1\,;\,0)$ et dirigée par $\vec{u}(1\,;\,1\,;\,1)$. Calculons $d(M, D)$ pour $M(1\,;\,2\,;\,3)$.

\textbf{Étape 1 :} $\overrightarrow{AM} = (1-0\,;\,2-1\,;\,3-0) = (1\,;\,1\,;\,3)$.

\textbf{Étape 2 :} produit vectoriel :
\[ \overrightarrow{AM} \wedge \vec{u} = \begin{pmatrix} 1 \\ 1 \\ 3 \end{pmatrix} \wedge \begin{pmatrix} 1 \\ 1 \\ 1 \end{pmatrix} = \begin{pmatrix} 1 \times 1 - 3 \times 1 \\ 3 \times 1 - 1 \times 1 \\ 1 \times 1 - 1 \times 1 \end{pmatrix} = \begin{pmatrix} -2 \\ 2 \\ 0 \end{pmatrix} \]
d'où $\left\|\overrightarrow{AM} \wedge \vec{u}\right\| = \sqrt{(-2)^2 + 2^2 + 0^2} = \sqrt{8} = 2\sqrt{2}$.

\textbf{Étape 3 :} $\|\vec{u}\| = \sqrt{1^2+1^2+1^2} = \sqrt{3}$, donc :
\[ d(M, D) = \frac{2\sqrt{2}}{\sqrt{3}} = \frac{2\sqrt{6}}{3} \approx 1{,}63. \]
\end{exemplebox}

\begin{attention}[Les trois oublis classiques]
\begin{itemize}
\item \textbf{Diviser par la mauvaise norme.} On divise par $\|\vec{u}\|$ (le vecteur directeur de la droite), jamais par $\|\overrightarrow{AM}\|$. Vérifie que ton résultat est homogène à une longueur.
\item \textbf{Oublier que $\vec{u}$ doit être non nul.} La formule n'a de sens que si $\vec{u} \neq \vec{0}$, ce qui est garanti dès que $D$ est une vraie droite.
\item \textbf{Se tromper dans le produit vectoriel.} La formule des coordonnées est $(ay - bz\,;\, bz \ldots)$ — non ! Rappel : pour $(a\,;\,b\,;\,c) \wedge (a'\,;\,b'\,;\,c')$, on obtient $(bc' - cb'\,;\, ca' - ac'\,;\, ab' - ba')$. Prends le temps de poser le calcul en colonnes.
\end{itemize}
\end{attention}

\courssub{Distance d'un point à un plan}

\textbf{Le problème et l'intuition.}\par
Soit $P$ un plan passant par un point $A$ et de vecteur normal $\vec{n}$ (non nul), et $M$ un point de l'espace. La distance de $M$ à $P$ est la longueur $MH$, où $H$ est le projeté orthogonal de $M$ sur $P$. Géométriquement, $MH$ est la \emph{composante} du vecteur $\overrightarrow{AM}$ dans la direction de la normale $\vec{n}$. Or le produit scalaire mesure exactement cette composante : $\overrightarrow{AM} \cdot \vec{n} = \pm \, MH \times \|\vec{n}\|$, le signe dépendant du côté du plan où se trouve $M$. D'où la formule, déjà rencontrée avec le produit scalaire.

\begin{popfigure}
\begin{center}
\begin{tikzpicture}[scale=1.15, >=Stealth]
% Plan P (parallélogramme en perspective)
\fill[popblueL, opacity=0.5] (0,0) -- (4.5,0) -- (6,1.8) -- (1.5,1.8) -- cycle;
\draw[popblue, thick] (0,0) -- (4.5,0) -- (6,1.8) -- (1.5,1.8) -- cycle;
\node[popblue!70!black] at (5.3,0.45) {$P$};
% Point A dans le plan
\fill[popdark] (2.6,0.9) circle (1.5pt) node[below left, popdark] {$A$};
% Point H
\fill[popblue] (3.6,1.1) circle (1.5pt) node[below right, popdark] {$H$};
% Point M au-dessus
\fill[popink] (3.6,3.6) circle (1.5pt) node[above, popdark] {$M$};
% Segment MH
\draw[popink, very thick, dashed] (3.6,3.6) -- (3.6,1.1) node[midway, right, popink] {$MH$};
% Vecteur normal n depuis A
\draw[->, poporange, very thick] (2.6,0.9) -- (2.6,2.9) node[right, poporange] {$\vec{n}$};
% Vecteur AM
\draw[->, popgreen, very thick] (2.6,0.9) -- (3.6,3.6) node[midway, above right, popgreen!60!black] {$\overrightarrow{AM}$};
% Angle droit en H
\draw[popink] (3.6,1.32) -- (3.42,1.32) -- (3.42,1.1);
\end{tikzpicture}
\end{center}
\legende{La distance de $M$ au plan $P$ est la longueur de la projection orthogonale de $\overrightarrow{AM}$ sur la normale $\vec{n}$ : $MH = \dfrac{|\overrightarrow{AM} \cdot \vec{n}|}{\|\vec{n}\|}$.}
\end{popfigure}

\begin{propriete}[Distance d'un point à un plan]
Soit $P$ un plan passant par $A$ et de vecteur normal $\vec{n} \neq \vec{0}$, et $M$ un point de l'espace. Alors :
\[ d(M, P) = \frac{\left|\overrightarrow{AM} \cdot \vec{n}\,\right|}{\|\vec{n}\|} \]
En particulier, si $P$ admet pour équation cartésienne $ax + by + cz + d = 0$ dans un repère orthonormé, alors $\vec{n}(a\,;\,b\,;\,c)$ est normal à $P$ et, pour $M(x_0\,;\,y_0\,;\,z_0)$ :
\[ d(M, P) = \frac{|ax_0 + by_0 + cz_0 + d|}{\sqrt{a^2 + b^2 + c^2}}. \]
\end{propriete}

\textbf{Justification.}\par
Soit $H$ le projeté orthogonal de $M$ sur $P$. Comme $\overrightarrow{HM}$ est colinéaire à $\vec{n}$, on a $\left|\overrightarrow{HM} \cdot \vec{n}\right| = HM \times \|\vec{n}\|$. Or $\overrightarrow{AM} = \overrightarrow{AH} + \overrightarrow{HM}$, et $\overrightarrow{AH} \cdot \vec{n} = 0$ car $\overrightarrow{AH}$ est dans le plan, donc perpendiculaire à la normale. Ainsi $\overrightarrow{AM} \cdot \vec{n} = \overrightarrow{HM} \cdot \vec{n}$, et en prenant les valeurs absolues : $\left|\overrightarrow{AM} \cdot \vec{n}\right| = MH \times \|\vec{n}\|$, ce qui donne la formule. Pour le cas de l'équation cartésienne, si $A(x_A\,;\,y_A\,;\,z_A) \in P$, alors $ax_A + by_A + cz_A + d = 0$, et $\overrightarrow{AM} \cdot \vec{n} = a(x_0 - x_A) + b(y_0 - y_A) + c(z_0 - z_A) = ax_0 + by_0 + cz_0 + d$.

\textbf{Où intervient le produit vectoriel ?}\par
Très souvent, le plan $P$ n'est pas donné par une équation ni par un vecteur normal, mais par \emph{trois points} $A$, $B$, $C$ non alignés. C'est là que le produit vectoriel entre en jeu : les vecteurs $\overrightarrow{AB}$ et $\overrightarrow{AC}$ dirigent le plan, donc $\vec{n} = \overrightarrow{AB} \wedge \overrightarrow{AC}$ est un vecteur \cle{normal} au plan. On enchaîne alors avec le produit scalaire.

\begin{methode}[Distance d'un point à un plan défini par trois points]
Pour calculer la distance de $M$ au plan $P = (ABC)$ :
\begin{enumerate}
\item \textbf{Obtenir un vecteur normal} : calculer $\vec{n} = \overrightarrow{AB} \wedge \overrightarrow{AC}$. (Au passage, si $\vec{n} = \vec{0}$, les points sont alignés et ne définissent pas un plan !)
\item \textbf{Calculer le produit scalaire} $\overrightarrow{AM} \cdot \vec{n}$ et prendre sa valeur absolue.
\item \textbf{Diviser par $\|\vec{n}\|$} : le quotient est $d(M, P)$.
\end{enumerate}
Variante : on peut aussi, à l'étape 1, en déduire une équation cartésienne $ax + by + cz + d = 0$ du plan (en écrivant que $\vec{n}$ est normal et que $A \in P$), puis appliquer la formule $\dfrac{|ax_0+by_0+cz_0+d|}{\sqrt{a^2+b^2+c^2}}$.
\end{methode}

\begin{exoresolu}[Distance d'un point à un plan défini par trois points]
Dans un repère orthonormé direct, on donne $A(1\,;\,0\,;\,0)$, $B(0\,;\,1\,;\,0)$, $C(0\,;\,0\,;\,1)$ et $M(2\,;\,3\,;\,4)$. Déterminer la distance de $M$ au plan $(ABC)$.
\tcblower
\textbf{Étape 1 : vecteur normal par produit vectoriel.}
$\overrightarrow{AB} = (-1\,;\,1\,;\,0)$ et $\overrightarrow{AC} = (-1\,;\,0\,;\,1)$.
\[ \vec{n} = \overrightarrow{AB} \wedge \overrightarrow{AC} = \begin{pmatrix} 1 \times 1 - 0 \times 0 \\ 0 \times (-1) - (-1) \times 1 \\ (-1) \times 0 - 1 \times (-1) \end{pmatrix} = \begin{pmatrix} 1 \\ 1 \\ 1 \end{pmatrix}. \]
$\vec{n} \neq \vec{0}$, donc $A$, $B$, $C$ définissent bien un plan.

\textbf{Étape 2 : équation du plan (chaîne complète).} Le plan a une équation de la forme $x + y + z + d = 0$. Comme $A(1\,;\,0\,;\,0) \in P$ : $1 + 0 + 0 + d = 0$, donc $d = -1$. D'où $P : x + y + z - 1 = 0$.

\textbf{Étape 3 : distance.}
\[ d(M, P) = \frac{|2 + 3 + 4 - 1|}{\sqrt{1^2+1^2+1^2}} = \frac{8}{\sqrt{3}} = \frac{8\sqrt{3}}{3} \approx 4{,}62. \]
Vérification par la méthode directe : $\overrightarrow{AM} = (1\,;\,3\,;\,4)$, $\overrightarrow{AM} \cdot \vec{n} = 1+3+4 = 8$, $\|\vec{n}\| = \sqrt{3}$, donc $d(M,P) = \dfrac{8}{\sqrt{3}}$. Les deux méthodes concordent.
\end{exoresolu}

\begin{attention}[Valeur absolue et norme : les deux gardiens de la formule]
\begin{itemize}
\item \textbf{La valeur absolue au numérateur est obligatoire.} Un produit scalaire peut être négatif (il dépend du côté du plan où se situe $M$), alors qu'une distance est toujours positive. Écrire $d(M,P) = \dfrac{\overrightarrow{AM} \cdot \vec{n}}{\|\vec{n}\|}$ sans valeur absolue est une erreur de rédaction sanctionnée au Bac.
\item \textbf{On divise toujours par la norme du vecteur normal} (ou du vecteur directeur pour une droite). Si $\vec{n}$ n'est pas unitaire, oublier cette division fausse le résultat.
\item \textbf{Contrôle rapide :} si $M \in P$, le numérateur doit s'annuler. Teste mentalement ce cas pour valider ta formule.
\end{itemize}
\end{attention}

\begin{exercice}[À toi de jouer]
Dans un repère orthonormé direct, on donne $A(1\,;\,-1\,;\,2)$, $\vec{u}(2\,;\,0\,;\,-1)$ et $M(3\,;\,1\,;\,0)$.
\begin{enumerate}
\item Calculer la distance de $M$ à la droite $D(A, \vec{u})$.
\item Soit $P$ le plan passant par $A$ et de vecteur normal $\vec{n}(1\,;\,2\,;\,-2)$. Calculer $d(M, P)$.
\end{enumerate}
\end{exercice}

\begin{corrige}[Exercice]
\begin{enumerate}
\item $\overrightarrow{AM} = (2\,;\,2\,;\,-2)$. Produit vectoriel :
\[ \overrightarrow{AM} \wedge \vec{u} = \begin{pmatrix} 2 \times (-1) - (-2) \times 0 \\ (-2) \times 2 - 2 \times (-1) \\ 2 \times 0 - 2 \times 2 \end{pmatrix} = \begin{pmatrix} -2 \\ -2 \\ -4 \end{pmatrix}. \]
Norme : $\sqrt{4+4+16} = \sqrt{24} = 2\sqrt{6}$. Et $\|\vec{u}\| = \sqrt{4+0+1} = \sqrt{5}$. Donc $d(M, D) = \dfrac{2\sqrt{6}}{\sqrt{5}} = \dfrac{2\sqrt{30}}{5}$.
\item $\overrightarrow{AM} \cdot \vec{n} = 2 \times 1 + 2 \times 2 + (-2) \times (-2) = 2 + 4 + 4 = 10$ ; $\|\vec{n}\| = \sqrt{1+4+4} = 3$. Donc $d(M, P) = \dfrac{|10|}{3} = \dfrac{10}{3}$.
\end{enumerate}
\end{corrige}

\begin{aretenir}[Les deux formules à connaître par cœur]
\begin{itemize}
\item Distance à une \textbf{droite} $D(A, \vec{u})$ : $\quad d(M, D) = \dfrac{\left\|\overrightarrow{AM} \wedge \vec{u}\,\right\|}{\|\vec{u}\|}$ \quad (produit \emph{vectoriel}, pas de valeur absolue car une norme est positive).
\item Distance à un \textbf{plan} $P(A, \vec{n})$ : $\quad d(M, P) = \dfrac{\left|\overrightarrow{AM} \cdot \vec{n}\,\right|}{\|\vec{n}\|}$ \quad (produit \emph{scalaire}, valeur absolue indispensable).
\item Plan donné par trois points ? Le produit vectoriel $\overrightarrow{AB} \wedge \overrightarrow{AC}$ fournit le vecteur normal, et la chaîne de calcul est lancée.
\end{itemize}
\end{aretenir}

\begin{savaistu}[Pourquoi ces formules sont partout]
Les distances point-droite et point-plan sont au cœur de la géométrie algorithmique moderne. En infographie, pour déterminer si un personnage de jeu vidéo touche un mur, le moteur calcule en permanence des distances point-plan. En topographie — par exemple lors du bornage des terrains à Yaoundé ou de l'implantation des pylônes d'électrification rurale —, on calcule des distances d'un point relevé au GPS à des lignes de référence. La formule $\dfrac{\|\overrightarrow{AM} \wedge \vec{u}\|}{\|\vec{u}\|}$ est aussi la clé du calcul du \emph{moment d'une force} en physique : le moment par rapport à un axe, c'est précisément la force multipliée par la distance de son point d'application à l'axe, le « bras de levier » que tu utilises chaque fois que tu ouvres une porte en poussant loin des gonds !
\end{savaistu}

\coursec{Aires, produit mixte et volume du tétraèdre}

Dans la section précédente, nous avons exploité le produit vectoriel pour mesurer des \emph{distances} : celle d'un point à une droite, celle d'un point à un plan. Mais le produit vectoriel cache un trésor supplémentaire : sa \cle{norme} mesure une \emph{aire}. Ouvre grand les yeux : dès que tu sais calculer $\vec{u} \wedge \vec{v}$, tu sais calculer l'aire d'un triangle ou d'un parallélogramme dans l'espace, et avec une petite opération de plus — le \cle{produit mixte} — tu sauras calculer le \emph{volume} d'un tétraèdre et tester si quatre points sont coplanaires. C'est l'un des outils les plus rentables du Baccalauréat.

\courssub{Aire d'un parallélogramme et d'un triangle}

\textbf{Intuition.} Prends deux baguettes représentant les vecteurs $\vec{u}$ et $\vec{v}$ partant du même point. Elles déterminent un parallélogramme. Si les baguettes sont presque alignées, le parallélogramme est presque plat : son aire est petite. Si elles sont perpendiculaires, l'aire est maximale (pour des longueurs fixées). Or la norme du produit vectoriel est $\|\vec{u} \wedge \vec{v}\| = \|\vec{u}\|\,\|\vec{v}\|\,|\sin\theta|$ : elle contient exactement ce facteur $\sin\theta$ qui mesure « l'écartement » des deux vecteurs. Ce n'est pas un hasard.

\begin{propriete}[Aire du parallélogramme et du triangle]
Soient $A$, $B$, $C$, $D$ quatre points de l'espace.
\begin{itemize}
\item L'aire du \cle{parallélogramme} $ABCD$ (avec $\overrightarrow{AB}$ et $\overrightarrow{AD}$ comme côtés consécutifs) est :
\[ \mathcal{S}_{ABCD} = \left\| \overrightarrow{AB} \wedge \overrightarrow{AD} \right\|. \]
\item L'aire du \cle{triangle} $ABC$ est :
\[ \mathcal{S}_{ABC} = \frac{1}{2} \left\| \overrightarrow{AB} \wedge \overrightarrow{AC} \right\|. \]
\end{itemize}
\end{propriete}

\begin{experience}[Pourquoi la norme du produit vectoriel donne-t-elle l'aire ?]
Considérons le parallélogramme construit sur $\overrightarrow{AB}$ et $\overrightarrow{AD}$, et notons $\theta$ l'angle géométrique $\widehat{BAD}$ (compris entre $0$ et $\pi$).
\begin{enumerate}
\item L'aire d'un parallélogramme est « base $\times$ hauteur ». Prenons $\|\overrightarrow{AB}\|$ comme base. La hauteur est la distance de $D$ à la droite $(AB)$, qui vaut $\|\overrightarrow{AD}\| \sin\theta$.
\item Donc $\mathcal{S} = \|\overrightarrow{AB}\| \times \|\overrightarrow{AD}\| \sin\theta$.
\item Or, par définition du produit vectoriel : $\left\|\overrightarrow{AB} \wedge \overrightarrow{AD}\right\| = \|\overrightarrow{AB}\|\,\|\overrightarrow{AD}\|\,|\sin\theta|$, et comme $\theta \in [0 ; \pi]$, $\sin\theta \geq 0$. On retrouve exactement $\mathcal{S}$. \hfill$\blacksquare$
\end{enumerate}
Pour le triangle, il suffit de remarquer que le triangle $ABC$ est exactement la \emph{moitié} du parallélogramme construit sur $\overrightarrow{AB}$ et $\overrightarrow{AC}$ : d'où le facteur $\dfrac{1}{2}$.
\end{experience}

\begin{attention}[Le facteur $\dfrac{1}{2}$ oublié]
C'est l'erreur classique du Bac : pour un \textbf{triangle}, n'oublie jamais le facteur $\dfrac{1}{2}$. $\|\overrightarrow{AB} \wedge \overrightarrow{AC}\|$ donne l'aire du \emph{parallélogramme}, pas celle du triangle. Retiens : \textbf{parallélogramme : norme entière ; triangle : moitié de la norme.}
\end{attention}

\begin{exemplebox}[Aire d'un triangle dans l'espace]
Dans un repère orthonormé direct $(O ; \vec{i}, \vec{j}, \vec{k})$, on donne $A(1;0;0)$, $B(0;2;0)$ et $C(0;0;3)$. Calculons l'aire du triangle $ABC$.

On a $\overrightarrow{AB}\begin{pmatrix} -1 \\ 2 \\ 0 \end{pmatrix}$ et $\overrightarrow{AC}\begin{pmatrix} -1 \\ 0 \\ 3 \end{pmatrix}$. Alors :
\[ \overrightarrow{AB} \wedge \overrightarrow{AC} = \begin{pmatrix} 2 \times 3 - 0 \times 0 \\ 0 \times (-1) - (-1) \times 3 \\ (-1) \times 0 - 2 \times (-1) \end{pmatrix} = \begin{pmatrix} 6 \\ 3 \\ 2 \end{pmatrix}. \]
Sa norme vaut $\sqrt{6^2 + 3^2 + 2^2} = \sqrt{36+9+4} = \sqrt{49} = 7$. Donc :
\[ \mathcal{S}_{ABC} = \frac{1}{2} \times 7 = 3,5 \text{ unités d'aire.} \]
Bonus : au passage, le vecteur $\begin{pmatrix} 6 \\ 3 \\ 2 \end{pmatrix}$ est normal au plan $(ABC)$ — tu reconnais la technique de la section précédente.
\end{exemplebox}

\courssub{Le produit mixte de trois vecteurs}

\begin{savaistu}[Un complément utile au Bac]
Le \cle{produit mixte} ne figure pas toujours en toutes lettres dans les intitulés du programme, mais il est un \emph{complément très utile} au Baccalauréat : il condense en une seule formule le calcul des volumes et les tests de coplanarité. Les correcteurs le connaissent bien et son usage correct est toujours valorisé.
\end{savaistu}

\begin{definition}[Produit mixte de trois vecteurs (complément Bac)]
Soient $\vec{u}$, $\vec{v}$, $\vec{w}$ trois vecteurs de l'espace. On appelle \cle{produit mixte} de $\vec{u}$, $\vec{v}$, $\vec{w}$ (dans cet ordre) le \textbf{nombre réel} :
\[ [\vec{u}, \vec{v}, \vec{w}] = \left( \vec{u} \wedge \vec{v} \right) \cdot \vec{w}. \]
C'est le produit scalaire du vecteur $\vec{u} \wedge \vec{v}$ avec le vecteur $\vec{w}$.
\end{definition}

\begin{attention}[Un réel, pas un vecteur !]
Le produit mixte est un \textbf{scalaire} (un nombre), car il se termine par un produit \emph{scalaire}. Ne confonds pas : $\vec{u} \wedge \vec{v}$ est un vecteur, $(\vec{u} \wedge \vec{v}) \cdot \vec{w}$ est un réel. Il peut être positif, négatif ou nul — son signe dépend de l'orientation du triplet $(\vec{u}, \vec{v}, \vec{w})$.
\end{attention}

\textbf{Interprétation géométrique.} Construisons le parallélépipède dont trois arêtes issues d'un même sommet sont $\vec{u}$, $\vec{v}$, $\vec{w}$. Sa base est le parallélogramme construit sur $\vec{u}$ et $\vec{v}$, d'aire $\|\vec{u} \wedge \vec{v}\|$. Sa hauteur est la valeur absolue de la composante de $\vec{w}$ sur la direction normale à la base, c'est-à-dire sur $\vec{u} \wedge \vec{v}$. Or :
\[ (\vec{u} \wedge \vec{v}) \cdot \vec{w} = \|\vec{u} \wedge \vec{v}\| \times \text{(projection algébrique de } \vec{w} \text{ sur } \vec{u} \wedge \vec{v}\text{)} = \pm \, (\text{aire de base}) \times (\text{hauteur}). \]

\begin{propriete}[Volume du parallélépipède]
La valeur absolue du produit mixte est le volume du parallélépipède construit sur $\vec{u}$, $\vec{v}$, $\vec{w}$ :
\[ V_{\text{parallélépipède}} = \left| [\vec{u}, \vec{v}, \vec{w}] \right| = \left| (\vec{u} \wedge \vec{v}) \cdot \vec{w} \right|. \]
\end{propriete}

\begin{popfigure}
\begin{tikzpicture}[scale=1.15, line join=round]
\coordinate (O) at (0,0);
\coordinate (U) at (3.2,0);
\coordinate (V) at (1.3,1.5);
\coordinate (W) at (0.5,2.4);
\coordinate (UV) at ($(U)+(V)$);
\coordinate (UW) at ($(U)+(W)$);
\coordinate (VW) at ($(V)+(W)$);
\coordinate (UVW) at ($(U)+(V)+(W)$);
% faces
\fill[popblueL,opacity=0.55] (O)--(U)--(UV)--(V)--cycle;
\fill[poporangeL,opacity=0.45] (O)--(U)--(UW)--(W)--cycle;
\fill[poppurpleL,opacity=0.45] (U)--(UV)--(UVW)--(UW)--cycle;
% arêtes cachées
\draw[dashed,popmuted] (O)--(V);
\draw[dashed,popmuted] (V)--(UV);
\draw[dashed,popmuted] (V)--(VW);
% arêtes visibles
\draw[popdark,thick] (O)--(U)--(UV)--(UVW)--(VW)--(W)--cycle;
\draw[popdark,thick] (U)--(UW)--(UVW);
\draw[popdark,thick] (O)--(W);
\draw[popdark,thick] (UW)--(W);
\draw[popdark,thick] (VW)--(UV);
% vecteurs
\draw[-{Stealth},very thick,popblue] (O)--(U) node[midway,below] {$\vec{u}$};
\draw[-{Stealth},very thick,poporange] (O)--(V) node[midway,left,pos=0.35] {$\vec{v}$};
\draw[-{Stealth},very thick,poppurple] (O)--(W) node[midway,left] {$\vec{w}$};
% produit vectoriel
\draw[-{Stealth},very thick,popink] (O)--(1.9,3.4) node[above] {$\vec{u}\wedge\vec{v}$};
\fill[popdark] (O) circle (1.6pt) node[below left] {$O$};
\end{tikzpicture}
\legende{Parallélépipède construit sur trois vecteurs : son volume est $|(\vec{u}\wedge\vec{v})\cdot\vec{w}|$. La face colorée en bleu est la base, d'aire $\|\vec{u}\wedge\vec{v}\|$.}
\end{popfigure}

\courssub{Volume du tétraèdre}

\begin{propriete}[Volume du tétraèdre — démonstration exigible]
Soient $A$, $B$, $C$, $D$ quatre points non coplanaires de l'espace. Le volume du tétraèdre $ABCD$ est :
\[ V_{ABCD} = \frac{1}{6} \left| \left( \overrightarrow{AB} \wedge \overrightarrow{AC} \right) \cdot \overrightarrow{AD} \right|. \]
\end{propriete}

\begin{experience}[Démonstration guidée : $V = \dfrac{1}{3} \times \text{aire de base} \times \text{hauteur}$]
On admet (résultat du programme) que le volume d'un tétraèdre est $V = \dfrac{1}{3} \times \mathcal{B} \times h$, où $\mathcal{B}$ est l'aire d'une face et $h$ la hauteur correspondante.
\begin{enumerate}
\item Choisissons la face $ABC$ comme base : $\mathcal{B} = \mathcal{S}_{ABC} = \dfrac{1}{2} \left\| \overrightarrow{AB} \wedge \overrightarrow{AC} \right\|$.
\item La hauteur $h$ est la distance de $D$ au plan $(ABC)$. Notons $\vec{n} = \overrightarrow{AB} \wedge \overrightarrow{AC}$ : c'est un vecteur normal au plan $(ABC)$. D'après la formule de la distance d'un point à un plan (section précédente) :
\[ h = \frac{\left| \vec{n} \cdot \overrightarrow{AD} \right|}{\|\vec{n}\|}. \]
\item On assemble :
\begin{align*}
V &= \frac{1}{3} \times \mathcal{B} \times h = \frac{1}{3} \times \frac{1}{2}\|\vec{n}\| \times \frac{|\vec{n}\cdot\overrightarrow{AD}|}{\|\vec{n}\|} \\
&= \frac{1}{6} \left| \vec{n} \cdot \overrightarrow{AD} \right| = \frac{1}{6} \left| \left( \overrightarrow{AB} \wedge \overrightarrow{AC} \right) \cdot \overrightarrow{AD} \right|.
\end{align*}
Les normes $\|\vec{n}\|$ se simplifient : c'est la beauté du calcul. \hfill$\blacksquare$
\end{enumerate}
\end{experience}

\begin{attention}[Le facteur $\dfrac{1}{6}$, pas $\dfrac{1}{3}$ !]
Piège redoutable : la formule « $\dfrac{1}{3} \times$ base $\times$ hauteur » contient déjà un tiers, et l'aire du triangle de base apporte un $\dfrac{1}{2}$ supplémentaire. Au total : $\dfrac{1}{3} \times \dfrac{1}{2} = \dfrac{1}{6}$. Le volume du tétraèdre est donc le \textbf{sixième} de la valeur absolue du produit mixte, tandis que le parallélépipède construit sur les mêmes vecteurs a pour volume le produit mixte entier. Un tétraèdre, c'est exactement un sixième du parallélépipède.
\end{attention}

\begin{popfigure}
\begin{tikzpicture}[scale=1.5, line join=round]
\coordinate (A) at (0,0);
\coordinate (B) at (3.4,0);
\coordinate (C) at (1.6,1.9);
\coordinate (D) at (1.2,3.1);
\coordinate (H) at (1.2,0.85);
% base
\fill[popblueL,opacity=0.6] (A)--(B)--(C)--cycle;
\draw[popdark,thick] (A)--(B)--(C)--cycle;
% arêtes vers D
\draw[popdark,thick] (A)--(D);
\draw[popdark,thick] (B)--(D);
\draw[popdark,thick] (C)--(D);
% hauteur
\draw[dashed,thick,popink] (D)--(H) node[midway,right] {$h$};
\draw[popink] (H) ++(0.12,0) -- ++(0,0.12) -- ++(-0.12,0);
\fill[popink] (H) circle (1.3pt) node[below right,popink] {$H$};
% produit vectoriel
\draw[-{Stealth},very thick,poporange] (A)--(0.9,2.6) node[above left] {$\overrightarrow{AB}\wedge\overrightarrow{AC}$};
% sommets
\fill[popdark] (A) circle (1.5pt) node[below left] {$A$};
\fill[popdark] (B) circle (1.5pt) node[below right] {$B$};
\fill[popdark] (C) circle (1.5pt) node[right] {$C$};
\fill[popdark] (D) circle (1.5pt) node[above] {$D$};
\end{tikzpicture}
\legende{Tétraèdre $ABCD$ : la hauteur $h$ est la distance de $D$ au plan $(ABC)$, et $\overrightarrow{AB}\wedge\overrightarrow{AC}$ est normal à ce plan.}
\end{popfigure}

\courssub{Application : tester la coplanarité de quatre points}

Quatre points $A$, $B$, $C$, $D$ sont \cle{coplanaires} si et seulement si le tétraèdre $ABCD$ est « aplati », c'est-à-dire de volume nul. Or le volume est nul si et seulement si le produit mixte est nul. D'où la méthode :

\begin{methode}[Tester la coplanarité de quatre points par le produit mixte]
Pour savoir si $A$, $B$, $C$, $D$ sont coplanaires :
\begin{enumerate}
\item Calculer les coordonnées de $\overrightarrow{AB}$, $\overrightarrow{AC}$, $\overrightarrow{AD}$.
\item Calculer $\vec{n} = \overrightarrow{AB} \wedge \overrightarrow{AC}$.
\item Calculer le produit mixte $p = \vec{n} \cdot \overrightarrow{AD}$.
\item Conclure : $A, B, C, D$ sont \textbf{coplanaires} $\iff p = 0$. Si $p \neq 0$, ils forment un vrai tétraèdre de volume $\dfrac{|p|}{6}$.
\end{enumerate}
\end{methode}

\begin{exoresolu}[Aire, volume et coplanarité avec un même tétraèdre]
Dans un repère orthonormé direct, on donne $A(1;0;0)$, $B(0;2;0)$, $C(0;0;3)$ et $D(1;1;1)$.
\begin{enumerate}
\item Calculer l'aire du triangle $ABC$.
\item Calculer le volume du tétraèdre $ABCD$.
\item Le point $E(2;2;0)$ est-il coplanaire avec $A$, $B$ et $C$ ?
\end{enumerate}
\tcblower
\textbf{1. Aire de $ABC$.} Nous l'avons calculée plus haut : $\overrightarrow{AB}\begin{pmatrix} -1 \\ 2 \\ 0 \end{pmatrix}$, $\overrightarrow{AC}\begin{pmatrix} -1 \\ 0 \\ 3 \end{pmatrix}$, et
\[ \overrightarrow{AB} \wedge \overrightarrow{AC} = \begin{pmatrix} 6 \\ 3 \\ 2 \end{pmatrix}, \qquad \left\| \overrightarrow{AB} \wedge \overrightarrow{AC} \right\| = \sqrt{49} = 7. \]
Donc $\mathcal{S}_{ABC} = \dfrac{7}{2} = 3,5$ unités d'aire.

\textbf{2. Volume de $ABCD$.} On a $\overrightarrow{AD}\begin{pmatrix} 0 \\ 1 \\ 1 \end{pmatrix}$. Le produit mixte vaut :
\[ \left( \overrightarrow{AB} \wedge \overrightarrow{AC} \right) \cdot \overrightarrow{AD} = 6 \times 0 + 3 \times 1 + 2 \times 1 = 5. \]
Donc $V_{ABCD} = \dfrac{|5|}{6} = \dfrac{5}{6}$ unité de volume.

Vérifions avec la formule « $\dfrac{1}{3} \times$ base $\times$ hauteur » : la distance de $D$ au plan $(ABC)$ est $h = \dfrac{|5|}{7} = \dfrac{5}{7}$, donc $V = \dfrac{1}{3} \times \dfrac{7}{2} \times \dfrac{5}{7} = \dfrac{5}{6}$. Cohérent !

\textbf{3. Coplanarité de $A$, $B$, $C$, $E$.} On a $\overrightarrow{AE}\begin{pmatrix} 1 \\ 2 \\ 0 \end{pmatrix}$. Le produit mixte vaut :
\[ \left( \overrightarrow{AB} \wedge \overrightarrow{AC} \right) \cdot \overrightarrow{AE} = 6 \times 1 + 3 \times 2 + 2 \times 0 = 12 \neq 0. \]
Les quatre points $A$, $B$, $C$, $E$ ne sont \textbf{pas coplanaires} : ils forment un tétraèdre de volume $\dfrac{12}{6} = 2$.
\end{exoresolu}

\begin{exemplebox}[Un exemple camerounais : le hangar pyramidal]
Une coopérative agricole de Bafoussam veut construire un hangar en forme de tétraèdre (toiture pyramidale triangulaire). Dans un repère local (unité : le mètre), les quatre sommets sont $A(0;0;0)$, $B(6;0;0)$, $C(0;4;0)$ et $D(0;0;3)$. Quel volume de haricots peut-on stocker ?

$\overrightarrow{AB}\begin{pmatrix} 6 \\ 0 \\ 0 \end{pmatrix}$, $\overrightarrow{AC}\begin{pmatrix} 0 \\ 4 \\ 0 \end{pmatrix}$, $\overrightarrow{AD}\begin{pmatrix} 0 \\ 0 \\ 3 \end{pmatrix}$. Alors $\overrightarrow{AB} \wedge \overrightarrow{AC} = \begin{pmatrix} 0 \\ 0 \\ 24 \end{pmatrix}$, et le produit mixte vaut $0 \times 0 + 0 \times 0 + 24 \times 3 = 72$. D'où :
\[ V = \frac{72}{6} = \SI{12}{m^3}. \]
\end{exemplebox}

\begin{aretenir}[L'essentiel de la section]
\begin{itemize}
\item Aire du triangle : $\mathcal{S}_{ABC} = \dfrac{1}{2}\left\|\overrightarrow{AB} \wedge \overrightarrow{AC}\right\|$ ; aire du parallélogramme : $\mathcal{S} = \left\|\overrightarrow{AB} \wedge \overrightarrow{AD}\right\|$.
\item Produit mixte : $[\vec{u},\vec{v},\vec{w}] = (\vec{u} \wedge \vec{v}) \cdot \vec{w}$, un \textbf{réel}.
\item Volume du parallélépipède : $|[\vec{u},\vec{v},\vec{w}]|$ ; volume du tétraèdre : $V = \dfrac{1}{6}\left|(\overrightarrow{AB}\wedge\overrightarrow{AC})\cdot\overrightarrow{AD}\right|$.
\item $A, B, C, D$ coplanaires $\iff (\overrightarrow{AB}\wedge\overrightarrow{AC})\cdot\overrightarrow{AD} = 0$.
\end{itemize}
\end{aretenir}

\begin{exercice}[Pour t'entraîner]
On donne $A(1;1;0)$, $B(2;0;1)$, $C(0;1;2)$ et $D(3;3;3)$.
\begin{enumerate}
\item Calculer l'aire du triangle $ABC$.
\item Calculer le volume du tétraèdre $ABCD$.
\item Déterminer le réel $t$ tel que le point $M(t ; 1 ; 1)$ soit coplanaire avec $A$, $B$ et $C$.
\end{enumerate}
\end{exercice}

\begin{corrige}[Exercice]
\textbf{1.} $\overrightarrow{AB}\begin{pmatrix} 1 \\ -1 \\ 1 \end{pmatrix}$, $\overrightarrow{AC}\begin{pmatrix} -1 \\ 0 \\ 2 \end{pmatrix}$.
\[ \overrightarrow{AB} \wedge \overrightarrow{AC} = \begin{pmatrix} (-1)(2) - (1)(0) \\ (1)(-1) - (1)(2) \\ (1)(0) - (-1)(-1) \end{pmatrix} = \begin{pmatrix} -2 \\ -3 \\ -1 \end{pmatrix}, \quad \|\cdot\| = \sqrt{4+9+1} = \sqrt{14}. \]
Donc $\mathcal{S}_{ABC} = \dfrac{\sqrt{14}}{2}$.

\textbf{2.} $\overrightarrow{AD}\begin{pmatrix} 2 \\ 2 \\ 3 \end{pmatrix}$. Produit mixte : $(-2)(2) + (-3)(2) + (-1)(3) = -4 - 6 - 3 = -13$. Donc $V = \dfrac{|-13|}{6} = \dfrac{13}{6}$.

\textbf{3.} $\overrightarrow{AM}\begin{pmatrix} t-1 \\ 0 \\ 1 \end{pmatrix}$. Coplanarité $\iff (\overrightarrow{AB}\wedge\overrightarrow{AC})\cdot\overrightarrow{AM} = 0$ :
\[ (-2)(t-1) + (-3)(0) + (-1)(1) = -2t + 2 - 1 = -2t + 1 = 0 \iff t = \frac{1}{2}. \]
Le point cherché est $M\left(\dfrac{1}{2} ; 1 ; 1\right)$.
\end{corrige}

\newpage

\coursec{Activité d'intégration}

\coursec{Produit vectoriel}

\courssub{Orientation de l'espace}

\begin{definition}[Base orthonormée directe]
Soit $(\vec{i},\vec{j},\vec{k})$ une base orthonormée de l'espace. On dit qu'elle est \cle{directe} si le triplet $(\vec{i},\vec{j},\vec{k})$ satisfait la \cle{règle de la main droite} (ou règle du bonhomme d'Ampère) : si l'on plie les doigts de la main droite de $\vec{i}$ vers $\vec{j}$, le pouce indique le sens de $\vec{k}$.
Analytiquement, cela se traduit par $\det(\vec{i},\vec{j},\vec{k}) = +1$.
Si $\det(\vec{i},\vec{j},\vec{k}) = -1$, la base est \cle{indirecte} (main gauche).
\end{definition}

\begin{experience}[Repérer l'orientation]
Munis-toi de ta main droite. Place l'index selon $\vec{i}$, le majeur selon $\vec{j}$ (perpendiculaire à l'index), le pouce tendu donne le sens de $\vec{k}$.
\begin{itemize}
\item Base canonique $(\vec{i},\vec{j},\vec{k})$ : directe.
\item Base $(\vec{j},\vec{i},\vec{k})$ : indirecte (permuter deux vecteurs change l'orientation).
\end{itemize}
Toute formule de coordonnées du produit vectoriel suppose une base \textbf{orthonormée directe}. En base indirecte, le résultat est changé en son opposé.
\end{experience}

\courssub{Définition et interprétation géométrique du produit vectoriel}

\begin{definition}[Produit vectoriel]
Soit $(\vec{i},\vec{j},\vec{k})$ une base orthonormée \textbf{directe} de l'espace. Pour deux vecteurs $\vec{u}$ et $\vec{v}$, on appelle \cle{produit vectoriel} de $\vec{u}$ et $\vec{v}$, noté $\vec{u}\wedge\vec{v}$, l'unique vecteur $\vec{w}$ tel que :
\begin{enumerate}
\item $\vec{w}$ est orthogonal à $\vec{u}$ et à $\vec{v}$ (donc orthogonal au plan $(\vec{u},\vec{v})$ si $\vec{u}$ et $\vec{v}$ ne sont pas colinéaires) ;
\item $\|\vec{w}\| = \|\vec{u}\|\,\|\vec{v}\|\,|\sin(\vec{u},\vec{v})|$ (aire du parallélogramme construit sur $\vec{u}$ et $\vec{v}$) ;
\item $(\vec{u},\vec{v},\vec{w})$ est une base directe (si $\vec{u},\vec{v}$ non colinéaires).
\end{enumerate}
Si $\vec{u}$ et $\vec{v}$ sont colinéaires, $\vec{u}\wedge\vec{v}=\vec{0}$.
\end{definition}

\begin{propriete}[Interprétation géométrique]
\begin{itemize}
\item \textbf{Direction :} $\vec{u}\wedge\vec{v}$ est un vecteur normal au plan $(\vec{u},\vec{v})$.
\item \textbf{Norme :} $\|\vec{u}\wedge\vec{v}\|$ est l'aire du parallélogramme ayant $\vec{u}$ et $\vec{v}$ pour côtés adjacents.
\item \textbf{Aire d'un triangle :} Pour trois points $A,B,C$ non alignés, l'aire du triangle $ABC$ est $\mathscr{A}_{ABC} = \dfrac{1}{2}\left\|\overrightarrow{AB}\wedge\overrightarrow{AC}\right\|$.
\end{itemize}
\end{propriete}

\courssub{Expression analytique dans une base orthonormée directe}

\begin{propriete}[Coordonnées du produit vectoriel]
Dans une base orthonormée \textbf{directe} $(\vec{i},\vec{j},\vec{k})$, pour $\vec{u}\begin{pmatrix}x\\y\\z\end{pmatrix}$ et $\vec{v}\begin{pmatrix}x'\\y'\\z'\end{pmatrix}$ :
\[
\vec{u}\wedge\vec{v} = \begin{vmatrix}
\vec{i} & \vec{j} & \vec{k} \\
x & y & z \\
x' & y' & z'
\end{vmatrix}
= \begin{pmatrix}
yz' - zy' \\
zx' - xz' \\
xy' - yx'
\end{pmatrix}.
\]
\emph{Attention :} Cette formule (développement du déterminant formel) n'est valable \textbf{que} dans une base orthonormée directe.
\end{propriete}

\courssub{Propriétés algébriques du produit vectoriel}

\begin{propriete}[Propriétés fondamentales]
Soient $\vec{u},\vec{v},\vec{w}$ trois vecteurs de l'espace et $\lambda\in\mathbb{R}$.
\begin{enumerate}
\item \textbf{Antisymétrie :} $\vec{u}\wedge\vec{v} = -\,\vec{v}\wedge\vec{u}$.
\item \textbf{Bilinéarité :}
      \begin{align*}
      (\vec{u}+\vec{v})\wedge\vec{w} &= \vec{u}\wedge\vec{w} + \vec{v}\wedge\vec{w}, \\
      \vec{u}\wedge(\vec{v}+\vec{w}) &= \vec{u}\wedge\vec{v} + \vec{u}\wedge\vec{w}, \\
      (\lambda\vec{u})\wedge\vec{v} &= \vec{u}\wedge(\lambda\vec{v}) = \lambda(\vec{u}\wedge\vec{v}).
      \end{align*}
\item \textbf{Colinéarité :} $\vec{u}\wedge\vec{v}=\vec{0} \iff \vec{u}$ et $\vec{v}$ sont colinéaires.
\item \textbf{Non-associativité :} En général $(\vec{u}\wedge\vec{v})\wedge\vec{w} \neq \vec{u}\wedge(\vec{v}\wedge\vec{w})$.
\end{enumerate}
\end{propriete}

\courssub{Applications : distances, aires et volumes}

\begin{aretenir}[Formulaire de référence -- Produit vectoriel et produit mixte]
Dans une base orthonormée \textbf{directe} $(\vec{i},\vec{j},\vec{k})$ :
\begin{itemize}
\item \textbf{Produit vectoriel :} $\vec{u}(x,y,z),\ \vec{v}(x',y',z') \implies \vec{u}\wedge\vec{v}\begin{pmatrix} yz'-zy' \\ zx'-xz' \\ xy'-yx' \end{pmatrix}$.
\item \textbf{Vecteur normal à un plan :} Si $\vec{u},\vec{v}$ directeurs du plan, $\vec{n}=\vec{u}\wedge\vec{v}$ est un vecteur normal.
\item \textbf{Équation cartésienne d'un plan :} Plan passant par $A(x_A,y_A,z_A)$ de vecteur normal $\vec{n}(a,b,c)$ : $a(x-x_A)+b(y-y_A)+c(z-z_A)=0$, soit $ax+by+cz+d=0$ avec $d=-(ax_A+by_A+cz_A)$.
\item \textbf{Aire d'un triangle $ABC$ :} $\mathscr{A} = \dfrac{1}{2}\left\|\overrightarrow{AB}\wedge\overrightarrow{AC}\right\|$.
\item \textbf{Distance d'un point $M$ à un plan $(P)$ :} Si $(P): ax+by+cz+d=0$ et $M(x_0,y_0,z_0)$,
      \[ d(M,(P)) = \frac{|ax_0+by_0+cz_0+d|}{\sqrt{a^2+b^2+c^2}}. \]
\item \textbf{Distance d'un point $M$ à une droite $(AB)$ :}
      \[ d(M,(AB)) = \frac{\left\|\overrightarrow{AM}\wedge\overrightarrow{AB}\right\|}{\left\|\overrightarrow{AB}\right\|}. \]
\item \textbf{Produit mixte :} $[\vec{u},\vec{v},\vec{w}] = \vec{u}\cdot(\vec{v}\wedge\vec{w}) = \det(\vec{u},\vec{v},\vec{w})$.
      Volume du parallélépipède : $\mathscr{V}_{\text{parall}} = |[\vec{u},\vec{v},\vec{w}]|$.
\item \textbf{Volume d'un tétraèdre $OABC$ :} $\mathscr{V} = \dfrac{1}{6}\left|[\overrightarrow{OA},\overrightarrow{OB},\overrightarrow{OC}]\right|$.
\end{itemize}
\end{aretenir}

\courssub{Activité d'intégration : Le hangar de l'ONCPA}

\courssub{Objectifs de l'activité}
À l'issue de cette activité, tu seras capable de mobiliser, dans une situation concrète de construction, l'ensemble des savoir-faire de la leçon :
\begin{itemize}
\item justifier qu'une base de l'espace est \cle{orthonormée directe} ;
\item calculer un \cle{produit vectoriel} par ses coordonnées et interpréter géométriquement le résultat ;
\item déterminer un \cle{vecteur normal} à un plan et en déduire une \cle{équation cartésienne} de plan ;
\item calculer l'\cle{aire d'un triangle} de l'espace ;
\item calculer la \cle{distance d'un point à un plan} et la \cle{distance d'un point à une droite} ;
\item calculer le \cle{volume d'un tétraèdre} à l'aide du \cle{produit mixte}.
\end{itemize}
L'activité se déroule en groupes de trois ou quatre élèves, avec restitution au tableau et vérification croisée des résultats entre groupes.

\courssub{Situation}
L'Office National de Commercialisation des Produits de Base (ONCPA) fait construire, dans la plaine des Mbo (région de l'Ouest), un hangar destiné au stockage du maïs après la récolte. Le bureau d'études a modélisé la structure dans l'espace rapporté à un repère orthonormé $(O\,;\vec{i},\vec{j},\vec{k})$, dont l'unité sur chaque axe est le mètre. Le point $O$ est un coin du hangar au niveau du sol, l'axe $(O\,;\vec{k})$ est vertical, dirigé vers le haut, et le sol est le plan d'équation $z=0$.

La \cle{toiture} du hangar est une plaque triangulaire de tôle dont les trois sommets (le \emph{faîtage}) ont pour coordonnées :
\[ A(6\,;0\,;4), \qquad B(0\,;6\,;4), \qquad C(0\,;0\,;6). \]

Un poteau de soutien est planté verticalement au point $P(2\,;2\,;0)$ du sol. Pour dimensionner les matériaux, l'ingénieur doit répondre aux questions suivantes :
\begin{enumerate}
\item La base $(\vec{i},\vec{j},\vec{k})$ fournie par le logiciel de conception est-elle bien \emph{directe} ? (Cela conditionne la validité de toutes les formules de produit vectoriel utilisées ensuite.)
\item Quelle est l'inclinaison du toit ? Plus précisément, on cherche un vecteur normal au plan du toit $(ABC)$ et une équation cartésienne de ce plan.
\item Quelle quantité de tôle faut-il commander, c'est-à-dire quelle est l'aire du triangle $ABC$ ? (Le fournisseur de Bafoussam vend la tôle au mètre carré.)
\item Une entretoise verticale doit relier le pied du poteau $P$ au plan du toit : quelle est la distance de $P$ au plan $(ABC)$ ?
\item La gouttière suit l'arête $(AB)$ : à quelle distance du pied du poteau $P$ passe-t-elle, c'est-à-dire quelle est la distance de $P$ à la droite $(AB)$ ?
\item Le coin de stockage délimité par le tétraèdre $OABC$ sera rempli de sacs de maïs : quel volume (en $\mathrm{m}^3$) peut-il contenir ?
\end{enumerate}

\begin{popfigure}
\begin{center}
\begin{tikzpicture}[scale=0.9, line join=round]
\coordinate (O) at (0,0);
\coordinate (A) at (4.2,0.6);
\coordinate (B) at (0.8,3.4);
\coordinate (C) at (-1.6,1.4);
\coordinate (P) at (1.4,0.9);
\fill[poppurpleL,opacity=0.75] (A)--(B)--(C)--cycle;
\draw[thick,popdark] (A)--(B)--(C)--cycle;
\draw[thick,poporange] (O)--(A) (O)--(B) (O)--(C);
\fill[popink] (O) circle(1.6pt) node[below left]{$O$};
\fill[popink] (A) circle(1.6pt) node[right]{$A(6\,;0\,;4)$};
\fill[popink] (B) circle(1.6pt) node[above]{$B(0\,;6\,;4)$};
\fill[popink] (C) circle(1.6pt) node[left]{$C(0\,;0\,;6)$};
\fill[popgreen] (P) circle(1.6pt) node[below]{$P(2\,;2\,;0)$};
\draw[dashed,popgreen,thick] (P)--(1.9,2.1) node[midway,right]{\scriptsize entretoise};
\end{tikzpicture}
\end{center}
\legende{Schéma simplifié du hangar : la toiture est le triangle $ABC$, le coin de stockage est le tétraèdre $OABC$.}
\end{popfigure}

\courssub{Consignes}
Travail en groupe. Chaque question doit être rédigée soigneusement, avec la formule utilisée, le calcul détaillé et une phrase de conclusion interprétant le résultat pour l'ingénieur.

\begin{enumerate}
\item \textbf{Orientation.} Rappeler la condition que doit vérifier le déterminant $\det(\vec{i},\vec{j},\vec{k})$ pour que la base $(\vec{i},\vec{j},\vec{k})$ soit directe. Calculer ce déterminant et conclure.
\item \textbf{Vecteur normal au toit.}
      \begin{enumerate}
      \item Déterminer les coordonnées des vecteurs $\overrightarrow{AB}$ et $\overrightarrow{AC}$.
      \item Calculer $\overrightarrow{AB}\wedge\overrightarrow{AC}$. Justifier que ce vecteur est normal au plan $(ABC)$.
      \item En déduire un vecteur normal $\vec{n}$ plus simple, puis une équation cartésienne du plan $(ABC)$.
      \end{enumerate}
\item \textbf{Quantité de tôle.} En utilisant le résultat de la question 2, calculer l'aire du triangle $ABC$. Donner la valeur exacte, puis une valeur approchée à $0{,}1\ \mathrm{m}^2$ près.
\item \textbf{Entretoise.} Calculer la distance du point $P(2\,;2\,;0)$ au plan $(ABC)$. Donner la valeur exacte, puis une valeur approchée au centimètre près.
\item \textbf{Gouttière.} Calculer $\overrightarrow{AP}\wedge\overrightarrow{AB}$, puis la distance du point $P$ à la droite $(AB)$.
\item \textbf{Volume de stockage.} Calculer le produit mixte $\left[\overrightarrow{OA},\overrightarrow{OB},\overrightarrow{OC}\right]$, puis le volume du tétraèdre $OABC$. Combien de sacs de maïs de $0{,}1\ \mathrm{m}^3$ peut-on théoriquement y ranger ?
\item \textbf{Restitution.} Comparer vos résultats avec ceux d'un autre groupe. En cas de désaccord, repérer ensemble l'étape fautive.
\end{enumerate}

\courssub{Démarche et corrigé commenté}

\begin{exoresolu}[Corrigé complet de l'activité]
Vérifier l'orientation de la base, déterminer le plan du toit, l'aire de tôle, les distances d'entretoise et de gouttière, et le volume de stockage du hangar de l'ONCPA.
\tcblower
\textbf{1. Orientation de la base.}
La base $(\vec{i},\vec{j},\vec{k})$ est directe si et seulement si $\det(\vec{i},\vec{j},\vec{k})>0$. Or, dans la base canonique, les vecteurs ont pour coordonnées $\vec{i}\begin{pmatrix}1\\0\\0\end{pmatrix}$, $\vec{j}\begin{pmatrix}0\\1\\0\end{pmatrix}$, $\vec{k}\begin{pmatrix}0\\0\\1\end{pmatrix}$.
\[ \det(\vec{i},\vec{j},\vec{k})=\begin{vmatrix} 1&0&0\\0&1&0\\0&0&1 \end{vmatrix}=1>0. \]
La base est donc orthonormée \textbf{directe} : toutes les formules du produit vectoriel s'appliquent sans changement de signe.

\textbf{2. Vecteur normal et équation du plan du toit.}

a) $\overrightarrow{AB}\begin{pmatrix}0-6\\6-0\\4-4\end{pmatrix}=\begin{pmatrix}-6\\6\\0\end{pmatrix}$ et $\overrightarrow{AC}\begin{pmatrix}0-6\\0-0\\6-4\end{pmatrix}=\begin{pmatrix}-6\\0\\2\end{pmatrix}$.

b) On applique la formule du produit vectoriel dans une base orthonormée directe :
\[ \overrightarrow{AB}\wedge\overrightarrow{AC}=\begin{pmatrix} 6\times 2-0\times 0 \\ 0\times(-6)-(-6)\times 2 \\ (-6)\times 0-6\times(-6) \end{pmatrix}=\begin{pmatrix}12\\12\\36\end{pmatrix}. \]
Par définition, $\overrightarrow{AB}\wedge\overrightarrow{AC}$ est orthogonal à $\overrightarrow{AB}$ et à $\overrightarrow{AC}$, donc à deux vecteurs directeurs non colinéaires du plan $(ABC)$ : c'est un \textbf{vecteur normal} au toit.

c) On simplifie par $12$ : $\vec{n}\begin{pmatrix}1\\1\\3\end{pmatrix}$ est aussi normal au plan. Le plan $(ABC)$ admet une équation de la forme $x+y+3z+d=0$. Comme $A(6\,;0\,;4)$ appartient au plan : $6+0+12+d=0$, d'où $d=-18$.
\[ (ABC)\ :\ x+y+3z-18=0. \]
\emph{Vérification croisée} : $B$ donne $0+6+12-18=0$ et $C$ donne $0+0+18-18=0$. Les trois sommets vérifient l'équation.

\textbf{3. Aire de la toiture.}
\[ \mathscr{A}=\frac{1}{2}\left\|\overrightarrow{AB}\wedge\overrightarrow{AC}\right\|=\frac{1}{2}\sqrt{12^2+12^2+36^2}=\frac{1}{2}\sqrt{1584}=\frac{1}{2}\times 12\sqrt{11}=6\sqrt{11}. \]
Donc $\mathscr{A}=6\sqrt{11}\ \mathrm{m}^2\approx 19{,}9\ \mathrm{m}^2$. L'ingénieur commandera environ $20\ \mathrm{m}^2$ de tôle (en prévoyant une marge pour les découpes).

\textbf{4. Distance de $P$ au plan du toit (entretoise).}
\[ d\big(P,(ABC)\big)=\frac{|2+2+3\times 0-18|}{\sqrt{1^2+1^2+3^2}}=\frac{|-14|}{\sqrt{11}}=\frac{14}{\sqrt{11}}=\frac{14\sqrt{11}}{11}\approx 4{,}22\ \mathrm{m}. \]
L'entretoise reliant le pied du poteau au toit mesure environ $4{,}22\ \mathrm{m}$ (soit $422\ \mathrm{cm}$).

\textbf{5. Distance de $P$ à la gouttière $(AB)$.}
On a $\overrightarrow{AP}\begin{pmatrix}2-6\\2-0\\0-4\end{pmatrix}=\begin{pmatrix}-4\\2\\-4\end{pmatrix}$. Alors :
\[ \overrightarrow{AP}\wedge\overrightarrow{AB}=\begin{pmatrix} 2\times 0-(-4)\times 6 \\ (-4)\times(-6)-(-4)\times 0 \\ (-4)\times 6-2\times(-6) \end{pmatrix}=\begin{pmatrix}24\\24\\-12\end{pmatrix}. \]
Sa norme : $\sqrt{24^2+24^2+(-12)^2}=\sqrt{576+576+144}=\sqrt{1296}=36$.
Par ailleurs $\left\|\overrightarrow{AB}\right\|=\sqrt{(-6)^2+6^2+0^2}=\sqrt{72}=6\sqrt{2}$.
Donc :
\[ d\big(P,(AB)\big)=\frac{36}{6\sqrt{2}}=\frac{6}{\sqrt{2}}=3\sqrt{2}\approx 4{,}24\ \mathrm{m}. \]
La gouttière passe à environ $4{,}24\ \mathrm{m}$ du pied du poteau.

\textbf{6. Volume du coin de stockage.}
On calcule le produit mixte (déterminant des trois vecteurs) :
\[ \left[\overrightarrow{OA},\overrightarrow{OB},\overrightarrow{OC}\right]=\det(\overrightarrow{OA},\overrightarrow{OB},\overrightarrow{OC})=\begin{vmatrix} 6 & 0 & 4 \\ 0 & 6 & 4 \\ 0 & 0 & 6 \end{vmatrix}. \]
La matrice est triangulaire supérieure, son déterminant est le produit de la diagonale : $6 \times 6 \times 6 = 216$.
Le volume du tétraèdre est :
\[ \mathscr{V}=\frac{1}{6}\times|216|=36\ \mathrm{m}^3. \]
À raison de $0{,}1\ \mathrm{m}^3$ par sac, le coin de stockage peut contenir théoriquement $36\div 0{,}1=360$ sacs de maïs.

\textbf{Conclusion pour l'ingénieur :} toiture de $6\sqrt{11}\approx 19{,}9\ \mathrm{m}^2$, entretoise de $\dfrac{14\sqrt{11}}{11}\approx 4{,}22\ \mathrm{m}$, gouttière à $3\sqrt{2}\approx 4{,}24\ \mathrm{m}$ du poteau, capacité du coin de stockage : $36\ \mathrm{m}^3$ (360 sacs).
\end{exoresolu}

\begin{attention}[Pièges rencontrés pendant l'activité]
\begin{itemize}
\item Ne pas oublier de vérifier que la base est \textbf{directe} avant d'appliquer la formule des coordonnées du produit vectoriel : dans une base indirecte, le résultat serait changé en son opposé.
\item L'aire du triangle exige le facteur $\dfrac{1}{2}$ : $\|\vec{u}\wedge\vec{v}\|$ est l'aire du \emph{parallélogramme} construit sur $\vec{u}$ et $\vec{v}$.
\item Le volume du tétraèdre exige le facteur $\dfrac{1}{6}$, et non $\dfrac{1}{3}$ (qui convient au prisme) ni $1$ (parallélépipède rectangle).
\item Dans la distance point-plan, ne pas oublier la valeur absolue au numérateur et la norme $\sqrt{a^2+b^2+c^2}$ au dénominateur.
\item Pour le produit mixte, l'ordre des vecteurs change le signe mais la valeur absolue pour le volume rend le résultat invariant par permutation circulaire.
\end{itemize}
\end{attention}

\courssub{Bilan de l'activité}

Cette activité a permis de mettre en évidence que le \cle{produit vectoriel} est un outil unique qui répond à des questions d'apparence très différentes :

\begin{itemize}
\item un \textbf{même calcul}, $\overrightarrow{AB}\wedge\overrightarrow{AC}$, a fourni à la fois le vecteur normal au plan du toit (question 2) et l'aire de la toiture (question 3) : la direction et la norme du produit vectoriel portent chacune une information géométrique ;
\item les formules de \textbf{distances} (point-plan, point-droite) ne sont que des réécritures du produit scalaire et du produit vectoriel : il n'y a rien à apprendre de nouveau, il suffit de savoir les mobiliser ;
\item le \textbf{produit mixte} ramène le volume d'un tétraèdre à un simple déterminant $3\times 3$ ;
\item enfin, la condition « base orthonormée \emph{directe} » n'est pas un détail théorique : c'est elle qui garantit la validité de tous les calculs menés ensuite, comme l'ingénieur vérifie le calibrage de ses instruments avant toute mesure.
\end{itemize}

Tu es désormais capable de traiter, de bout en bout, un problème d'espace issu d'une situation professionnelle réelle : c'est exactement le type de mobilisation de compétences attendu au Baccalauréat.

\begin{savaistu}[Le produit vectoriel dans l'industrie et la 3D]
Au-delà du hangar de l'ONCPA, le produit vectoriel est au cœur de la \emph{infographie 3D} (calcul des normales pour l'éclairage, le \emph{shading}, la détection de collisions), de la \emph{robotique} (calcul des couples moteurs, cinématique des bras articulés), et de la \emph{navigation} (produit vectoriel pour corriger la dérive d'un drone ou d'un bateau). Au Cameroun, les logiciels de CAO (Conception Assistée par Ordinateur) utilisés par les bureaux d'études de Douala ou Yaoundé pour dimensionner les barrages (Lom Pangar, Memve'ele), les tours de télécommunication (MTN, Orange) ou les routes, reposent tous sur ces formules vectorielles.
\end{savaistu}

\newpage

\coursec{Méthodes et exercices résolus}

\coursec{Orientation de l'espace}

\begin{definition}[Base orthonormée directe et indirecte]
Soit $(\vec{i},\vec{j},\vec{k})$ une base orthonormée de l'espace. On dit qu'elle est \cle{directe} si le triplet $(\vec{i},\vec{j},\vec{k})$ satisfait la \textbf{règle de la main droite} (ou règle du bonhomme d'Ampère) :
\begin{itemize}
\item Pouce $\rightarrow \vec{i}$ (axe $x$),
\item Index $\rightarrow \vec{j}$ (axe $y$),
\item Majeur $\rightarrow \vec{k}$ (axe $z$).
\end{itemize}
Si l'on doit utiliser la main gauche pour superposer les doigts aux vecteurs, la base est \cle{indirecte}.
\end{definition}

\begin{popfigure}
\begin{tikzpicture}[scale=1.2, >=Stealth]
% Main droite
\draw[popblue, thick] (0,0) -- (1,0) node[midway, below] {$\vec{i}$};
\draw[poporange, thick] (0,0) -- (0,1) node[midway, left] {$\vec{j}$};
\draw[popgreen, thick] (0,0) -- (-0.3,-0.3) node[midway, below left] {$\vec{k}$};
% Schéma main (simplifié)
\draw[thick, popmuted] (2,0) .. controls (2.5,0.5) and (3,0.5) .. (3,1) node[right] {Pouce ($\vec{i}$)};
\draw[thick, popmuted] (2,0) .. controls (2.5,-0.2) and (3,-0.2) .. (3,-0.5) node[right] {Index ($\vec{j}$)};
\draw[thick, popmuted] (2,0) .. controls (2.2,0.2) and (2.5,0.5) .. (2.5,1) node[above] {Majeur ($\vec{k}$)};
\node at (1.5, -1.2) {Base \textbf{directe} (main droite)};
% Base indirecte
\begin{scope}[xshift=6cm]
\draw[popblue, thick] (0,0) -- (1,0) node[midway, below] {$\vec{i}$};
\draw[poporange, thick] (0,0) -- (0,1) node[midway, left] {$\vec{j}$};
\draw[popgreen, thick] (0,0) -- (0.3,-0.3) node[midway, below right] {$\vec{k}$};
\draw[thick, popmuted] (2,0) .. controls (2.5,0.5) and (3,0.5) .. (3,1) node[right] {Pouce ($\vec{i}$)};
\draw[thick, popmuted] (2,0) .. controls (2.5,-0.2) and (3,-0.2) .. (3,-0.5) node[right] {Index ($\vec{j}$)};
\draw[thick, popmuted, dashed] (2,0) .. controls (2.2,0.2) and (2.5,0.5) .. (2.5,1) node[above] {Majeur ($\vec{k}$) \textcolor{red}{(impossible)}};
\node at (1.5, -1.2) {Base \textbf{indirecte} (main gauche nécessaire)};
\end{scope}
\end{tikzpicture}
\legende{Règle de la main droite : $(\vec{i},\vec{j},\vec{k})$ direct $\iff$ pouce/index/majeur alignés sur $\vec{i},\vec{j},\vec{k}$.}
\end{popfigure}

\begin{savaistu}[Repères au Cameroun : antennes et cartographie]
Les antennes relais \textbf{MTN} et \textbf{Orange} sont orientées dans l'espace par des azimuths (angle horizontal) et des tilts (inclinaison verticale). Les ingénieurs télécoms utilisent des repères directs pour calculer les zones de couverture (faisceaux). En cartographie (SIG), le système de projection UTM (utilisé pour le cadastre à Yaoundé ou Douala) repose sur un repère orthonormé direct local pour garantir que les aires mesurées sur la carte correspondent aux aires réelles au sol sans inversion de signe.
\end{savaistu}

\begin{experience}[Découvrir l'orientation par le produit vectoriel]
Dans un repère orthonormé direct $(O;\vec{i},\vec{j},\vec{k})$, calculer $\vec{i}\wedge\vec{j}$, $\vec{j}\wedge\vec{k}$, $\vec{k}\wedge\vec{i}$.
\begin{enumerate}
\item $\vec{i}(1;0;0)$, $\vec{j}(0;1;0)$ $\Rightarrow$ $\vec{i}\wedge\vec{j} = (0;0;1) = \vec{k}$.
\item $\vec{j}\wedge\vec{k} = \vec{i}$.
\item $\vec{k}\wedge\vec{i} = \vec{j}$.
\end{enumerate}
On retrouve le cycle direct $\vec{i}\to\vec{j}\to\vec{k}\to\vec{i}$. Si la base était indirecte, $\vec{i}\wedge\vec{j} = -\vec{k}$. \textbf{Conclusion :} La formule analytique du produit vectoriel (règle du gamma) \textbf{ne est valable que dans une base orthonormée directe}.
\end{experience}

\coursec{Définition et interprétation géométrique du produit vectoriel}

\begin{definition}[Produit vectoriel de deux vecteurs]
Soit $\vec{u}$ et $\vec{v}$ deux vecteurs de l'espace. On appelle \cle{produit vectoriel} de $\vec{u}$ par $\vec{v}$, noté $\vec{u}\wedge\vec{v}$, le vecteur $\vec{w}$ tel que :
\begin{enumerate}
\item \textbf{Direction :} $\vec{w}$ est orthogonal à $\vec{u}$ et à $\vec{v}$ (donc orthogonal au plan $(\vec{u},\vec{v})$).
\item \textbf{Sens :} Le triplet $(\vec{u},\vec{v},\vec{w})$ est \textbf{direct} (règle de la main droite : $\vec{u}$ pouce, $\vec{v}$ index $\Rightarrow$ $\vec{w}$ majeur).
\item \textbf{Norme :} $\|\vec{w}\| = \|\vec{u}\|\,\|\vec{v}\|\,|\sin(\vec{u},\vec{v})|$.
\end{enumerate}
Si $\vec{u}=\vec{0}$ ou $\vec{v}=\vec{0}$ ou $\vec{u}\parallel\vec{v}$, on pose $\vec{u}\wedge\vec{v}=\vec{0}$.
\end{definition}

\begin{popfigure}
\begin{tikzpicture}[scale=1.5, >=Stealth]
% Vecteurs u et v dans le plan horizontal
\draw[popblue, thick, ->] (0,0,0) -- (2,0,0) node[below] {$\vec{u}$};
\draw[poporange, thick, ->] (0,0,0) -- (1,1.5,0) node[left] {$\vec{v}$};
% Plan (parallélogramme)
\draw[popblue!30, fill=popblue!10, opacity=0.6] (0,0,0) -- (2,0,0) -- (3,1.5,0) -- (1,1.5,0) -- cycle;
% Produit vectoriel (vertical)
\draw[popgreen, very thick, ->] (0,0,0) -- (0,0,1.8) node[above] {$\vec{u}\wedge\vec{v}$};
% Angle
\draw[poppurple, thick] (0.3,0,0) arc (0:56:0.3) node[midway, right=2pt] {$\theta$};
% Main droite indicative
\draw[popmuted, thin] (2.5, -0.5) node[align=center] {Main droite :\\ pouce $\vec{u}$, index $\vec{v}$\\ $\Rightarrow$ majeur $\vec{u}\wedge\vec{v}$};
\end{tikzpicture}
\legende{Interprétation géométrique : $\vec{u}\wedge\vec{v}$ orthogonal au plan, sens donné par la main droite, norme = aire du parallélogramme.}
\end{popfigure}

\begin{propriete}[Propriétés géométriques fondamentales]
\begin{itemize}
\item \textbf{Aire :} $\|\vec{u}\wedge\vec{v}\|$ est l'aire du parallélogramme construit sur $\vec{u}$ et $\vec{v}$.
\item \textbf{Orthogonalité :} $(\vec{u}\wedge\vec{v})\cdot\vec{u}=0$ et $(\vec{u}\wedge\vec{v})\cdot\vec{v}=0$.
\item \textbf{Colinéarité :} $\vec{u}\wedge\vec{v}=\vec{0} \iff \vec{u}$ et $\vec{v}$ sont colinéaires (angle $0$ ou $\pi$, $\sin=0$).
\item \textbf{Antisymétrie :} $\vec{v}\wedge\vec{u} = -(\vec{u}\wedge\vec{v})$ (changer l'ordre inverse le sens, triplet devient indirect).
\end{itemize}
\end{propriete}

\coursec{Expression analytique dans une base orthonormée directe}

\begin{propriete}[Formule analytique (Règle du gamma)]
Dans un repère orthonormé \textbf{direct} $(O;\vec{i},\vec{j},\vec{k})$, pour $\vec{u}(x;y;z)$ et $\vec{v}(x';y';z')$ :
\[
\vec{u}\wedge\vec{v} =
\begin{pmatrix}
y z' - z y' \\
z x' - x z' \\
x y' - y x'
\end{pmatrix}
=
\begin{vmatrix}
\vec{i} & \vec{j} & \vec{k} \\
x & y & z \\
x' & y' & z'
\end{vmatrix}.
\]
\end{propriete}

\begin{methode}[Calculer $\vec{u}\wedge\vec{v}$ à partir des coordonnées]
\begin{enumerate}
\item \textbf{Vérifier que le repère est orthonormé direct.} (Condition d'application obligatoire, à mentionner dans la rédaction).
\item Écrire les coordonnées en colonnes et appliquer la règle du \emph{gamma} (ou calculer le déterminant formel).
\item \textbf{Astuce de mémorisation} : pour la $i$-ème coordonnée, masquer la $i$-ème ligne et faire le produit en croix sur les deux lignes restantes en respectant l'ordre cyclique $(x,y,z) \to (y,z,x) \to (z,x,y)$.
  \begin{itemize}
  \item 1ère coord. : masque $x$ $\to$ $y z' - z y'$
  \item 2ème coord. : masque $y$ $\to$ $z x' - x z'$ \textbf{(attention : $z x'$ en premier !)}
  \item 3ème coord. : masque $z$ $\to$ $x y' - y x'$
  \end{itemize}
\item \textbf{Contrôle systématique (réflexe Bac) :} Vérifier $(\vec{u}\wedge\vec{v})\cdot\vec{u}=0$ et $(\vec{u}\wedge\vec{v})\cdot\vec{v}=0$.
\end{enumerate}
\end{methode}

\begin{exoresolu}[Produit vectoriel et orthogonalité]
L'espace est muni d'un repère orthonormé direct $(O;\vec{i},\vec{j},\vec{k})$. On donne $\vec{u}(2;-1;3)$ et $\vec{v}(1;4;-2)$.
\begin{enumerate}
\item Calculer $\vec{u}\wedge\vec{v}$.
\item Vérifier que $\vec{u}\wedge\vec{v}$ est orthogonal à $\vec{u}$ et à $\vec{v}$.
\item Calculer $\|\vec{u}\wedge\vec{v}\|$ et en déduire l'aire du parallélogramme construit sur $\vec{u}$ et $\vec{v}$.
\end{enumerate}
\tcblower
\textbf{1.} Le repère est orthonormé direct (donnée de l'énoncé), on applique la formule :
\[
\vec{u}\wedge\vec{v}
=
\begin{pmatrix}
(-1)(-2) - (3)(4) \\
(3)(1) - (2)(-2) \\
(2)(4) - (-1)(1)
\end{pmatrix}
=
\begin{pmatrix}
2 - 12 \\
3 + 4 \\
8 + 1
\end{pmatrix}
=
\begin{pmatrix}
-10 \\ 7 \\ 9
\end{pmatrix}.
\]

\textbf{2.} Contrôle par produit scalaire :
\begin{align*}
(\vec{u}\wedge\vec{v})\cdot\vec{u} &= (-10)(2) + (7)(-1) + (9)(3) = -20 - 7 + 27 = 0,\\
(\vec{u}\wedge\vec{v})\cdot\vec{v} &= (-10)(1) + (7)(4) + (9)(-2) = -10 + 28 - 18 = 0.
\end{align*}
Le vecteur $\vec{u}\wedge\vec{v}$ est bien orthogonal à $\vec{u}$ et à $\vec{v}$.

\textbf{3.} $\|\vec{u}\wedge\vec{v}\| = \sqrt{(-10)^2 + 7^2 + 9^2} = \sqrt{100 + 49 + 81} = \sqrt{230}$.
L'aire du parallélogramme est $\mathcal{A} = \|\vec{u}\wedge\vec{v}\| = \sqrt{230}$ unités d'aire.

\textbf{Conclusion :} $\vec{u}\wedge\vec{v}(-10;7;9)$, orthogonal à $\vec{u}$ et $\vec{v}$, aire $= \sqrt{230}$.
\end{exoresolu}

\begin{aretenir}[Propriétés algébriques (à connaître pour le Bac)]
\begin{itemize}
\item \textbf{Antisymétrie :} $\vec{v}\wedge\vec{u} = -(\vec{u}\wedge\vec{v})$.
\item \textbf{Bilinéarité :}
  \begin{align*}
  (\vec{u}_1+\vec{u}_2)\wedge\vec{v} &= \vec{u}_1\wedge\vec{v} + \vec{u}_2\wedge\vec{v}, \\
  \vec{u}\wedge(\vec{v}_1+\vec{v}_2) &= \vec{u}\wedge\vec{v}_1 + \vec{u}\wedge\vec{v}_2, \\
  (\lambda\vec{u})\wedge\vec{v} &= \vec{u}\wedge(\lambda\vec{v}) = \lambda(\vec{u}\wedge\vec{v}).
  \end{align*}
\item \textbf{Non-associativité :} $(\vec{u}\wedge\vec{v})\wedge\vec{w} \neq \vec{u}\wedge(\vec{v}\wedge\vec{w})$ en général.
\item \textbf{Identité de Jacobi (culture) :} $\vec{u}\wedge(\vec{v}\wedge\vec{w}) + \vec{v}\wedge(\vec{w}\wedge\vec{u}) + \vec{w}\wedge(\vec{u}\wedge\vec{v}) = \vec{0}$.
\item \textbf{Colinéarité :} $\vec{u}\wedge\vec{v}=\vec{0} \iff \vec{u}\parallel\vec{v}$.
\end{itemize}
\end{aretenir}

\coursec{Applications : Vecteur normal, distances, aires, volume}

\courssub{Déterminer un vecteur normal à un plan}

\begin{methode}[Trouver un vecteur normal à un plan $(ABC)$]
\begin{enumerate}
\item Calculer les coordonnées de deux vecteurs \textbf{non colinéaires} du plan, par exemple $\overrightarrow{AB}$ et $\overrightarrow{AC}$.
\item \textbf{Justifier la non-colinéarité} (montrer qu'il n'existe pas $k$ tel que $\overrightarrow{AC}=k\overrightarrow{AB}$, ou que le produit vectoriel est non nul).
\item Calculer $\vec{n} = \overrightarrow{AB}\wedge\overrightarrow{AC}$. Par définition, $\vec{n}$ est orthogonal à $\overrightarrow{AB}$ et $\overrightarrow{AC}$, donc normal au plan $(ABC)$.
\item Si une équation cartésienne est demandée : le plan a pour équation $ax+by+cz+d=0$ avec $\vec{n}(a;b;c)$. On trouve $d$ en utilisant les coordonnées d'un point du plan (ex: $A$).
\end{enumerate}
\end{methode}

\begin{attention}[Condition d'application cruciale]
Le produit vectoriel ne donne un vecteur normal \textbf{que si} les deux vecteurs de départ sont \cle{non colinéaires}. Si $\overrightarrow{AB}\wedge\overrightarrow{AC}=\vec{0}$, les points $A,B,C$ sont alignés : ils ne définissent pas un plan unique.
\end{attention}

\begin{exoresolu}[Vecteur normal et équation de plan]
Dans un repère orthonormé direct, on donne $A(1;0;2)$, $B(3;1;-1)$ et $C(0;2;1)$.
\begin{enumerate}
\item Démontrer que $A$, $B$, $C$ définissent un plan.
\item Déterminer un vecteur normal au plan $(ABC)$.
\item En déduire une équation cartésienne du plan $(ABC)$.
\end{enumerate}
\tcblower
\textbf{1.} $\overrightarrow{AB}(2;1;-3)$ et $\overrightarrow{AC}(-1;2;-1)$.
Test de colinéarité : si $\overrightarrow{AC}=k\overrightarrow{AB}$, alors $-1=2k \Rightarrow k=-\frac12$. Mais $2 = k\times 1 = -\frac12$ : contradiction. Les vecteurs ne sont pas colinéaires, les points ne sont pas alignés, ils définissent un plan.

\textbf{2.} $\vec{n} = \overrightarrow{AB}\wedge\overrightarrow{AC} =
\begin{pmatrix}
(1)(-1) - (-3)(2) \\
(-3)(-1) - (2)(-1) \\
(2)(2) - (1)(-1)
\end{pmatrix}
=
\begin{pmatrix}
5 \\ 5 \\ 5
\end{pmatrix}
$.
On simplifie : $\vec{n}'(1;1;1)$ est aussi un vecteur normal (tout multiple non nul convient).

\textbf{3.} Équation : $x+y+z+d=0$. $A(1;0;2)\in(ABC) \Rightarrow 1+0+2+d=0 \Rightarrow d=-3$.
Équation cartésienne : $x+y+z-3=0$.

\textbf{Conclusion :} Vecteur normal $\vec{n}(1;1;1)$, plan : $x+y+z-3=0$.
\end{exoresolu}

\courssub{Calculer la distance d'un point à une droite}

\begin{methode}[Distance d'un point $M$ à une droite $(AB)$]
\begin{enumerate}
\item Calculer $\overrightarrow{AB}$ (directeur de la droite) et $\overrightarrow{AM}$.
\item Calculer $\overrightarrow{AB}\wedge\overrightarrow{AM}$ et sa norme.
\item Appliquer la formule :
\[
d\big(M,(AB)\big) = \frac{\|\overrightarrow{AB}\wedge\overrightarrow{AM}\|}{\|\overrightarrow{AB}\|}.
\]
\item \textbf{Justification (à savoir rédiger) :} $\|\overrightarrow{AB}\wedge\overrightarrow{AM}\|$ est l'aire du parallélogramme construit sur $\overrightarrow{AB}$ et $\overrightarrow{AM}$. Cette aire vaut aussi $\text{base}\times\text{hauteur} = AB \times d(M,(AB))$.
\end{enumerate}
\end{methode}

\begin{aretenir}[Test d'alignement]
$M \in (AB) \iff \overrightarrow{AB}\wedge\overrightarrow{AM} = \vec{0} \iff \overrightarrow{AB}$ et $\overrightarrow{AM}$ sont colinéaires.
\end{aretenir}

\begin{exoresolu}[Distance d'un point à une droite]
Dans un repère orthonormé direct $(O;\vec{i},\vec{j},\vec{k})$ d'unité $\SI{1}{cm}$, $A(2;1;0)$, $B(0;3;1)$, $M(1;-1;4)$.
\begin{enumerate}
\item Montrer que $M \notin (AB)$.
\item Calculer $d(M,(AB))$ (valeur exacte et approchée $10^{-2}$).
\end{enumerate}
\tcblower
\textbf{1.} $\overrightarrow{AB}(-2;2;1)$, $\overrightarrow{AM}(-1;-2;4)$.
Rapports : $\frac{-1}{-2}=0,5$ ; $\frac{-2}{2}=-1$. Non égaux $\Rightarrow$ non colinéaires $\Rightarrow M\notin(AB)$.

\textbf{2.} $\overrightarrow{AB}\wedge\overrightarrow{AM} =
\begin{pmatrix}
(2)(4)-(1)(-2) \\
(1)(-1)-(-2)(4) \\
(-2)(-2)-(2)(-1)
\end{pmatrix}
=
\begin{pmatrix}
10 \\ 7 \\ 6
\end{pmatrix}
$.
$\|\overrightarrow{AB}\wedge\overrightarrow{AM}\| = \sqrt{100+49+36} = \sqrt{185}$.
$\|\overrightarrow{AB}\| = \sqrt{4+4+1} = 3$.
$d(M,(AB)) = \frac{\sqrt{185}}{3} \approx \frac{13,60}{3} \approx \SI{4,53}{cm}$.

\textbf{Conclusion :} Distance exacte $\frac{\sqrt{185}}{3}$ cm, approchée $\SI{4,53}{cm}$.
\end{exoresolu}

\courssub{Calculer la distance d'un point à un plan}

\begin{methode}[Distance d'un point $M$ à un plan $(ABC)$]
\textbf{Variante 1 — Par projection sur un vecteur normal.}
\begin{enumerate}
\item Déterminer $\vec{n} = \overrightarrow{AB}\wedge\overrightarrow{AC}$ (normal au plan).
\item Calculer $\overrightarrow{AM}$ et le produit scalaire $\vec{n}\cdot\overrightarrow{AM}$.
\item Appliquer : $d(M,(ABC)) = \frac{|\vec{n}\cdot\overrightarrow{AM}|}{\|\vec{n}\|}$.
\end{enumerate}
\textbf{Variante 2 — Par équation cartésienne.}
Si le plan a pour équation $ax+by+cz+d=0$ et $M(x_M;y_M;z_M)$ :
\[
d(M,\mathcal{P}) = \frac{|ax_M+by_M+cz_M+d|}{\sqrt{a^2+b^2+c^2}}.
\]
\end{methode}

\begin{attention}[Piège classique : la valeur absolue]
Une distance est \textbf{toujours positive ou nulle}. Ne jamais oublier les barres de valeur absolue au numérateur. Un résultat négatif est une erreur grave.
\end{attention}

\begin{exoresolu}[Distance d'un point à un plan]
On reprend le plan $(ABC)$ d'équation $x+y+z-3=0$ ($A(1;0;2)$, $B(3;1;-1)$, $C(0;2;1)$) et $M(4;-2;5)$.
\begin{enumerate}
\item Calculer $d(M,(ABC))$ par projection.
\item Retrouver le résultat par l'équation cartésienne.
\end{enumerate}
\tcblower
\textbf{1.} $\vec{n}(1;1;1)$, $\overrightarrow{AM}(3;-2;3)$.
$\vec{n}\cdot\overrightarrow{AM} = 3-2+3 = 4$. $\|\vec{n}\| = \sqrt{3}$.
$d = \frac{|4|}{\sqrt{3}} = \frac{4\sqrt{3}}{3}$.

\textbf{2.} Formule cartésienne : $a=b=c=1, d=-3$.
$d = \frac{|4 + (-2) + 5 - 3|}{\sqrt{3}} = \frac{|4|}{\sqrt{3}} = \frac{4\sqrt{3}}{3}$.
Les deux méthodes concordent.

\textbf{Conclusion :} $d(M,(ABC)) = \frac{4\sqrt{3}}{3} \approx 2,31$ unités de longueur.
\end{exoresolu}

\courssub{Calculer l'aire d'un triangle}

\begin{methode}[Aire d'un triangle $ABC$ de l'espace]
\begin{enumerate}
\item Calculer deux vecteurs côtés issus du même sommet : $\overrightarrow{AB}$ et $\overrightarrow{AC}$.
\item Calculer $\overrightarrow{AB}\wedge\overrightarrow{AC}$ et sa norme (aire du parallélogramme).
\item Diviser par $2$ :
\[
\mathcal{A}_{ABC} = \frac{1}{2}\,\|\overrightarrow{AB}\wedge\overrightarrow{AC}\|.
\]
\end{enumerate}
\end{methode}

\begin{attention}[Erreur n°1 au Bac : oublier le $\frac12$]
L'aire du \textbf{triangle} est la \emph{moitié} de la norme du produit vectoriel. Oublier le facteur $\frac12$ donne l'aire du parallélogramme : c'est l'erreur la plus fréquente sur ce chapitre.
\end{attention}

\begin{exoresolu}[Aire d'une voile triangulaire]
Un artisan de Douala confectionne une voile d'ombrage. Repère orthonormé direct (unité : mètre), $A(0;0;3)$, $B(4;0;2)$, $C(0;3;1)$.
\begin{enumerate}
\item Calculer l'aire de la voile $ABC$.
\item Coût du tissu : $\num{2500}$ FCFA/$\text{m}^2$. Estimer le coût ($\sqrt{217}\approx 14,73$).
\end{enumerate}
\tcblower
\textbf{1.} $\overrightarrow{AB}(4;0;-1)$, $\overrightarrow{AC}(0;3;-2)$.
$\overrightarrow{AB}\wedge\overrightarrow{AC} =
\begin{pmatrix}
(0)(-2)-(-1)(3) \\
(-1)(0)-(4)(-2) \\
(4)(3)-(0)(0)
\end{pmatrix}
=
\begin{pmatrix}
3 \\ 8 \\ 12
\end{pmatrix}
$.
$\|\overrightarrow{AB}\wedge\overrightarrow{AC}\| = \sqrt{9+64+144} = \sqrt{217}$.
$\mathcal{A}_{ABC} = \frac{\sqrt{217}}{2} \text{ m}^2 \approx \frac{14,73}{2} \approx \SI{7,37}{m^2}$.

\textbf{2.} Coût $\approx 7,37 \times 2500 = \num{18425}$ FCFA $\approx \num{18500}$ FCFA (avec marge).

\textbf{Conclusion :} Aire $\frac{\sqrt{217}}{2}\ \text{m}^2 \approx \SI{7,37}{m^2}$, coût $\approx \num{18400}$ FCFA.
\end{exoresolu}

\courssub{Produit mixte et volume d'un tétraèdre}

\begin{definition}[Produit mixte]
Soit $\vec{u},\vec{v},\vec{w}$ trois vecteurs. On appelle \cle{produit mixte} de $\vec{u},\vec{v},\vec{w}$ le scalaire :
\[
[\vec{u},\vec{v},\vec{w}] = (\vec{u}\wedge\vec{v})\cdot\vec{w}.
\]
Géométriquement, $|[\vec{u},\vec{v},\vec{w}]|$ est le volume du parallélépipède construit sur $\vec{u},\vec{v},\vec{w}$.
\end{definition}

\begin{propriete}[Propriétés du produit mixte]
\begin{itemize}
\item Invariance par permutation circulaire : $[\vec{u},\vec{v},\vec{w}] = [\vec{v},\vec{w},\vec{u}] = [\vec{w},\vec{u},\vec{v}]$.
\item Changement de signe par transposition : $[\vec{u},\vec{v},\vec{w}] = -[\vec{v},\vec{u},\vec{w}]$.
\item $[\vec{u},\vec{v},\vec{w}] = 0 \iff \vec{u},\vec{v},\vec{w}$ sont coplanaires (test de coplanarité).
\end{itemize}
\end{propriete}

\begin{methode}[Volume d'un tétraèdre $ABCD$]
\begin{enumerate}
\item Choisir un sommet $A$, calculer $\overrightarrow{AB},\overrightarrow{AC},\overrightarrow{AD}$.
\item Calculer le produit mixte : $V_{\text{par}} = \big|(\overrightarrow{AB}\wedge\overrightarrow{AC})\cdot\overrightarrow{AD}\big|$ (volume du parallélépipède).
\item Volume du tétraèdre : $\mathcal{V}_{ABCD} = \frac{1}{6} V_{\text{par}} = \frac{1}{6}\big|(\overrightarrow{AB}\wedge\overrightarrow{AC})\cdot\overrightarrow{AD}\big|$.
\end{enumerate}
\end{methode}

\begin{aretenir}[Formules de volumes]
\begin{itemize}
\item Parallélépipède : $V = |$produit mixte$|$.
\item Tétraèdre : $V = \frac{1}{6}|$produit mixte$|$.
\item Prisme (base triangle) : $V = \frac{1}{2}|$produit mixte$|$.
\item Tétraèdre (base $\times$ hauteur) : $V = \frac{1}{3}\times\mathcal{A}_{\text{base}}\times h$.
\end{itemize}
\end{aretenir}

\begin{exoresolu}[Volume d'un tétraèdre et coplanarité]
Repère orthonormé direct. $A(1;0;0)$, $B(0;2;0)$, $C(0;0;3)$, $D(2;1;4)$.
\begin{enumerate}
\item Volume du tétraèdre $ABCD$.
\item $E(3;2;-3)$ est-il coplanaire avec $A,B,C$ ?
\end{enumerate}
\tcblower
\textbf{1.} $\overrightarrow{AB}(-1;2;0)$, $\overrightarrow{AC}(-1;0;3)$, $\overrightarrow{AD}(1;1;4)$.
$\overrightarrow{AB}\wedge\overrightarrow{AC} =
\begin{pmatrix}
6 \\ 3 \\ 2
\end{pmatrix}
$.
Produit mixte : $(6;3;2)\cdot(1;1;4) = 6+3+8 = 17$.
$\mathcal{V}_{ABCD} = \frac{1}{6}|17| = \frac{17}{6}$ unités de volume.

\textbf{2.} $\overrightarrow{AE}(2;2;-3)$.
Produit mixte : $(6;3;2)\cdot(2;2;-3) = 12+6-6 = 12 \neq 0$.
Les vecteurs ne sont pas coplanaires $\Rightarrow E$ n'est pas dans le plan $(ABC)$.
Volume $ABCE = \frac{12}{6} = 2$ unités de volume.

\textbf{Conclusion :} $\mathcal{V}_{ABCD}=\frac{17}{6}$, $E$ non coplanaire avec $A,B,C$.
\end{exoresolu}

\coursec{Exercice de synthèse type Bac}

\begin{exoresolu}[Problème complet : plan, distance, aire, volume]
Repère orthonormé direct $(O;\vec{i},\vec{j},\vec{k})$. $A(1;1;1)$, $B(2;-1;0)$, $C(0;1;3)$, $D(3;0;2)$.
\begin{enumerate}
\item Démontrer que $A,B,C$ définissent un plan et donner un vecteur normal $\vec{n}$.
\item Calculer l'aire du triangle $ABC$.
\item Calculer la distance de $D$ au plan $(ABC)$.
\item En déduire le volume du tétraèdre $ABCD$ par $\mathcal{V}=\frac13\mathcal{A}_{\text{base}}h$, puis vérifier par produit mixte.
\end{enumerate}
\tcblower
\textbf{1.} $\overrightarrow{AB}(1;-2;-1)$, $\overrightarrow{AC}(-1;0;2)$.
Non colinéaires ($\frac{-1}{1}=-1 \neq \frac{0}{-2}=0$) $\Rightarrow$ plan défini.
$\vec{n} = \overrightarrow{AB}\wedge\overrightarrow{AC} =
\begin{pmatrix}
(-2)(2)-(-1)(0) \\
(-1)(-1)-(1)(2) \\
(1)(0)-(-2)(-1)
\end{pmatrix}
=
\begin{pmatrix}
-4 \\ -1 \\ -2
\end{pmatrix}
$.

\textbf{2.} $\|\vec{n}\| = \sqrt{16+1+4} = \sqrt{21}$.
$\mathcal{A}_{ABC} = \frac{1}{2}\sqrt{21}$.

\textbf{3.} $\overrightarrow{AD}(2;-1;1)$.
$\vec{n}\cdot\overrightarrow{AD} = (-4)(2)+(-1)(-1)+(-2)(1) = -8+1-2 = -9$.
$d(D,(ABC)) = \frac{|-9|}{\sqrt{21}} = \frac{9}{\sqrt{21}} = \frac{3\sqrt{21}}{7}$.

\textbf{4.} Par base $\times$ hauteur :
$\mathcal{V} = \frac{1}{3}\times\frac{\sqrt{21}}{2}\times\frac{9}{\sqrt{21}} = \frac{9}{6} = \frac{3}{2}$.
Par produit mixte :
$\mathcal{V} = \frac{1}{6}\big|(\overrightarrow{AB}\wedge\overrightarrow{AC})\cdot\overrightarrow{AD}\big| = \frac{1}{6}|-9| = \frac{3}{2}$. \checkmark

\textbf{Conclusion :} Volume $= \frac{3}{2}$ unités de volume. Cohérence des deux méthodes vérifiée.
\end{exoresolu}

\begin{attention}[Les 5 erreurs qui coûtent des points au Bac]
\begin{enumerate}
\item \textbf{Oublier la condition « base orthonormée directe ».} La formule analytique (gamma/déterminant) n'est valable que dans ce cas. Citez toujours : « Dans le repère orthonormé direct donné... ». Sinon, le signe du résultat est faux.
\item \textbf{Erreur sur la 2ème coordonnée.} Formule : $(yz'-zy'\,;\,zx'-xz'\,;\,xy'-yx')$. La coordonnée du milieu est $zx'-xz'$ (et non $xz'-zx'$). \textbf{Réflexe :} Vérifier systématiquement $(\vec{u}\wedge\vec{v})\cdot\vec{u}=0$.
\item \textbf{Oublier les facteurs $\frac12$ ou $\frac16$.} Aire triangle $=\frac12\|\vec{u}\wedge\vec{v}\|$ (pas la norme seule). Volume tétraèdre $=\frac16|$produit mixte$|$ (le $\frac13$ est pour la formule base$\times$hauteur).
\item \textbf{Oublier la valeur absolue.} Distance $=\frac{|\vec{n}\cdot\overrightarrow{AM}|}{\|\vec{n}\|}$, Volume $=\frac16|$produit mixte$|$. Un résultat négatif est une erreur de cours sanctionnée.
\item \textbf{Appliquer sans vérifier les hypothèses.}
  \begin{itemize}
  \item Plan : prouver que les deux vecteurs sont non colinéaires (sinon pas de plan).
  \item Distance point-droite : vérifier que le point n'est pas sur la droite (sinon distance nulle, calcul inutile).
  \item La justification de non-colinéarité est \textbf{systématiquement attendue} dans la rédaction.
  \end{itemize}
\end{enumerate}
\end{attention}

\begin{savaistu}[Le produit vectoriel dans l'ingénierie au Cameroun]
\begin{itemize}
\item \textbf{Électrification rurale (EDC/ENEO) :} Le calcul des moments de force (couple) sur les pylônes haute tension utilise $\vec{M} = \overrightarrow{AB}\wedge\vec{F}$. La direction du couple (axe de rotation) est donnée par le produit vectoriel.
\item \textbf{Topographie / Cadastre :} L'aire d'une parcelle non plane (projetée) ou le volume de déblais/remblais pour une route (ex: route Yaoundé-Douala) se calcule par somme de tétraèdres (méthode des volumes par produit mixte).
\item \textbf{Robotique / Impression 3D (FabLabs Yaoundé/Douala) :} L'orientation de l'effecteur (pince, buse) est définie par une base directe. Le produit vectoriel permet de passer du repère outil au repère monde et de calculer les vitesses angulaires.
\end{itemize}
\end{savaistu}

\newpage

\coursec{Exercices}

\courssub{Je vérifie mes connaissances}

\begin{exercice}[QCM : les fondamentaux]
Pour chaque question, une seule réponse est exacte. Recopie la bonne réponse sur ta copie.

\textbf{1.} Dans une base orthonormée \emph{directe} $(\vec{\imath}, \vec{\jmath}, \vec{k})$, on a :
\begin{itemize}
\item[a)] $\vec{\imath} \wedge \vec{\jmath} = -\vec{k}$
\item[b)] $\vec{\imath} \wedge \vec{\jmath} = \vec{k}$
\item[c)] $\vec{\imath} \wedge \vec{\jmath} = \vec{0}$
\end{itemize}

\textbf{2.} Si $\vec{u}$ et $\vec{v}$ sont deux vecteurs colinéaires non nuls, alors $\vec{u} \wedge \vec{v}$ est égal à :
\begin{itemize}
\item[a)] $\vec{u}$ \quad b) $\|\vec{u}\|\,\|\vec{v}\|\,\vec{k}$ \quad c) $\vec{0}$
\end{itemize}

\textbf{3.} Soient $\vec{u}(1 ; 2 ; 3)$ et $\vec{v}(0 ; 1 ; -1)$. Le produit vectoriel $\vec{u} \wedge \vec{v}$ a pour coordonnées :
\begin{itemize}
\item[a)] $(-5 ; 1 ; 1)$ \quad b) $(5 ; -1 ; -1)$ \quad c) $(-5 ; -1 ; 1)$
\end{itemize}

\textbf{4.} L'aire du triangle $ABC$ est égale à :
\begin{itemize}
\item[a)] $\dfrac{1}{2}\left|\overrightarrow{AB} \cdot \overrightarrow{AC}\right|$
\item[b)] $\dfrac{1}{2}\left\|\overrightarrow{AB} \wedge \overrightarrow{AC}\right\|$
\item[c)] $\left\|\overrightarrow{AB} \wedge \overrightarrow{AC}\right\|$
\end{itemize}
\end{exercice}

\begin{exercice}[Vrai ou faux, en justifiant]
Réponds par \textbf{Vrai} ou \textbf{Faux} en justifiant chaque réponse par une propriété du cours ou un contre-exemple.
\begin{enumerate}
\item Pour tous vecteurs $\vec{u}$ et $\vec{v}$ de l'espace : $\vec{u} \wedge \vec{v} = \vec{v} \wedge \vec{u}$.
\item Pour tous vecteurs $\vec{u}$, $\vec{v}$ et $\vec{w}$ : $\vec{u} \wedge (\vec{v} + \vec{w}) = \vec{u} \wedge \vec{v} + \vec{u} \wedge \vec{w}$.
\item Le produit vectoriel de deux vecteurs non nuls et non colinéaires est un vecteur orthogonal à chacun d'eux.
\item Si $\vec{u} \wedge \vec{v} = \vec{0}$, alors $\vec{u} = \vec{0}$ ou $\vec{v} = \vec{0}$.
\item Pour tout réel $k$ et tous vecteurs $\vec{u}$ et $\vec{v}$ : $(k\vec{u}) \wedge \vec{v} = k(\vec{u} \wedge \vec{v})$.
\end{enumerate}
\end{exercice}

\begin{exercice}[Orientation et définitions]
\begin{enumerate}
\item Énonce la règle du bonhomme d'Ampère (ou de la main droite) permettant de dire qu'un trièdre $(\vec{u}, \vec{v}, \vec{w})$ est direct.
\item Donne la définition complète du produit vectoriel $\vec{u} \wedge \vec{v}$ de deux vecteurs non colinéaires (direction, sens, norme).
\item Rappelle la formule donnant la distance d'un point $M$ à une droite $D$ passant par $A$ et dirigée par $\vec{u}$, à l'aide du produit vectoriel.
\item Rappelle la formule du volume d'un tétraèdre $ABCD$ à l'aide du produit mixte.
\end{enumerate}
\end{exercice}

\begin{exercice}[Calculs directs]
L'espace est muni d'une base orthonormée directe. On donne $\vec{u}(2 ; -1 ; 3)$, $\vec{v}(1 ; 4 ; -2)$ et $\vec{w}(0 ; 1 ; 5)$.
\begin{enumerate}
\item Calculer $\vec{u} \wedge \vec{v}$, puis $\vec{v} \wedge \vec{u}$. Que constate-t-on ?
\item Calculer $\|\vec{u} \wedge \vec{v}\|$.
\item Calculer $\vec{u} \wedge (2\vec{v})$ et $(\vec{u} + \vec{w}) \wedge \vec{v}$.
\item Vérifier que $\vec{u} \wedge \vec{v}$ est orthogonal à $\vec{u}$ et à $\vec{v}$ à l'aide du produit scalaire.
\end{enumerate}
\end{exercice}

\courssub{Je m'entraîne}

\begin{exercice}[Vecteur normal et équation de plan]
L'espace est muni d'un repère orthonormé direct $(O ; \vec{\imath}, \vec{\jmath}, \vec{k})$. On considère les points $A(1 ; 0 ; 2)$, $B(3 ; 1 ; -1)$ et $C(0 ; 2 ; 1)$.
\begin{enumerate}
\item Calculer les coordonnées du vecteur $\overrightarrow{AB} \wedge \overrightarrow{AC}$.
\item Justifier que ce vecteur est normal au plan $(ABC)$.
\item En déduire une équation cartésienne du plan $(ABC)$.
\item Le point $D(5 ; 3 ; -4)$ appartient-il au plan $(ABC)$ ?
\end{enumerate}
\end{exercice}

\begin{exercice}[Distance d'un point à une droite]
Dans un repère orthonormé direct, la droite $D$ passe par le point $A(2 ; 1 ; 3)$ et est dirigée par le vecteur $\vec{u}(1 ; -2 ; 2)$. On considère le point $M(4 ; 3 ; 1)$.
\begin{enumerate}
\item Calculer $\overrightarrow{AM} \wedge \vec{u}$, puis sa norme.
\item Calculer $\|\vec{u}\|$.
\item En déduire la distance du point $M$ à la droite $D$.
\item Déterminer les coordonnées du projeté orthogonal $H$ de $M$ sur $D$ (on pourra chercher $H$ sous la forme $H(2 + t ; 1 - 2t ; 3 + 2t)$ avec $\overrightarrow{MH} \cdot \vec{u} = 0$), puis vérifier que $MH = d(M, D)$.
\end{enumerate}
\end{exercice}

\begin{exercice}[Distance d'un point à un plan]
Dans un repère orthonormé direct, on donne les points $A(1 ; 2 ; -1)$, $B(2 ; 0 ; 1)$, $C(-1 ; 1 ; 2)$ et le point $M(3 ; 1 ; 4)$.
\begin{enumerate}
\item Déterminer un vecteur $\vec{n}$ normal au plan $(ABC)$ à l'aide du produit vectoriel.
\item En déduire une équation cartésienne du plan $(ABC)$.
\item Calculer la distance du point $M$ au plan $(ABC)$.
\item Calculer la distance de l'origine $O$ au plan $(ABC)$.
\end{enumerate}
\end{exercice}

\begin{exercice}[Coplanarité et produit mixte]
Dans un repère orthonormé direct, on considère les points $A(1 ; 1 ; 1)$, $B(2 ; 3 ; 0)$, $C(0 ; 1 ; 2)$ et $D(3 ; 5 ; -1)$.
\begin{enumerate}
\item Calculer $\overrightarrow{AB} \wedge \overrightarrow{AC}$.
\item Calculer le produit mixte $\left(\overrightarrow{AB} \wedge \overrightarrow{AC}\right) \cdot \overrightarrow{AD}$.
\item Que peut-on en conclure pour les points $A$, $B$, $C$ et $D$ ? Justifier.
\item Donner une équation cartésienne du plan contenant ces quatre points.
\end{enumerate}
\end{exercice}

\courssub{Je me prépare au Bac}

\begin{exercice}[Type Bac — Plan, aire et distance]
L'espace est muni d'un repère orthonormé direct $(O ; \vec{\imath}, \vec{\jmath}, \vec{k})$. On considère les points :
\[ A(1 ; 2 ; -1), \qquad B(2 ; 0 ; 1), \qquad C(-1 ; 1 ; 2). \]
\begin{enumerate}
\item Calculer les coordonnées des vecteurs $\overrightarrow{AB}$ et $\overrightarrow{AC}$.
\item Calculer $\overrightarrow{AB} \wedge \overrightarrow{AC}$. En déduire que les points $A$, $B$ et $C$ ne sont pas alignés.
\item Justifier que le vecteur obtenu est normal au plan $(ABC)$, puis déterminer une équation cartésienne du plan $(ABC)$.
\item Calculer l'aire du triangle $ABC$.
\item Calculer la distance du point $O$ au plan $(ABC)$.
\item En déduire le volume du tétraèdre $OABC$.
\item Calculer les longueurs $AB$, $AC$ et $BC$. Le triangle $ABC$ est-il isocèle ? Rectangle ? Justifier.
\end{enumerate}
\end{exercice}

\begin{exercice}[Type Bac — Le hangar de stockage]
Une coopérative agricole de l'Ouest Cameroun construit un hangar dont la charpente est modélisée, dans un repère orthonormé direct $(O ; \vec{\imath}, \vec{\jmath}, \vec{k})$ d'unité \SI{1}{m}, par un cube $ABCDEFGH$ d'arête $a = 4$, avec :
\[ A(0;0;0),\; B(4;0;0),\; D(0;4;0),\; E(0;0;4), \]
$ABCD$ étant la face au sol et $EFGH$ la face supérieure, avec $F$ au-dessus de $B$, $G$ au-dessus de $C$ et $H$ au-dessus de $D$.
\begin{enumerate}
\item Donner les coordonnées des points $C$, $F$, $G$ et $H$.
\item On s'intéresse au triangle $BEG$, formé par trois câbles de renfort de la charpente.
\begin{enumerate}
\item Calculer les coordonnées de $\overrightarrow{EB}$ et $\overrightarrow{EG}$, puis $\overrightarrow{EB} \wedge \overrightarrow{EG}$.
\item En déduire l'aire du triangle $BEG$ en $\mathrm{m}^2$.
\item Montrer que le triangle $BEG$ est équilatéral.
\end{enumerate}
\item Calculer le produit mixte $\left(\overrightarrow{EB} \wedge \overrightarrow{EG}\right) \cdot \overrightarrow{EH}$ et en déduire le volume du tétraèdre $BEGH$.
\item Déterminer la distance du point $H$ au plan $(BEG)$.
\item Déterminer la distance du point $H$ à la droite $(EG)$.
\end{enumerate}
\end{exercice}

\begin{exercice}[Type Bac — Problème de synthèse : le tétraèdre]
L'espace est muni d'un repère orthonormé direct $(O ; \vec{\imath}, \vec{\jmath}, \vec{k})$. Un tétraèdre $ABCD$ a pour sommets :
\[ A(0 ; 0 ; 0), \qquad B(3 ; 0 ; 0), \qquad C(1 ; 2 ; 0), \qquad D(1 ; 1 ; 4). \]
\textbf{Partie A : volume et hauteur}
\begin{enumerate}
\item Calculer $\overrightarrow{AB} \wedge \overrightarrow{AC}$ et en déduire l'aire du triangle $ABC$.
\item Calculer le volume $V$ du tétraèdre $ABCD$ à l'aide du produit mixte.
\item En utilisant l'aire de la base $ABC$, retrouver la hauteur $h$ du tétraèdre issue de $D$.
\item Déterminer une équation cartésienne du plan $(ABC)$, puis calculer la distance de $D$ au plan $(ABC)$. Comparer avec la question 3.
\end{enumerate}
\textbf{Partie B : aires des faces}
\begin{enumerate}[resume]
\item Calculer l'aire du triangle $ABD$.
\item Calculer l'aire du triangle $ACD$.
\item Calculer l'aire du triangle $BCD$.
\item En déduire l'aire totale de la surface du tétraèdre.
\end{enumerate}
\textbf{Partie C : identité de Lagrange}
\begin{enumerate}[resume]
\item Démontrer que pour tous vecteurs $\vec{u}$ et $\vec{v}$ de l'espace :
\[ \|\vec{u} \wedge \vec{v}\|^2 + (\vec{u} \cdot \vec{v})^2 = \|\vec{u}\|^2 \, \|\vec{v}\|^2. \]
(On pourra poser $\theta$ une mesure de l'angle géométrique entre $\vec{u}$ et $\vec{v}$.)
\item Utiliser cette identité pour retrouver rapidement l'aire du triangle $ABC$ sans calculer de coordonnées de produit vectoriel.
\end{enumerate}
\end{exercice}



\newpage

\coursec{Corrigés des exercices}

\begin{corrige}[Exercice 1 — QCM : les fondamentaux]
\textbf{1.} Réponse \textbf{b)} : dans une base orthonormée directe $(\vec{\imath}, \vec{\jmath}, \vec{k})$, on a $\vec{\imath} \wedge \vec{\jmath} = \vec{k}$, $\vec{\jmath} \wedge \vec{k} = \vec{\imath}$ et $\vec{k} \wedge \vec{\imath} = \vec{\jmath}$ (permutations circulaires). La réponse a) correspond à une base \emph{indirecte}.

\textbf{2.} Réponse \textbf{c)} : deux vecteurs sont colinéaires si et seulement si leur produit vectoriel est nul. En effet, si $\vec{u}$ et $\vec{v}$ sont colinéaires, l'angle géométrique $\theta$ vaut $0$ ou $\pi$, donc $\sin\theta = 0$ et $\|\vec{u} \wedge \vec{v}\| = \|\vec{u}\|\,\|\vec{v}\|\sin\theta = 0$.

\textbf{3.} Réponse \textbf{a)} : avec $\vec{u}(1;2;3)$ et $\vec{v}(0;1;-1)$ :
\[ \vec{u} \wedge \vec{v} = \begin{pmatrix} 2\times(-1) - 3\times 1 \\ 3\times 0 - 1\times(-1) \\ 1\times 1 - 2\times 0 \end{pmatrix} = \begin{pmatrix} -2-3 \\ 0+1 \\ 1-0 \end{pmatrix} = \begin{pmatrix} -5 \\ 1 \\ 1 \end{pmatrix}. \]
Vérification : $(\vec{u}\wedge\vec{v})\cdot\vec{u} = -5+2+3 = 0$ et $(\vec{u}\wedge\vec{v})\cdot\vec{v} = 0+1-1 = 0$. Le résultat est bien orthogonal aux deux vecteurs.

\textbf{4.} Réponse \textbf{b)} : l'aire du triangle $ABC$ est la moitié de l'aire du parallélogramme construit sur $\overrightarrow{AB}$ et $\overrightarrow{AC}$, c'est-à-dire $\mathcal{A} = \dfrac{1}{2}\left\|\overrightarrow{AB} \wedge \overrightarrow{AC}\right\|$. La réponse a) mélange produit scalaire et produit vectoriel ; la réponse c) donne l'aire du parallélogramme, pas celle du triangle.
\end{corrige}

\begin{corrige}[Exercice 2 — Vrai ou faux, en justifiant]
\begin{enumerate}
\item \textbf{Faux.} Le produit vectoriel est \emph{antisymétrique} : $\vec{u} \wedge \vec{v} = -(\vec{v} \wedge \vec{u})$. Par exemple $\vec{\imath} \wedge \vec{\jmath} = \vec{k}$ mais $\vec{\jmath} \wedge \vec{\imath} = -\vec{k}$. L'égalité proposée n'a lieu que si $\vec{u} \wedge \vec{v} = \vec{0}$, c'est-à-dire si $\vec{u}$ et $\vec{v}$ sont colinéaires.
\item \textbf{Vrai.} C'est la distributivité du produit vectoriel sur l'addition vectorielle, conséquence directe de la bilinéarité.
\item \textbf{Vrai.} Par définition, $\vec{u} \wedge \vec{v}$ est orthogonal à $\vec{u}$ et à $\vec{v}$. L'hypothèse \og non nuls et non colinéaires \fg{} garantit en outre que ce vecteur est non nul, puisque $\sin\theta \neq 0$ pour $\theta \in \mathopen]0 ; \pi\mathclose[$.
\item \textbf{Faux.} $\vec{u} \wedge \vec{v} = \vec{0}$ signifie que $\vec{u}$ et $\vec{v}$ sont \emph{colinéaires}, pas nécessairement nuls. Contre-exemple : $\vec{u}(1;0;0)$ et $\vec{v}(2;0;0)$ sont tous deux non nuls, et pourtant $\vec{u} \wedge \vec{v} = \vec{0}$ car $\vec{v} = 2\vec{u}$.
\item \textbf{Vrai.} C'est l'homogénéité du produit vectoriel, conséquence de la bilinéarité : $(k\vec{u}) \wedge \vec{v} = \vec{u} \wedge (k\vec{v}) = k(\vec{u} \wedge \vec{v})$ pour tout réel $k$.
\end{enumerate}
\end{corrige}

\begin{corrige}[Exercice 3 — Orientation et définitions]
\begin{enumerate}
\item \textbf{Règle du bonhomme d'Ampère} : on dispose le bonhomme d'Ampère le long du premier vecteur $\vec{u}$, les pieds à l'origine commune, et il regarde dans la direction du deuxième vecteur $\vec{v}$ en tournant dans le sens du plus petit angle de $\vec{u}$ vers $\vec{v}$ ; le trièdre $(\vec{u}, \vec{v}, \vec{w})$ est direct si $\vec{w}$ pointe du côté de sa main gauche. De façon équivalente (règle de la main droite) : le pouce suivant $\vec{u}$, l'index suivant $\vec{v}$, le majeur indique alors la direction de $\vec{w}$ pour un trièdre direct.
\item Soient $\vec{u}$ et $\vec{v}$ deux vecteurs non colinéaires. Le produit vectoriel $\vec{u} \wedge \vec{v}$ est l'unique vecteur tel que :
\begin{itemize}
\item \textbf{direction} : $\vec{u} \wedge \vec{v}$ est orthogonal à $\vec{u}$ et à $\vec{v}$ ;
\item \textbf{sens} : le trièdre $(\vec{u}, \vec{v}, \vec{u} \wedge \vec{v})$ est direct ;
\item \textbf{norme} : $\|\vec{u} \wedge \vec{v}\| = \|\vec{u}\|\,\|\vec{v}\|\,\sin\theta$, où $\theta \in [0 ; \pi]$ est une mesure de l'angle géométrique entre $\vec{u}$ et $\vec{v}$.
\end{itemize}
Si $\vec{u}$ et $\vec{v}$ sont colinéaires, on pose $\vec{u} \wedge \vec{v} = \vec{0}$.
\item La distance du point $M$ à la droite $D$ passant par $A$ et de vecteur directeur $\vec{u}$ est :
\[ d(M, D) = \dfrac{\left\|\overrightarrow{AM} \wedge \vec{u}\right\|}{\|\vec{u}\|}. \]
\item Le volume du tétraèdre $ABCD$ est :
\[ V = \dfrac{1}{6}\left|\left(\overrightarrow{AB} \wedge \overrightarrow{AC}\right) \cdot \overrightarrow{AD}\right|, \]
où la quantité $\left(\overrightarrow{AB} \wedge \overrightarrow{AC}\right) \cdot \overrightarrow{AD}$ est le \cle{produit mixte} des trois vecteurs.
\end{enumerate}
\end{corrige}

\begin{corrige}[Exercice 4 — Calculs directs]
\begin{enumerate}
\item $\vec{u} \wedge \vec{v} = \begin{pmatrix} (-1)\times(-2) - 3\times 4 \\ 3\times 1 - 2\times(-2) \\ 2\times 4 - (-1)\times 1 \end{pmatrix} = \begin{pmatrix} 2-12 \\ 3+4 \\ 8+1 \end{pmatrix} = \begin{pmatrix} -10 \\ 7 \\ 9 \end{pmatrix}$.

Par antisymétrie, $\vec{v} \wedge \vec{u} = -(\vec{u} \wedge \vec{v}) = \begin{pmatrix} 10 \\ -7 \\ -9 \end{pmatrix}$. On constate bien que $\vec{v} \wedge \vec{u} = -(\vec{u} \wedge \vec{v})$ : le produit vectoriel est antisymétrique.
\item $\|\vec{u} \wedge \vec{v}\| = \sqrt{(-10)^2 + 7^2 + 9^2} = \sqrt{100 + 49 + 81} = \sqrt{230}$.
\item Par homogénéité : $\vec{u} \wedge (2\vec{v}) = 2(\vec{u} \wedge \vec{v}) = (-20 ; 14 ; 18)$.

Par distributivité : $(\vec{u} + \vec{w}) \wedge \vec{v} = \vec{u} \wedge \vec{v} + \vec{w} \wedge \vec{v}$. Or
\[ \vec{w} \wedge \vec{v} = \begin{pmatrix} 1\times(-2) - 5\times 4 \\ 5\times 1 - 0\times(-2) \\ 0\times 4 - 1\times 1 \end{pmatrix} = \begin{pmatrix} -2-20 \\ 5-0 \\ 0-1 \end{pmatrix} = \begin{pmatrix} -22 \\ 5 \\ -1 \end{pmatrix}, \]
donc $(\vec{u} + \vec{w}) \wedge \vec{v} = (-10 - 22 ; 7 + 5 ; 9 - 1) = (-32 ; 12 ; 8)$.
\item $(\vec{u} \wedge \vec{v}) \cdot \vec{u} = (-10)\times 2 + 7\times(-1) + 9\times 3 = -20 - 7 + 27 = 0$.

$(\vec{u} \wedge \vec{v}) \cdot \vec{v} = (-10)\times 1 + 7\times 4 + 9\times(-2) = -10 + 28 - 18 = 0$.

Le vecteur $\vec{u} \wedge \vec{v}$ est bien orthogonal à $\vec{u}$ et à $\vec{v}$ : cette vérification systématique est un excellent réflexe pour détecter une erreur de calcul.
\end{enumerate}
\end{corrige}

\begin{corrige}[Exercice 5 — Vecteur normal et équation de plan]
\begin{enumerate}
\item $\overrightarrow{AB}(2 ; 1 ; -3)$ et $\overrightarrow{AC}(-1 ; 2 ; -1)$. Donc
\[ \overrightarrow{AB} \wedge \overrightarrow{AC} = \begin{pmatrix} 1\times(-1) - (-3)\times 2 \\ (-3)\times(-1) - 2\times(-1) \\ 2\times 2 - 1\times(-1) \end{pmatrix} = \begin{pmatrix} -1+6 \\ 3+2 \\ 4+1 \end{pmatrix} = \begin{pmatrix} 5 \\ 5 \\ 5 \end{pmatrix}. \]
\item $\overrightarrow{AB}$ et $\overrightarrow{AC}$ sont deux vecteurs non colinéaires du plan $(ABC)$ (leur produit vectoriel est non nul). Par définition, $\overrightarrow{AB} \wedge \overrightarrow{AC}$ est orthogonal à chacun d'eux, donc à toute combinaison linéaire de ces vecteurs, c'est-à-dire à tout vecteur du plan $(ABC)$ : c'est un vecteur normal au plan $(ABC)$.
\item On peut prendre le vecteur normal simplifié $\vec{n}(1 ; 1 ; 1)$. Une équation du plan est $x + y + z + d = 0$. Comme $A(1 ; 0 ; 2) \in (ABC)$ : $1 + 0 + 2 + d = 0$, soit $d = -3$. D'où :
\[ (ABC) : x + y + z - 3 = 0. \]
Vérification rapide : $B$ : $3+1-1-3 = 0$ ; $C$ : $0+2+1-3 = 0$.
\item Pour $D(5 ; 3 ; -4)$ : $5 + 3 - 4 - 3 = 1 \neq 0$. Donc $D \notin (ABC)$.
\end{enumerate}
\end{corrige}

\begin{corrige}[Exercice 6 — Distance d'un point à une droite]
\begin{enumerate}
\item $\overrightarrow{AM}(2 ; 2 ; -2)$. Alors
\[ \overrightarrow{AM} \wedge \vec{u} = \begin{pmatrix} 2\times 2 - (-2)\times(-2) \\ (-2)\times 1 - 2\times 2 \\ 2\times(-2) - 2\times 1 \end{pmatrix} = \begin{pmatrix} 4-4 \\ -2-4 \\ -4-2 \end{pmatrix} = \begin{pmatrix} 0 \\ -6 \\ -6 \end{pmatrix}, \]
et $\left\|\overrightarrow{AM} \wedge \vec{u}\right\| = \sqrt{0 + 36 + 36} = \sqrt{72} = 6\sqrt{2}$.
\item $\|\vec{u}\| = \sqrt{1 + 4 + 4} = \sqrt{9} = 3$.
\item $d(M, D) = \dfrac{\left\|\overrightarrow{AM} \wedge \vec{u}\right\|}{\|\vec{u}\|} = \dfrac{6\sqrt{2}}{3} = 2\sqrt{2}$.
\item Un point $H$ de $D$ s'écrit $H(2 + t ; 1 - 2t ; 3 + 2t)$, donc $\overrightarrow{MH}(2 + t - 4 ; 1 - 2t - 3 ; 3 + 2t - 1) = (t - 2 ; -2 - 2t ; 2 + 2t)$. Le projeté orthogonal est caractérisé par $\overrightarrow{MH} \cdot \vec{u} = 0$ :
\begin{align*}
(t - 2)\times 1 + (-2 - 2t)\times(-2) + (2 + 2t)\times 2 &= 0 \\
t - 2 + 4 + 4t + 4 + 4t &= 0 \\
9t + 6 &= 0 \iff t = -\dfrac{2}{3}.
\end{align*}
Donc $H\left(2 - \dfrac{2}{3} ; 1 + \dfrac{4}{3} ; 3 - \dfrac{4}{3}\right)$, c'est-à-dire $H\left(\dfrac{4}{3} ; \dfrac{7}{3} ; \dfrac{5}{3}\right)$.

Vérification : $\overrightarrow{MH}\left(-\dfrac{8}{3} ; -\dfrac{2}{3} ; \dfrac{2}{3}\right)$, donc
\[ MH = \sqrt{\dfrac{64}{9} + \dfrac{4}{9} + \dfrac{4}{9}} = \sqrt{\dfrac{72}{9}} = \sqrt{8} = 2\sqrt{2} = d(M, D). \]
Les deux méthodes concordent.
\end{enumerate}
\end{corrige}

\begin{corrige}[Exercice 7 — Distance d'un point à un plan]
\begin{enumerate}
\item $\overrightarrow{AB}(1 ; -2 ; 2)$ et $\overrightarrow{AC}(-2 ; -1 ; 3)$. Donc
\[ \vec{n} = \overrightarrow{AB} \wedge \overrightarrow{AC} = \begin{pmatrix} (-2)\times 3 - 2\times(-1) \\ 2\times(-2) - 1\times 3 \\ 1\times(-1) - (-2)\times(-2) \end{pmatrix} = \begin{pmatrix} -6+2 \\ -4-3 \\ -1-4 \end{pmatrix} = \begin{pmatrix} -4 \\ -7 \\ -5 \end{pmatrix}. \]
\item Une équation du plan est $-4x - 7y - 5z + d = 0$. Avec $A(1 ; 2 ; -1)$ : $-4 - 14 + 5 + d = 0$, soit $d = 13$. En multipliant par $-1$ :
\[ (ABC) : 4x + 7y + 5z - 13 = 0. \]
\item $d(M, (ABC)) = \dfrac{|4\times 3 + 7\times 1 + 5\times 4 - 13|}{\sqrt{4^2 + 7^2 + 5^2}} = \dfrac{|12 + 7 + 20 - 13|}{\sqrt{16 + 49 + 25}} = \dfrac{26}{\sqrt{90}} = \dfrac{26}{3\sqrt{10}} = \dfrac{13\sqrt{10}}{15}$.
\item $d(O, (ABC)) = \dfrac{|4\times 0 + 7\times 0 + 5\times 0 - 13|}{\sqrt{90}} = \dfrac{13}{3\sqrt{10}} = \dfrac{13\sqrt{10}}{30}$.
\end{enumerate}
\textbf{Point de vigilance} : dans la formule de la distance, n'oublie pas la valeur absolue au numérateur, et le dénominateur est la norme du vecteur normal, $\sqrt{a^2+b^2+c^2}$, et non $\sqrt{a+b+c}$.
\end{corrige}

\begin{corrige}[Exercice 8 — Coplanarité et produit mixte]
\begin{enumerate}
\item $\overrightarrow{AB}(1 ; 2 ; -1)$ et $\overrightarrow{AC}(-1 ; 0 ; 1)$. Donc
\[ \overrightarrow{AB} \wedge \overrightarrow{AC} = \begin{pmatrix} 2\times 1 - (-1)\times 0 \\ (-1)\times(-1) - 1\times 1 \\ 1\times 0 - 2\times(-1) \end{pmatrix} = \begin{pmatrix} 2-0 \\ 1-1 \\ 0+2 \end{pmatrix} = \begin{pmatrix} 2 \\ 0 \\ 2 \end{pmatrix}. \]
\item $\overrightarrow{AD}(2 ; 4 ; -2)$. Le produit mixte vaut :
\[ \left(\overrightarrow{AB} \wedge \overrightarrow{AC}\right) \cdot \overrightarrow{AD} = 2\times 2 + 0\times 4 + 2\times(-2) = 4 - 4 = 0. \]
\item Le produit mixte est nul, donc les vecteurs $\overrightarrow{AB}$, $\overrightarrow{AC}$ et $\overrightarrow{AD}$ sont coplanaires : les quatre points $A$, $B$, $C$ et $D$ sont \textbf{coplanaires}. De façon équivalente, le tétraèdre $ABCD$ est aplati, de volume nul.
\item Le vecteur $\vec{n}(2 ; 0 ; 2)$, ou plus simplement $\vec{n}'(1 ; 0 ; 1)$, est normal au plan. Une équation est $x + z + d = 0$ ; avec $A(1 ; 1 ; 1)$ : $1 + 1 + d = 0$, soit $d = -2$. D'où :
\[ x + z - 2 = 0. \]
Vérification : $B$ : $2 + 0 - 2 = 0$ ; $C$ : $0 + 2 - 2 = 0$ ; $D$ : $3 - 1 - 2 = 0$. Les quatre points appartiennent bien à ce plan.
\end{enumerate}
\end{corrige}

\begin{corrige}[Exercice 9 — Type Bac : plan, aire et distance]
\begin{enumerate}
\item $\overrightarrow{AB}(2-1 ; 0-2 ; 1-(-1)) = (1 ; -2 ; 2)$ et $\overrightarrow{AC}(-1-1 ; 1-2 ; 2-(-1)) = (-2 ; -1 ; 3)$.
\item $\overrightarrow{AB} \wedge \overrightarrow{AC} = \begin{pmatrix} (-2)\times 3 - 2\times(-1) \\ 2\times(-2) - 1\times 3 \\ 1\times(-1) - (-2)\times(-2) \end{pmatrix} = \begin{pmatrix} -6+2 \\ -4-3 \\ -1-4 \end{pmatrix} = \begin{pmatrix} -4 \\ -7 \\ -5 \end{pmatrix}$.

Ce vecteur est non nul, donc $\overrightarrow{AB}$ et $\overrightarrow{AC}$ ne sont pas colinéaires : les points $A$, $B$ et $C$ ne sont pas alignés.
\item $\overrightarrow{AB} \wedge \overrightarrow{AC}$ est orthogonal à deux vecteurs non colinéaires du plan $(ABC)$ : c'est un vecteur normal au plan $(ABC)$. Une équation est $-4x - 7y - 5z + d = 0$ ; avec $A(1 ; 2 ; -1)$ : $-4 - 14 + 5 + d = 0$, soit $d = 13$. D'où :
\[ (ABC) : 4x + 7y + 5z - 13 = 0. \]
\item $\mathcal{A}_{ABC} = \dfrac{1}{2}\left\|\overrightarrow{AB} \wedge \overrightarrow{AC}\right\| = \dfrac{1}{2}\sqrt{16 + 49 + 25} = \dfrac{\sqrt{90}}{2} = \dfrac{3\sqrt{10}}{2}$.
\item $d(O, (ABC)) = \dfrac{|{-13}|}{\sqrt{90}} = \dfrac{13}{3\sqrt{10}} = \dfrac{13\sqrt{10}}{30}$.
\item $V_{OABC} = \dfrac{1}{3} \times \mathcal{A}_{ABC} \times d(O, (ABC)) = \dfrac{1}{3} \times \dfrac{3\sqrt{10}}{2} \times \dfrac{13}{3\sqrt{10}} = \dfrac{13}{6}$.

Vérification par le produit mixte : $\overrightarrow{AO}(-1 ; -2 ; 1)$, et $\left(\overrightarrow{AB} \wedge \overrightarrow{AC}\right) \cdot \overrightarrow{AO} = (-4)\times(-1) + (-7)\times(-2) + (-5)\times 1 = 4 + 14 - 5 = 13$, donc $V = \dfrac{|13|}{6} = \dfrac{13}{6}$. Les deux méthodes concordent.
\item $AB = \sqrt{1 + 4 + 4} = 3$ ; $AC = \sqrt{4 + 1 + 9} = \sqrt{14}$ ; $\overrightarrow{BC}(-3 ; 1 ; 1)$, donc $BC = \sqrt{9 + 1 + 1} = \sqrt{11}$.

Les trois longueurs $3$, $\sqrt{14}$ et $\sqrt{11}$ sont deux à deux distinctes : le triangle $ABC$ n'est \textbf{pas isocèle}. De plus $AB^2 + BC^2 = 9 + 11 = 20 \neq 14 = AC^2$, $AB^2 + AC^2 = 23 \neq 11 = BC^2$ et $AC^2 + BC^2 = 25 \neq 9 = AB^2$ : d'après la contraposée du théorème de Pythagore, le triangle n'est \textbf{pas rectangle}. C'est un triangle quelconque (scalène).
\end{enumerate}
\textbf{Barème indicatif (sur 5 points)} : 0,5 pt pour la question 1 ; 1 pt pour la question 2 (0,5 pt calcul + 0,5 pt justification de non-alignement) ; 1 pt pour la question 3 ; 0,5 pt pour la question 4 ; 0,5 pt pour la question 5 ; 0,75 pt pour la question 6 ; 0,75 pt pour la question 7.

\textbf{Points de vigilance} : pour la non-alignement, il faut citer explicitement le critère \og produit vectoriel non nul $\iff$ vecteurs non colinéaires \fg{} ; pour le volume, la formule $V = \frac{1}{3}\mathcal{B}\times h$ doit être écrite avant l'application numérique ; pour la nature du triangle, tester \emph{toutes} les configurations de Pythagore avant de conclure.
\end{corrige}

\begin{corrige}[Exercice 10 — Type Bac : le hangar de stockage]
\begin{enumerate}
\item $C(4 ; 4 ; 0)$, $F(4 ; 0 ; 4)$, $G(4 ; 4 ; 4)$, $H(0 ; 4 ; 4)$.
\item \begin{enumerate}
\item $\overrightarrow{EB}(4 ; 0 ; -4)$ et $\overrightarrow{EG}(4 ; 4 ; 0)$. Alors
\[ \overrightarrow{EB} \wedge \overrightarrow{EG} = \begin{pmatrix} 0\times 0 - (-4)\times 4 \\ (-4)\times 4 - 4\times 0 \\ 4\times 4 - 0\times 4 \end{pmatrix} = \begin{pmatrix} 16 \\ -16 \\ 16 \end{pmatrix}. \]
\item $\mathcal{A}_{BEG} = \dfrac{1}{2}\left\|\overrightarrow{EB} \wedge \overrightarrow{EG}\right\| = \dfrac{1}{2}\sqrt{256 + 256 + 256} = \dfrac{16\sqrt{3}}{2} = 8\sqrt{3} \ \mathrm{m}^2$.
\item $EB = \sqrt{16 + 0 + 16} = 4\sqrt{2}$ ; $EG = \sqrt{16 + 16 + 0} = 4\sqrt{2}$ ; $\overrightarrow{BG}(0 ; 4 ; 4)$, donc $BG = \sqrt{0 + 16 + 16} = 4\sqrt{2}$. Les trois côtés sont égaux : le triangle $BEG$ est équilatéral. (Ce sont trois diagonales de faces du cube, toutes de longueur $a\sqrt{2}$.)
\end{enumerate}
\item $\overrightarrow{EH}(0 ; 4 ; 0)$. Le produit mixte vaut :
\[ \left(\overrightarrow{EB} \wedge \overrightarrow{EG}\right) \cdot \overrightarrow{EH} = 16\times 0 + (-16)\times 4 + 16\times 0 = -64. \]
Donc $V_{BEGH} = \dfrac{|-64|}{6} = \dfrac{64}{6} = \dfrac{32}{3} \ \mathrm{m}^3$.
\item Le vecteur $\overrightarrow{EB} \wedge \overrightarrow{EG}(16 ; -16 ; 16)$, ou plus simplement $\vec{n}(1 ; -1 ; 1)$, est normal au plan $(BEG)$. Une équation est $x - y + z + d = 0$ ; avec $B(4 ; 0 ; 0)$ : $4 + d = 0$, soit $d = -4$. D'où $(BEG) : x - y + z - 4 = 0$.
\[ d(H, (BEG)) = \dfrac{|0 - 4 + 4 - 4|}{\sqrt{1 + 1 + 1}} = \dfrac{4}{\sqrt{3}} = \dfrac{4\sqrt{3}}{3} \ \mathrm{m}. \]
Vérification de cohérence : $V = \dfrac{1}{3} \times \mathcal{A}_{BEG} \times d(H, (BEG)) = \dfrac{1}{3} \times 8\sqrt{3} \times \dfrac{4\sqrt{3}}{3} = \dfrac{1}{3} \times \dfrac{96}{3} = \dfrac{32}{3}$ : on retrouve le volume de la question 3.
\item $d(H, (EG)) = \dfrac{\left\|\overrightarrow{EH} \wedge \overrightarrow{EG}\right\|}{\left\|\overrightarrow{EG}\right\|}$. Or
\[ \overrightarrow{EH} \wedge \overrightarrow{EG} = \begin{pmatrix} 4\times 0 - 0\times 4 \\ 0\times 4 - 0\times 0 \\ 0\times 4 - 4\times 4 \end{pmatrix} = \begin{pmatrix} 0 \\ 0 \\ -16 \end{pmatrix}, \]
de norme $16$, et $\left\|\overrightarrow{EG}\right\| = 4\sqrt{2}$. Donc
\[ d(H, (EG)) = \dfrac{16}{4\sqrt{2}} = \dfrac{4}{\sqrt{2}} = 2\sqrt{2} \ \mathrm{m}. \]
\end{enumerate}
\textbf{Barème indicatif (sur 5 points)} : 0,5 pt pour la question 1 ; 1,25 pt pour la question 2 (0,5 pt produit vectoriel, 0,25 pt aire, 0,5 pt équilatéral) ; 0,75 pt pour la question 3 ; 1,25 pt pour la question 4 (équation + distance) ; 1,25 pt pour la question 5.

\textbf{Points de vigilance} : bien choisir l'origine des vecteurs ($E$ pour le triangle $BEG$) ; dans le produit mixte, la valeur absolue est indispensable car le résultat peut être négatif (ici $-64$) ; ne pas oublier les unités ($\mathrm{m}^2$, $\mathrm{m}^3$, $\mathrm{m}$) dans un exercice contextualisé ; la vérification volume $= \frac{1}{3}\mathcal{B}h$ est un excellent moyen de contrôle à signaler au correcteur.
\end{corrige}

\begin{corrige}[Exercice 11 — Type Bac : problème de synthèse, le tétraèdre]
\textbf{Partie A : volume et hauteur}
\begin{enumerate}
\item $\overrightarrow{AB}(3 ; 0 ; 0)$ et $\overrightarrow{AC}(1 ; 2 ; 0)$. Donc
\[ \overrightarrow{AB} \wedge \overrightarrow{AC} = \begin{pmatrix} 0\times 0 - 0\times 2 \\ 0\times 1 - 3\times 0 \\ 3\times 2 - 0\times 1 \end{pmatrix} = \begin{pmatrix} 0 \\ 0 \\ 6 \end{pmatrix}. \]
Aire de $ABC$ : $\mathcal{A}_{ABC} = \dfrac{1}{2}\times 6 = 3$.
\item $\overrightarrow{AD}(1 ; 1 ; 4)$. Produit mixte : $\left(\overrightarrow{AB} \wedge \overrightarrow{AC}\right) \cdot \overrightarrow{AD} = 0 + 0 + 6\times 4 = 24$. Donc
\[ V = \dfrac{|24|}{6} = 4. \]
\item $V = \dfrac{1}{3} \times \mathcal{A}_{ABC} \times h$, donc $4 = \dfrac{1}{3}\times 3 \times h = h$. La hauteur issue de $D$ est $h = 4$.
\item Le vecteur $\overrightarrow{AB} \wedge \overrightarrow{AC}(0 ; 0 ; 6)$ est normal au plan $(ABC)$ : le plan a une équation de la forme $z = d$, et comme $A(0;0;0)$ lui appartient, $(ABC) : z = 0$ (c'est le plan $(xOy)$). Alors
\[ d(D, (ABC)) = \dfrac{|4|}{\sqrt{0^2+0^2+1^2}} = 4. \]
On retrouve bien $h = 4$ : la distance de $D$ au plan de base est la hauteur du tétraèdre.
\end{enumerate}
\textbf{Partie B : aires des faces}
\begin{enumerate}[resume]
\item $\overrightarrow{AB}(3 ; 0 ; 0)$ et $\overrightarrow{AD}(1 ; 1 ; 4)$ :
\[ \overrightarrow{AB} \wedge \overrightarrow{AD} = \begin{pmatrix} 0\times 4 - 0\times 1 \\ 0\times 1 - 3\times 4 \\ 3\times 1 - 0\times 1 \end{pmatrix} = \begin{pmatrix} 0 \\ -12 \\ 3 \end{pmatrix}, \quad \mathcal{A}_{ABD} = \dfrac{1}{2}\sqrt{144 + 9} = \dfrac{\sqrt{153}}{2} = \dfrac{3\sqrt{17}}{2}. \]
\item $\overrightarrow{AC}(1 ; 2 ; 0)$ et $\overrightarrow{AD}(1 ; 1 ; 4)$ :
\[ \overrightarrow{AC} \wedge \overrightarrow{AD} = \begin{pmatrix} 2\times 4 - 0\times 1 \\ 0\times 1 - 1\times 4 \\ 1\times 1 - 2\times 1 \end{pmatrix} = \begin{pmatrix} 8 \\ -4 \\ -1 \end{pmatrix}, \quad \mathcal{A}_{ACD} = \dfrac{1}{2}\sqrt{64 + 16 + 1} = \dfrac{\sqrt{81}}{2} = \dfrac{9}{2}. \]
\item $\overrightarrow{BC}(-2 ; 2 ; 0)$ et $\overrightarrow{BD}(-2 ; 1 ; 4)$ :
\[ \overrightarrow{BC} \wedge \overrightarrow{BD} = \begin{pmatrix} 2\times 4 - 0\times 1 \\ 0\times(-2) - (-2)\times 4 \\ (-2)\times 1 - 2\times(-2) \end{pmatrix} = \begin{pmatrix} 8 \\ 8 \\ 2 \end{pmatrix}, \quad \mathcal{A}_{BCD} = \dfrac{1}{2}\sqrt{64 + 64 + 4} = \dfrac{\sqrt{132}}{2} = \sqrt{33}. \]
\item Aire totale :
\[ \mathcal{A}_{\text{totale}} = 3 + \dfrac{3\sqrt{17}}{2} + \dfrac{9}{2} + \sqrt{33} = \dfrac{15 + 3\sqrt{17}}{2} + \sqrt{33}. \]
\end{enumerate}
\textbf{Partie C : identité de Lagrange}
\begin{enumerate}[resume]
\item Soit $\theta \in [0 ; \pi]$ une mesure de l'angle géométrique entre $\vec{u}$ et $\vec{v}$. Par définition :
\[ \|\vec{u} \wedge \vec{v}\| = \|\vec{u}\|\,\|\vec{v}\|\,\sin\theta \qquad \text{et} \qquad \vec{u} \cdot \vec{v} = \|\vec{u}\|\,\|\vec{v}\|\,\cos\theta. \]
(Comme $\theta \in [0;\pi]$, on a $\sin\theta \geq 0$, ce qui légitime la première égalité.) Alors :
\begin{align*}
\|\vec{u} \wedge \vec{v}\|^2 + (\vec{u} \cdot \vec{v})^2 &= \|\vec{u}\|^2\,\|\vec{v}\|^2\sin^2\theta + \|\vec{u}\|^2\,\|\vec{v}\|^2\cos^2\theta \\
&= \|\vec{u}\|^2\,\|\vec{v}\|^2\left(\sin^2\theta + \cos^2\theta\right) = \|\vec{u}\|^2\,\|\vec{v}\|^2.
\end{align*}
\item Pour le triangle $ABC$ : $\overrightarrow{AB} \cdot \overrightarrow{AC} = 3\times 1 + 0\times 2 + 0\times 0 = 3$, $\left\|\overrightarrow{AB}\right\|^2 = 9$ et $\left\|\overrightarrow{AC}\right\|^2 = 1 + 4 = 5$. Par l'identité de Lagrange :
\[ \left\|\overrightarrow{AB} \wedge \overrightarrow{AC}\right\|^2 = 9\times 5 - 3^2 = 45 - 9 = 36, \]
donc $\left\|\overrightarrow{AB} \wedge \overrightarrow{AC}\right\| = 6$ et $\mathcal{A}_{ABC} = \dfrac{6}{2} = 3$. On retrouve le résultat de la question 1.
\end{enumerate}
\textbf{Barème indicatif (sur 7 points)} : Partie A : 2,5 pts (0,75 pt pour la question 1 ; 0,75 pt pour la question 2 ; 0,5 pt pour la question 3 ; 0,5 pt pour la question 4) ; Partie B : 2,5 pts (0,5 pt par face + 0,5 pt pour l'aire totale) ; Partie C : 2 pts (1,25 pt pour la démonstration, 0,75 pt pour l'application).

\textbf{Points de vigilance} : la démonstration de l'identité de Lagrange est une restitution de connaissances exigible : il faut citer les deux définitions (norme du produit vectoriel avec $\sin\theta$, produit scalaire avec $\cos\theta$) et conclure par $\sin^2\theta + \cos^2\theta = 1$ ; pour l'aire totale, laisser le résultat sous forme exacte avec radicaux simplifiés ($\sqrt{153} = 3\sqrt{17}$, $\sqrt{132} = 2\sqrt{33}$) ; dans la Partie B, attention à bien prendre des vecteurs de \emph{même origine} pour chaque face (par exemple $\overrightarrow{BC}$ et $\overrightarrow{BD}$ pour la face $BCD$).
\end{corrige}

\newpage

\coursec{Fiche bilan}

\coursec{Orientation de l'espace et repérage}

\begin{experience}[La règle du bonhomme d'Ampère (main droite)]
Plongeons-nous dans une situation concrète : imaginons un pylône de télécommunication \textbf{MTN} ou \textbf{Orange} installé sur le toit d'un immeuble à Yaoundé. Pour orienter les antennes, les techniciens utilisent un repère spatial. Comment s'assurer que tout le monde parle du même « haut », « avant », « droite » ?
\begin{enumerate}
  \item Ouvre ta main \textbf{droite}. Tend le pouce, l'index et le majeur de façon à ce qu'ils soient deux à deux perpendiculaires.
  \item Le \textbf{pouce} pointe vers l'axe $Ox$ (vecteur $\vec{\imath}$), l'\textbf{index} vers $Oy$ (vecteur $\vec{\jmath}$), le \textbf{majeur} vers $Oz$ (vecteur $\vec{k}$).
  \item Si tu arrives à placer tes doigts ainsi sans te tordre le poignet de façon douloureuse, le repère $(\vec{\imath},\vec{\jmath},\vec{k})$ est dit \cle{direct}.
  \item Essaie avec la main gauche : c'est impossible de superposer les trois doigts sur les trois axes dans le même ordre. Ce repère serait \cle{indirect}.
\end{enumerate}
C'est la \cle{règle de la main droite} (ou règle du bonhomme d'Ampère). Elle fixe une convention universelle pour l'orientation de l'espace.
\end{experience}

\begin{definition}[Base orthonormée directe]
Soit $(\vec{\imath},\vec{\jmath},\vec{k})$ une base de l'espace.
\begin{itemize}
  \item Elle est \cle{orthonormée} si les trois vecteurs sont deux à deux orthogonaux et de norme $1$ : $\|\vec{\imath}\|=\|\vec{\jmath}\|=\|\vec{k}\|=1$ et $\vec{\imath}\cdot\vec{\jmath}=\vec{\jmath}\cdot\vec{k}=\vec{k}\cdot\vec{\imath}=0$.
  \item Elle est \cle{directe} si le triplet $(\vec{\imath},\vec{\jmath},\vec{k})$ satisfait la règle de la main droite.
\end{itemize}
Dans tout le chapitre, sauf mention contraire, nous travaillerons dans un repère \textbf{orthonormé direct} $(O;\vec{\imath},\vec{\jmath},\vec{k})$.
\end{definition}

\begin{attention}[Piège : Base non orthonormée ou non directe]
Les formules analytiques du produit vectoriel (coordonnées, déterminant) ne sont valables \textbf{que dans une base orthonormée directe}. Si l'énoncé donne une base orthonormée indirecte, le signe du produit vectoriel change. Si la base n'est pas orthonormée, il faut passer par la définition géométrique ou changer de base. \textbf{Vérifie toujours la nature de la base au début de l'exercice.}
\end{attention}

\begin{popfigure}
\begin{tikzpicture}[>=Stealth, scale=1.2]
% Axes
\draw[->, thick, popblue] (0,0,0) -- (3,0,0) node[anchor=north east] {$x$};
\draw[->, thick, poporange] (0,0,0) -- (0,3,0) node[anchor=north west] {$y$};
\draw[->, thick, popgreen] (0,0,0) -- (0,0,3) node[anchor=south] {$z$};
% Right hand schematic
\begin{scope}[shift={(3.5,0,0)}, scale=0.8]
\draw[thick, popdark] (0,0) -- (1,0.2) -- (1.5,0) -- (1, -0.2) -- cycle; % Palm
\draw[thick, popdark] (1,0.2) -- (1.8, 1) node[above right, font=\small] {Index $\vec{\jmath}$ ($Oy$)};
\draw[thick, popdark] (1.5,0) -- (2.5, 0) node[right, font=\small] {Majeur $\vec{k}$ ($Oz$)};
\draw[thick, popdark] (1,-0.2) -- (1.8, -0.8) node[below right, font=\small] {Pouce $\vec{\imath}$ ($Ox$)};
\end{scope}
% Basis vectors labels
\node[popblue] at (3.2,0,0) {$\vec{\imath}$};
\node[poporange] at (0,3.2,0) {$\vec{\jmath}$};
\node[popgreen] at (0,0,3.2) {$\vec{k}$};
\end{tikzpicture}
\legende{Repère orthonormé direct et règle de la main droite.}
\end{popfigure}

\coursec{Définition et interprétation géométrique du produit vectoriel}

\begin{definition}[Produit vectoriel de deux vecteurs]
Soient $\vec{u}$ et $\vec{v}$ deux vecteurs de l'espace.
\begin{itemize}
  \item Si $\vec{u}$ et $\vec{v}$ sont \cle{colinéaires} (l'un est multiple de l'autre, ou l'un est nul), on définit $\vec{u}\wedge\vec{v}=\vec{0}$.
  \item Si $\vec{u}$ et $\vec{v}$ sont \textbf{non colinéaires}, le \cle{produit vectoriel} $\vec{u}\wedge\vec{v}$ est l'unique vecteur $\vec{w}$ tel que :
  \begin{enumerate}
    \item $\vec{w}$ est \textbf{orthogonal} à $\vec{u}$ et à $\vec{v}$ (donc normal au plan dirigé par $\vec{u}$ et $\vec{v}$).
    \item La norme de $\vec{w}$ est $\|\vec{w}\| = \|\vec{u}\|\,\|\vec{v}\|\,|\sin(\vec{u},\vec{v})|$.
    \item Le triplet $(\vec{u},\vec{v},\vec{w})$ est \textbf{direct} (règle de la main droite : pouce $\vec{u}$, index $\vec{v}$, majeur $\vec{w}$).
  \end{enumerate}
\end{itemize}
\end{definition}

\begin{propriete}[Interprétation géométrique : Aire du parallélogramme]
Soient $\vec{u}$ et $\vec{v}$ deux vecteurs non nuls. La norme de leur produit vectoriel est égale à l'aire du parallélogramme construit sur $\vec{u}$ et $\vec{v}$ :
\[
\mathcal{A}_{\text{parallélogramme}} = \|\vec{u}\wedge\vec{v}\| = \|\vec{u}\|\,\|\vec{v}\|\,|\sin(\vec{u},\vec{v})|.
\]
\emph{Preuve :} L'aire d'un parallélogramme est base $\times$ hauteur. Base $=\|\vec{u}\|$, hauteur $=\|\vec{v}\||\sin\theta|$. CQFD.
\end{propriete}

\begin{exemplebox}[Cas particuliers dans le repère canonique]
Dans la base orthonormée directe $(\vec{\imath},\vec{\jmath},\vec{k})$ :
\begin{align*}
\vec{\imath}\wedge\vec{\jmath} &= \vec{k}, &\quad \vec{\jmath}\wedge\vec{k} &= \vec{\imath}, &\quad \vec{k}\wedge\vec{\imath} &= \vec{\jmath}, \\
\vec{\jmath}\wedge\vec{\imath} &= -\vec{k}, &\quad \vec{k}\wedge\vec{\jmath} &= -\vec{\imath}, &\quad \vec{\imath}\wedge\vec{k} &= -\vec{\jmath}.
\end{align*}
On retient le cycle direct $\vec{\imath}\to\vec{\jmath}\to\vec{k}\to\vec{\imath}$ (signe $+$) et le cycle inverse (signe $-$).
\end{exemplebox}

\begin{attention}[Produit vectoriel vs Produit scalaire]
\begin{itemize}
  \item Produit scalaire $\vec{u}\cdot\vec{v}$ : donne un \textbf{réel} (mesure l'alignement, $\cos$).
  \item Produit vectoriel $\vec{u}\wedge\vec{v}$ : donne un \textbf{vecteur} (mesure l'orthogonalité, $\sin$, orientation).
  \item Ne confonds jamais les notations : $\cdot$ (point) vs $\wedge$ (caret / chapeau).
\end{itemize}
\end{attention}

\coursec{Expression analytique dans une base orthonormée directe}

\begin{propriete}[Coordonnées du produit vectoriel]
Soient $\vec{u}\begin{pmatrix}a\\b\\c\end{pmatrix}$ et $\vec{v}\begin{pmatrix}a'\\b'\\c'\end{pmatrix}$ dans une base orthonormée directe.
\[
\vec{u}\wedge\vec{v} = \begin{pmatrix}
bc' - cb' \\
ca' - ac' \\
ab' - ba'
\end{pmatrix}
\]
\emph{Mnémotechnique (déterminant formel $3\times 3$) :}
\[
\vec{u}\wedge\vec{v} = \det\begin{pmatrix}
\vec{\imath} & \vec{\jmath} & \vec{k} \\
a & b & c \\
a' & b' & c'
\end{pmatrix}
= \vec{\imath}\begin{vmatrix}b&c\\b'&c'\end{vmatrix}
- \vec{\jmath}\begin{vmatrix}a&c\\a'&c'\end{vmatrix}
+ \vec{k}\begin{vmatrix}a&b\\a'&b'\end{vmatrix}
\]
\emph{Attention au signe du second composant :} c'est $ca'-ac'$ (et non $ac'-ca'$). Le développement du déterminant donne $-(ac'-ca') = ca'-ac'$.
\end{propriete}

\begin{methode}[Calculer $\vec{u}\wedge\vec{v}$ par coordonnées]
\begin{enumerate}
  \item Vérifier que la base est orthonormée \textbf{directe}.
  \item Écrire les coordonnées : $\vec{u}(a;b;c)$, $\vec{v}(a';b';c')$.
  \item Appliquer la formule « en croix » :
  \begin{itemize}
    \item 1\textsuperscript{ère} coord. : $b c' - c b'$ (on « oublie » $a,a'$, produit croisé $b,c$).
    \item 2\textsuperscript{ème} coord. : $c a' - a c'$ (on « oublie » $b,b'$, produit croisé $c,a$ \textbf{dans cet ordre}).
    \item 3\textsuperscript{ème} coord. : $a b' - b a'$ (on « oublie » $c,c'$, produit croisé $a,b$).
  \end{itemize}
  \item Vérifier l'orthogonalité : $(\vec{u}\wedge\vec{v})\cdot\vec{u}=0$ et $(\vec{u}\wedge\vec{v})\cdot\vec{v}=0$ (contrôle rapide).
\end{enumerate}
\end{methode}

\begin{exoresolu}[Calcul analytique]
Soient $\vec{u}\begin{pmatrix}2\\-1\\3\end{pmatrix}$ et $\vec{v}\begin{pmatrix}0\\4\\-2\end{pmatrix}$ dans une base orthonormée directe. Calculer $\vec{u}\wedge\vec{v}$.
\tcblower
\textbf{Solution.}
On applique la formule :
\begin{align*}
\vec{u}\wedge\vec{v} &= \begin{pmatrix}
(-1)(-2) - (3)(4) \\
(3)(0) - (2)(-2) \\
(2)(4) - (-1)(0)
\end{pmatrix} \\
&= \begin{pmatrix}
2 - 12 \\
0 + 4 \\
8 - 0
\end{pmatrix} = \begin{pmatrix}-10\\4\\8\end{pmatrix}.
\end{align*}
\textbf{Vérification (orthogonalité) :}
$\begin{pmatrix}-10\\4\\8\end{pmatrix}\cdot\begin{pmatrix}2\\-1\\3\end{pmatrix} = -20 -4 + 24 = 0$ \quad $\checkmark$ \\
$\begin{pmatrix}-10\\4\\8\end{pmatrix}\cdot\begin{pmatrix}0\\4\\-2\end{pmatrix} = 0 + 16 - 16 = 0$ \quad $\checkmark$
\end{exoresolu}

\coursec{Propriétés algébriques du produit vectoriel}

\begin{propriete}[Antisymétrie]
Pour tous vecteurs $\vec{u},\vec{v}$ :
\[
\vec{v}\wedge\vec{u} = -(\vec{u}\wedge\vec{v}).
\]
\emph{Conséquence :} $\vec{u}\wedge\vec{u}=\vec{0}$.
\emph{Preuve géométrique :} $(\vec{u},\vec{v},\vec{u}\wedge\vec{v})$ direct $\Rightarrow$ $(\vec{v},\vec{u},\vec{u}\wedge\vec{v})$ indirect $\Rightarrow$ $(\vec{v},\vec{u},-(\vec{u}\wedge\vec{v}))$ direct.
\end{propriete}

\begin{propriete}[Bilinéarité]
Pour tous vecteurs $\vec{u},\vec{v},\vec{w}$ et réels $\lambda,\mu$ :
\begin{align*}
(\lambda\vec{u}+\mu\vec{v})\wedge\vec{w} &= \lambda(\vec{u}\wedge\vec{w}) + \mu(\vec{v}\wedge\vec{w}), \\
\vec{w}\wedge(\lambda\vec{u}+\mu\vec{v}) &= \lambda(\vec{w}\wedge\vec{u}) + \mu(\vec{w}\wedge\vec{v}).
\end{align*}
\emph{Le produit vectoriel est linéaire par rapport à chacune de ses variables.}
\end{propriete}

\begin{attention}[Non-associativité !]
Le produit vectoriel \textbf{n'est pas associatif} :
\[
(\vec{u}\wedge\vec{v})\wedge\vec{w} \neq \vec{u}\wedge(\vec{v}\wedge\vec{v}) \quad \text{en général}.
\]
Exemple : $(\vec{\imath}\wedge\vec{\jmath})\wedge\vec{\jmath} = \vec{k}\wedge\vec{\jmath} = -\vec{\imath}$, alors que $\vec{\imath}\wedge(\vec{\jmath}\wedge\vec{\jmath}) = \vec{\imath}\wedge\vec{0} = \vec{0}$.
\textbf{Ne jamais écrire $\vec{u}\wedge\vec{v}\wedge\vec{w}$ sans parenthèses.}
\end{attention}

\begin{propriete}[Produit vectoriel nul et colinéarité]
\[
\vec{u}\wedge\vec{v} = \vec{0} \quad\Longleftrightarrow\quad \vec{u} \text{ et } \vec{v} \text{ sont colinéaires}.
\]
\emph{Preuve :} Si colinéaires, $\sin\theta=0$ donc norme nulle. Réciproquement, si norme nulle, $\|\vec{u}\|\|\vec{v}\||\sin\theta|=0$, donc $\vec{u}=\vec{0}$ ou $\vec{v}=\vec{0}$ ou $\sin\theta=0$ (colinéaires).
\end{propriete}

\begin{propriete}[Orthogonalité]
$\vec{u}\wedge\vec{v}$ est orthogonal à $\vec{u}$ et à $\vec{v}$.
\[
(\vec{u}\wedge\vec{v})\cdot\vec{u} = 0 \quad\text{et}\quad (\vec{u}\wedge\vec{v})\cdot\vec{v} = 0.
\]
C'est la propriété fondamentale pour trouver un \cle{vecteur normal} à un plan.
\end{propriete}

\begin{exoresolu}[Utilisation des propriétés pour simplifier]
Soient $\vec{u},\vec{v},\vec{w}$ tels que $\vec{u}\wedge\vec{v}=\begin{pmatrix}1\\2\\-1\end{pmatrix}$, $\vec{u}\wedge\vec{w}=\begin{pmatrix}0\\1\\3\end{pmatrix}$.
Calculer $(2\vec{u}-3\vec{v})\wedge(\vec{v}+4\vec{w})$.
\tcblower
\textbf{Solution.}
On développe par bilinéarité et antisymétrie :
\begin{align*}
(2\vec{u}-3\vec{v})\wedge(\vec{v}+4\vec{w})
&= 2\vec{u}\wedge\vec{v} + 8\vec{u}\wedge\vec{w} - 3\vec{v}\wedge\vec{v} - 12\vec{v}\wedge\vec{w} \\
&= 2\vec{u}\wedge\vec{v} + 8\vec{u}\wedge\vec{w} - \vec{0} + 12\vec{w}\wedge\vec{v} \quad (\text{car } \vec{v}\wedge\vec{v}=\vec{0}, \vec{v}\wedge\vec{w}=-\vec{w}\wedge\vec{v}) \\
&= 2\begin{pmatrix}1\\2\\-1\end{pmatrix} + 8\begin{pmatrix}0\\1\\3\end{pmatrix} + 12(\vec{w}\wedge\vec{v}).
\end{align*}
On ne peut pas aller plus loin sans connaître $\vec{w}\wedge\vec{v}$ (ou $\vec{v}\wedge\vec{w}$). Si l'énoncé ne le donne pas, on s'arrête là.
\end{exoresolu}

\coursec{Applications : Vecteur normal, équation de plan, distances}

\courssub{Vecteur normal à un plan et équation cartésienne}

\begin{propriete}[Vecteur normal à un plan]
Soit un plan $\mathcal{P}$ dirigé par deux vecteurs non colinéaires $\vec{u}$ et $\vec{v}$.
Alors $\vec{n} = \vec{u}\wedge\vec{v}$ est un \textbf{vecteur normal} à $\mathcal{P}$.
Tout vecteur normal à $\mathcal{P}$ est colinéaire à $\vec{u}\wedge\vec{v}$.
\end{propriete}

\begin{methode}[Équation cartésienne d'un plan $\mathcal{P}=(ABC)$]
\begin{enumerate}
  \item Calculer deux vecteurs directeurs non colinéaires, ex : $\overrightarrow{AB}$ et $\overrightarrow{AC}$.
  \item Calculer $\vec{n} = \overrightarrow{AB}\wedge\overrightarrow{AC} \begin{pmatrix}a\\b\\c\end{pmatrix}$.
  \item L'équation cartésienne est de la forme $ax+by+cz+d=0$.
  \item Déterminer $d$ en remplaçant $x,y,z$ par les coordonnées d'un point du plan (ex: $A(x_A,y_A,z_A)$) : $d = -(ax_A+by_A+cz_A)$.
\end{enumerate}
\end{methode}

\begin{exoresolu}[Équation de plan par produit vectoriel]
Donner une équation cartésienne du plan $\mathcal{P}$ passant par $A(1;2;-1)$, $B(0;1;3)$, $C(2;0;1)$.
\tcblower
\textbf{Solution.}
\begin{enumerate}
  \item $\overrightarrow{AB}=\begin{pmatrix}-1\\-1\\4\end{pmatrix}$, $\overrightarrow{AC}=\begin{pmatrix}1\\-2\\2\end{pmatrix}$.
  \item $\vec{n} = \overrightarrow{AB}\wedge\overrightarrow{AC} = \begin{pmatrix} (-1)(2) - (4)(-2) \\ (4)(1) - (-1)(2) \\ (-1)(-2) - (-1)(1) \end{pmatrix} = \begin{pmatrix} -2+8 \\ 4+2 \\ 2+1 \end{pmatrix} = \begin{pmatrix}6\\6\\3\end{pmatrix}$.
  \item On peut simplifier $\vec{n}$ par $3$ : $\vec{n}'=\begin{pmatrix}2\\2\\1\end{pmatrix}$. Équation : $2x+2y+z+d=0$.
  \item Point $A(1;2;-1)$ : $2(1)+2(2)+(-1)+d=0 \Rightarrow 2+4-1+d=0 \Rightarrow d=-5$.
  \item \textbf{Réponse :} $\mathcal{P} : 2x+2y+z-5=0$.
\end{enumerate}
\end{exoresolu}

\courssub{Distance d'un point à une droite}

\begin{propriete}[Distance point-droite]
Soit $(D)$ une droite de vecteur directeur $\vec{u}\neq\vec{0}$ passant par $A$. Soit $M$ un point de l'espace.
La distance de $M$ à $(D)$ est :
\[
d(M,(D)) = \frac{\|\overrightarrow{AM}\wedge\vec{u}\|}{\|\vec{u}\|}.
\]
\emph{Preuve :} Aire du parallélogramme $= \|\overrightarrow{AM}\wedge\vec{u}\| = \text{base}\times\text{hauteur} = \|\vec{u}\|\times d(M,(D))$.
\end{propriete}

\begin{methode}[Calculer $d(M,(D))$]
\begin{enumerate}
  \item Identifier un point $A$ sur $(D)$ et un vecteur directeur $\vec{u}$.
  \item Calculer $\overrightarrow{AM}$.
  \item Calculer $\vec{w} = \overrightarrow{AM}\wedge\vec{u}$.
  \item Calculer $\|\vec{w}\|$ et $\|\vec{u}\|$.
  \item Appliquer la formule : $d = \frac{\|\vec{w}\|}{\|\vec{u}\|}$.
\end{enumerate}
\end{methode}

\begin{exoresolu}[Distance point-droite]
$(D)$ passe par $A(1;0;2)$ et a pour directeur $\vec{u}\begin{pmatrix}1\\2\\2\end{pmatrix}$. $M(3;1;-1)$. Calculer $d(M,(D))$.
\tcblower
\textbf{Solution.}
$\overrightarrow{AM}=\begin{pmatrix}2\\1\\-3\end{pmatrix}$.
$\overrightarrow{AM}\wedge\vec{u} = \begin{pmatrix} (1)(2) - (-3)(2) \\ (-3)(1) - (2)(2) \\ (2)(2) - (1)(1) \end{pmatrix} = \begin{pmatrix} 2+6 \\ -3-4 \\ 4-1 \end{pmatrix} = \begin{pmatrix}8\\-7\\3\end{pmatrix}$.
$\|\overrightarrow{AM}\wedge\vec{u}\| = \sqrt{64+49+9} = \sqrt{122}$.
$\|\vec{u}\| = \sqrt{1+4+4} = 3$.
$d(M,(D)) = \frac{\sqrt{122}}{3}$.
\end{exoresolu}

\courssub{Distance d'un point à un plan}

\begin{propriete}[Distance point-plan]
Soit $\mathcal{P}$ un plan d'équation cartésienne $ax+by+cz+d=0$ (avec $(a,b,c)\neq(0,0,0)$).
Soit $M(x_M;y_M;z_M)$ un point.
La distance de $M$ à $\mathcal{P}$ est :
\[
d(M,\mathcal{P}) = \frac{|ax_M+by_M+cz_M+d|}{\sqrt{a^2+b^2+c^2}}.
\]
\emph{Preuve :} $\vec{n}(a;b;c)$ est normal à $\mathcal{P}$. Pour $H$ projeté orthogonal de $M$ sur $\mathcal{P}$, $\overrightarrow{MH}$ est colinéaire à $\vec{n}$. On projette $\overrightarrow{AM}$ ($A\in\mathcal{P}$) sur $\vec{n}$.
\end{propriete}

\begin{attention}[Condition d'application de la formule point-plan]
L'équation du plan \textbf{doit} être sous forme cartésienne $ax+by+cz+d=0$ (second membre nul). Si l'équation est $ax+by+cz = k$, réécris-la $ax+by+cz-k=0$ (donc $d=-k$) avant d'appliquer la formule.
\end{attention}

\begin{exoresolu}[Distance point-plan]
$\mathcal{P} : 2x-y+2z-3=0$, $M(1;2;-1)$. Calculer $d(M,\mathcal{P})$.
\tcblower
\textbf{Solution.}
Numérateur : $|2(1) - 1(2) + 2(-1) - 3| = |2 - 2 - 2 - 3| = |-5| = 5$.
Dénominateur : $\sqrt{2^2+(-1)^2+2^2} = \sqrt{4+1+4} = 3$.
$d(M,\mathcal{P}) = \frac{5}{3}$.
\end{exoresolu}

\coursec{Aires, produit mixte et volumes}

\courssub{Aire d'un triangle et d'un parallélogramme}

\begin{propriete}[Aire d'un triangle]
Soit $ABC$ un triangle non dégénéré.
\[
\mathcal{A}_{ABC} = \frac{1}{2}\|\overrightarrow{AB}\wedge\overrightarrow{AC}\|.
\]
\emph{Preuve :} Le triangle est la moitié du parallélogramme construit sur $\overrightarrow{AB}$ et $\overrightarrow{AC}$.
\end{propriete}

\begin{exoresolu}[Aire triangle]
$A(1;0;0)$, $B(0;2;0)$, $C(0;0;3)$. Calculer l'aire du triangle $ABC$.
\tcblower
\textbf{Solution.}
$\overrightarrow{AB}=\begin{pmatrix}-1\\2\\0\end{pmatrix}$, $\overrightarrow{AC}=\begin{pmatrix}-1\\0\\3\end{pmatrix}$.
$\overrightarrow{AB}\wedge\overrightarrow{AC} = \begin{pmatrix} 2\cdot3 - 0\cdot0 \\ 0\cdot(-1) - (-1)\cdot3 \\ (-1)\cdot0 - 2\cdot(-1) \end{pmatrix} = \begin{pmatrix}6\\3\\2\end{pmatrix}$.
$\|\overrightarrow{AB}\wedge\overrightarrow{AC}\| = \sqrt{36+9+4} = \sqrt{49} = 7$.
$\mathcal{A}_{ABC} = \frac{1}{2}\times 7 = 3,5$ unités d'aire.
\end{exoresolu}

\courssub{Produit mixte et volume du tétraèdre}

\begin{definition}[Produit mixte]
Soient $\vec{u},\vec{v},\vec{w}$ trois vecteurs. Le \cle{produit mixte} de $\vec{u},\vec{v},\vec{w}$ (dans cet ordre) est le réel :
\[
[\vec{u},\vec{v},\vec{w}] = (\vec{u}\wedge\vec{v})\cdot\vec{w}.
\]
Dans une base orthonormée directe, il est égal au déterminant des coordonnées :
\[
[\vec{u},\vec{v},\vec{w}] = \det(\vec{u},\vec{v},\vec{w}) = \begin{vmatrix}
a & a' & a'' \\
b & b' & b'' \\
c & c' & c''
\end{vmatrix}.
\]
\end{definition}

\begin{propriete}[Propriétés du produit mixte]
\begin{itemize}
  \item \textbf{Invariance par permutation circulaire :} $[\vec{u},\vec{v},\vec{w}] = [\vec{v},\vec{w},\vec{u}] = [\vec{w},\vec{u},\vec{v}]$.
  \item \textbf{Changement de signe par transposition :} $[\vec{v},\vec{u},\vec{w}] = -[\vec{u},\vec{v},\vec{w}]$.
  \item \textbf{Zéro $\Leftrightarrow$ Coplanarité :} $[\vec{u},\vec{v},\vec{w}] = 0 \Longleftrightarrow \vec{u},\vec{v},\vec{w}$ sont coplanaires (vecteurs libres) ou points $A,B,C,D$ coplanaires pour $\overrightarrow{AB},\overrightarrow{AC},\overrightarrow{AD}$.
\end{itemize}
\end{propriete}

\begin{propriete}[Volumes]
\begin{itemize}
  \item Volume du parallélépipède construit sur $\vec{u},\vec{v},\vec{w}$ : $\mathcal{V}_{\text{parall}} = |[\vec{u},\vec{v},\vec{w}]| = |\det(\vec{u},\vec{v},\vec{w})|$.
  \item Volume du tétraèdre $ABCD$ : $\mathcal{V}_{ABCD} = \frac{1}{6}|\det(\overrightarrow{AB},\overrightarrow{AC},\overrightarrow{AD})|$.
\end{itemize}
\emph{Preuve :} Volume = Aire base $\times$ hauteur. Base = parallélogramme $\|\vec{u}\wedge\vec{v}\|$. Hauteur = projection de $\vec{w}$ sur la normale $\frac{|(\vec{u}\wedge\vec{v})\cdot\vec{w}|}{\|\vec{u}\wedge\vec{v}\|}$. Produit = $|(\vec{u}\wedge\vec{v})\cdot\vec{w}|$. Tétraèdre = 1/6 du parallélépipède.
\end{propriete}

\begin{methode}[Volume d'un tétraèdre $ABCD$]
\begin{enumerate}
  \item Choisir un sommet (ex: $A$) et former les trois vecteurs $\overrightarrow{AB},\overrightarrow{AC},\overrightarrow{AD}$.
  \item Calculer le déterminant $3\times 3$ de leurs coordonnées.
  \item Volume $= \frac{1}{6}|\text{déterminant}|$.
  \item Si le volume est nul, les 4 points sont coplanaires.
\end{enumerate}
\end{methode}

\begin{exoresolu}[Volume tétraèdre et coplanarité]
$A(1;1;1)$, $B(2;3;4)$, $C(0;1;-1)$, $D(1;2;1)$.
1. Calculer le volume du tétraèdre $ABCD$.
2. En déduire si $A,B,C,D$ sont coplanaires.
\tcblower
\textbf{Solution.}
1. Vecteurs :
$\overrightarrow{AB}=\begin{pmatrix}1\\2\\3\end{pmatrix}$,
$\overrightarrow{AC}=\begin{pmatrix}-1\\0\\-2\end{pmatrix}$,
$\overrightarrow{AD}=\begin{pmatrix}0\\1\\0\end{pmatrix}$.
Déterminant :
\[
\det = \begin{vmatrix}
1 & -1 & 0 \\
2 & 0 & 1 \\
3 & -2 & 0
\end{vmatrix}
= 1\begin{vmatrix}0&1\\-2&0\end{vmatrix} - (-1)\begin{vmatrix}2&1\\3&0\end{vmatrix} + 0
= 1(0 - (-2)) + 1(0 - 3) = 2 - 3 = -1.
\]
Volume $\mathcal{V} = \frac{1}{6}|-1| = \frac{1}{6}$.
2. Comme $\mathcal{V}\neq 0$, les points ne sont \textbf{pas coplanaires}.
\end{exoresolu}

\begin{savaistu}[Le produit vectoriel au Cameroun : Topographie et Télécoms]
Dans la conception des réseaux de fibre optique \textbf{CAMTEL} ou l'implantation de pylônes \textbf{MTN/Orange} en zone rurale (ex: région de l'Est ou du Nord), les ingénieurs topographes utilisent le produit vectoriel pour :
\begin{itemize}
  \item Calculer l'aire exacte d'une parcelle de terrain (triangle ou polygone découpé en triangles) à partir de coordonnées GPS (converties en repère local orthonormé).
  \item Déterminer l'orientation précise d'une antenne (vecteur normal au plan de l'antenne) pour optimiser la couverture (azimut, tilt).
  \item Vérifier la verticalité d'un pylône : le vecteur directeur du pylône doit être orthogonal au plan horizontal (produit scalaire nul) et le produit vectoriel de deux haubans donne la normale au plan de stabilisation.
\end{itemize}
Le produit mixte permet de calculer le volume de terrassement nécessaire pour les fondations (tétraèdres ou prismes).
\end{savaistu}

\coursec{Bilan synthétique et préparation au Baccalauréat}

\begin{aretenir}[Formulaire officiel Bac - Produit Vectoriel]
\begin{center}
\begin{tabularx}{\linewidth}{|l|X|}
\hline
\rowcolor{popblueL} \textbf{Notion} & \textbf{Formule / Propriété clé} \\
\hline
Base directe & Règle main droite : Pouce $\vec{\imath}$, Index $\vec{\jmath}$, Majeur $\vec{k}$. \\
\hline
Définition $\vec{u}\wedge\vec{v}$ & $\vec{w}\perp\vec{u},\vec{w}\perp\vec{v}$, $\|\vec{w}\|=\|\vec{u}\|\|\vec{v}\||\sin\theta|$, $(\vec{u},\vec{v},\vec{w})$ direct. \\
\hline
Colinéarité & $\vec{u}\wedge\vec{v}=\vec{0} \Leftrightarrow \vec{u}\parallel\vec{v}$. \\
\hline
Coordonnées (BOND) & $\begin{pmatrix}a\\b\\c\end{pmatrix}\wedge\begin{pmatrix}a'\\b'\\c'\end{pmatrix}=\begin{pmatrix}bc'-cb'\\ca'-ac'\\ab'-ba'\end{pmatrix}$. \\
\hline
Antisymétrie & $\vec{v}\wedge\vec{u}=-(\vec{u}\wedge\vec{v})$. \\
\hline
Bilinéarité & $(\lambda\vec{u}+\mu\vec{v})\wedge\vec{w}=\lambda(\vec{u}\wedge\vec{w})+\mu(\vec{v}\wedge\vec{w})$. \\
\hline
Non-associativité & $(\vec{u}\wedge\vec{v})\wedge\vec{w}\neq\vec{u}\wedge(\vec{v}\wedge\vec{w})$ en général. \\
\hline
Vecteur normal plan & $\vec{n}=\overrightarrow{AB}\wedge\overrightarrow{AC}$. \\
\hline
Distance Point-Droite & $d(M,(D))=\dfrac{\|\overrightarrow{AM}\wedge\vec{u}\|}{\|\vec{u}\|}$, $A\in(D)$, $\vec{u}$ directeur. \\
\hline
Distance Point-Plan & $d(M,\mathcal{P})=\dfrac{|ax_M+by_M+cz_M+d|}{\sqrt{a^2+b^2+c^2}}$, $\mathcal{P}:ax+by+cz+d=0$. \\
\hline
Aire Triangle $ABC$ & $\mathcal{A}=\dfrac{1}{2}\|\overrightarrow{AB}\wedge\overrightarrow{AC}\|$. \\
\hline
Produit Mixte & $[\vec{u},\vec{v},\vec{w}]=(\vec{u}\wedge\vec{v})\cdot\vec{w}=\det(\vec{u},\vec{v},\vec{w})$. \\
\hline
Volume Tétraèdre & $\mathcal{V}=\dfrac{1}{6}|\det(\overrightarrow{AB},\overrightarrow{AC},\overrightarrow{AD})|$. \\
\hline
Coplanarité 4 points & $A,B,C,D$ coplanaires $\Leftrightarrow \det(\overrightarrow{AB},\overrightarrow{AC},\overrightarrow{AD})=0$. \\
\hline
\end{tabularx}
\end{center}
\end{aretenir}

\begin{aretenir}[Checklist finale avant l'épreuve]
\begin{itemize}
  \item[$\square$] Je sais utiliser la main droite pour déterminer si une base est directe ou indirecte.
  \item[$\square$] Je récite la définition géométrique du produit vectoriel (3 conditions : orthogonalité, norme, orientation).
  \item[$\square$] Je calcule les coordonnées sans erreur de signe (vigilance sur la 2\textsuperscript{ème} coordonnée : $ca'-ac'$).
  \item[$\square$] J'utilise l'antisymétrie et la bilinéarité pour simplifier les expressions littérales.
  \item[$\square$] Je ne tombe pas dans le piège de l'associativité (parenthèses obligatoires).
  \item[$\square$] Je trouve un vecteur normal à un plan par $\overrightarrow{AB}\wedge\overrightarrow{AC}$ et j'en déduis l'équation cartésienne.
  \item[$\square$] Je calcule $d(M,(D))$ par $\frac{\|\overrightarrow{AM}\wedge\vec{u}\|}{\|\vec{u}\|}$ (choix judicieux de $A$).
  \item[$\square$] Je calcule $d(M,\mathcal{P})$ par la formule cartésienne (vérification : second membre $=0$).
  \item[$\square$] Je n'oublie \textbf{jamais} le facteur $\frac{1}{2}$ (aire triangle) ni $\frac{1}{6}$ (volume tétraèdre).
  \item[$\square$] Je calcule un déterminant $3\times 3$ (Sarrus ou développement) pour le produit mixte/volume.
  \item[$\square$] Je vérifie la positivité de mes résultats (distance, aire, volume $\ge 0$).
  \item[$\square$] Je vérifie que la base de l'énoncé est bien \textbf{orthonormée directe} avant d'appliquer les formules analytiques.
\end{itemize}
\end{aretenir}

\coursec{Exercices d'entraînement type Bac}

\begin{exercice}[Exercice 1 : Produit vectoriel et plan]
Dans un repère orthonormé direct $(O;\vec{\imath},\vec{\jmath},\vec{k})$, on donne les points $A(1;2;-1)$, $B(3;0;1)$, $C(2;-1;2)$.
\begin{enumerate}
  \item Calculer $\overrightarrow{AB}\wedge\overrightarrow{AC}$. En déduire une équation cartésienne du plan $(ABC)$.
  \item Soit la droite $(D)$ d'équations paramétriques $\begin{cases} x=1+2t \\ y=-1+t \\ z=3-t \end{cases}$. Déterminer l'intersection de $(D)$ avec le plan $(ABC)$.
  \item Calculer la distance du point $M(0;0;0)$ au plan $(ABC)$.
\end{enumerate}
\end{exercice}

\begin{exercice}[Exercice 2 : Distances et tétraèdre]
On considère les points $A(1;1;0)$, $B(2;-1;1)$, $C(0;2;1)$, $D(1;0;3)$.
\begin{enumerate}
  \item Montrer que les droites $(AB)$ et $(CD)$ ne sont pas coplanaires (droites « gauches »).
  \item Calculer la distance entre les droites $(AB)$ et $(CD)$.
  \item Calculer le volume du tétraèdre $ABCD$.
  \item Calculer l'aire du triangle $BCD$ et en déduire la distance de $A$ au plan $(BCD)$.
\end{enumerate}
\end{exercice}

\begin{exercice}[Exercice 3 : Problème concret - Antenne relais]
Une antenne relais est fixée au sommet $S$ d'un pylône vertical de hauteur $30$ m. La base du pylône est au point $O(0;0;0)$. Trois haubans maintiennent le pylône, attachés au sol aux points $A(20;0;0)$, $B(-10;10\sqrt{3};0)$, $C(-10;-10\sqrt{3};0)$.
\begin{enumerate}
  \item Vérifier que $A,B,C$ sont dans un plan horizontal et calculer l'aire du triangle $ABC$.
  \item Déterminer un vecteur normal au plan $(SAB)$.
  \item Calculer la distance du point $C$ au plan $(SAB)$. Interpréter ce résultat physiquement.
  \item Calculer le volume du tétraèdre $SABC$.
\end{enumerate}
\end{exercice}

\coursec{Corrigés détaillés}

\begin{corrige}[Exercice 1]
\begin{enumerate}
  \item $\overrightarrow{AB}=\begin{pmatrix}2\\-2\\2\end{pmatrix}$, $\overrightarrow{AC}=\begin{pmatrix}1\\-3\\3\end{pmatrix}$.
  $\overrightarrow{AB}\wedge\overrightarrow{AC} = \begin{pmatrix} (-2)(3)-(2)(-3) \\ (2)(1)-(2)(3) \\ (2)(-3)-(-2)(1) \end{pmatrix} = \begin{pmatrix} -6+6 \\ 2-6 \\ -6+2 \end{pmatrix} = \begin{pmatrix}0\\-4\\-4\end{pmatrix}$.
  Vecteur normal simplifié $\vec{n}=\begin{pmatrix}0\\1\\1\end{pmatrix}$. Équation : $y+z+d=0$. Point $A(1;2;-1)$ : $2-1+d=0 \Rightarrow d=-1$.
  Plan $(ABC) : y+z-1=0$.
  \item $(D) : \vec{u}=\begin{pmatrix}2\\1\\-1\end{pmatrix}$, $M_0(1;-1;3)$. Intersection $I$ : $\exists t, I(1+2t; -1+t; 3-t) \in (ABC)$.
  $(-1+t) + (3-t) - 1 = 0 \Rightarrow 1 = 0$ \textbf{Impossible}.
  La droite $(D)$ est \textbf{parallèle} au plan $(ABC)$ (car $\vec{u}\cdot\vec{n} = 2(0)+1(1)+(-1)(1)=0$) et ne le coupe pas.
  \item $d(O,(ABC)) = \frac{|0+0-1|}{\sqrt{0^2+1^2+1^2}} = \frac{1}{\sqrt{2}} = \frac{\sqrt{2}}{2}$.
\end{enumerate}
\end{corrige}

\begin{corrige}[Exercice 2]
\begin{enumerate}
  \item $\overrightarrow{AB}=\begin{pmatrix}1\\-2\\1\end{pmatrix}$, $\overrightarrow{CD}=\begin{pmatrix}1\\-2\\2\end{pmatrix}$. $\overrightarrow{AB}$ et $\overrightarrow{CD}$ non colinéaires.
  Vecteur $\overrightarrow{AC}=\begin{pmatrix}-1\\1\\1\end{pmatrix}$.
  Déterminant $\det(\overrightarrow{AB},\overrightarrow{CD},\overrightarrow{AC}) = \begin{vmatrix}1&1&-1\\-2&-2&1\\1&2&1\end{vmatrix} = 1(-2-2) -1(-2-1) -1(-4+2) = -4+3+2=1\neq 0$.
  Droites gauches (non coplanaires). \quad $\checkmark$
  \item Distance entre droites gauches $(AB)$ et $(CD)$ :
  $d = \frac{|[\overrightarrow{AB},\overrightarrow{CD},\overrightarrow{AC}]|}{\|\overrightarrow{AB}\wedge\overrightarrow{CD}\|}$.
  Numérateur $= |1| = 1$.
  $\overrightarrow{AB}\wedge\overrightarrow{CD} = \begin{pmatrix} (-2)(2)-(1)(-2) \\ (1)(1)-(1)(2) \\ (1)(-2)-(-2)(1) \end{pmatrix} = \begin{pmatrix} -4+2 \\ 1-2 \\ -2+2 \end{pmatrix} = \begin{pmatrix}-2\\-1\\0\end{pmatrix}$.
  Norme $= \sqrt{4+1} = \sqrt{5}$.
  $d = \frac{1}{\sqrt{5}} = \frac{\sqrt{5}}{5}$.
  \item Volume $ABCD = \frac{1}{6}|\det| = \frac{1}{6}$.
  \item Aire $BCD$ : $\overrightarrow{BC}=\begin{pmatrix}-2\\3\\0\end{pmatrix}$, $\overrightarrow{BD}=\begin{pmatrix}-1\\1\\2\end{pmatrix}$.
  $\overrightarrow{BC}\wedge\overrightarrow{BD} = \begin{pmatrix} 6-0 \\ 0-(-4) \\ -2-(-3) \end{pmatrix} = \begin{pmatrix}6\\4\\1\end{pmatrix}$. Norme $= \sqrt{36+16+1}=\sqrt{53}$.
  $\mathcal{A}_{BCD} = \frac{\sqrt{53}}{2}$.
  Distance $A$ au plan $(BCD)$ : $h = \frac{3\mathcal{V}}{\mathcal{A}_{BCD}} = \frac{3(1/6)}{\sqrt{53}/2} = \frac{1/2}{\sqrt{53}/2} = \frac{1}{\sqrt{53}}$.
\end{enumerate}
\end{corrige}

\begin{corrige}[Exercice 3 : Antenne relais]
\begin{enumerate}
  \item $A,B,C$ ont $z=0$ $\Rightarrow$ plan horizontal $z=0$ (sol).
  $\overrightarrow{AB}=\begin{pmatrix}-30\\10\sqrt{3}\\0\end{pmatrix}$, $\overrightarrow{AC}=\begin{pmatrix}-30\\-10\sqrt{3}\\0\end{pmatrix}$.
  $\overrightarrow{AB}\wedge\overrightarrow{AC} = \begin{pmatrix}0\\0\\ (-30)(-10\sqrt{3}) - (10\sqrt{3})(-30) \end{pmatrix} = \begin{pmatrix}0\\0\\ 300\sqrt{3}+300\sqrt{3} \end{pmatrix} = \begin{pmatrix}0\\0\\600\sqrt{3}\end{pmatrix}$.
  Norme $= 600\sqrt{3}$. Aire $= \frac{1}{2}\times 600\sqrt{3} = 300\sqrt{3} \approx 519,6 \text{ m}^2$.
  \item $S(0;0;30)$. $\overrightarrow{SA}=\begin{pmatrix}20\\0\\-30\end{pmatrix}$, $\overrightarrow{SB}=\begin{pmatrix}-10\\10\sqrt{3}\\-30\end{pmatrix}$.
  $\vec{n}_{SAB} = \overrightarrow{SA}\wedge\overrightarrow{SB} = \begin{pmatrix} 0(-30)-(-30)(10\sqrt{3}) \\ (-30)(-10)-(20)(-30) \\ 20(10\sqrt{3})-0(-10) \end{pmatrix} = \begin{pmatrix} 300\sqrt{3} \\ 300+600 \\ 200\sqrt{3} \end{pmatrix} = \begin{pmatrix} 300\sqrt{3} \\ 900 \\ 200\sqrt{3} \end{pmatrix}$.
  Simplifié par $100$ : $\vec{n}=\begin{pmatrix}3\sqrt{3}\\9\\2\sqrt{3}\end{pmatrix}$.
  \item Plan $(SAB)$ : $3\sqrt{3}x + 9y + 2\sqrt{3}z + d = 0$. Point $S(0;0;30)$ : $0+0+60\sqrt{3}+d=0 \Rightarrow d=-60\sqrt{3}$.
  Équation : $3\sqrt{3}x + 9y + 2\sqrt{3}z - 60\sqrt{3} = 0$.
  Distance $C(-10;-10\sqrt{3};0)$ :
  Numérateur $|3\sqrt{3}(-10) + 9(-10\sqrt{3}) + 0 - 60\sqrt{3}| = |-30\sqrt{3} - 90\sqrt{3} - 60\sqrt{3}| = 180\sqrt{3}$.
  Dénominateur $\sqrt{(3\sqrt{3})^2 + 9^2 + (2\sqrt{3})^2} = \sqrt{27+81+12} = \sqrt{120} = 2\sqrt{30}$.
  $d(C,(SAB)) = \frac{180\sqrt{3}}{2\sqrt{30}} = \frac{90\sqrt{3}}{\sqrt{30}} = 90\sqrt{\frac{1}{10}} = \frac{90}{\sqrt{10}} = 9\sqrt{10} \approx 28,5 \text{ m}$.
  \emph{Interprétation :} C'est la longueur du segment perpendiculaire depuis l'ancrage $C$ jusqu'au plan formé par le pylône et les haubans $SA, SB$. Cela représente la « marge » spatiale.
  \item Volume $SABC = \frac{1}{6}|\det(\overrightarrow{SA},\overrightarrow{SB},\overrightarrow{SC})|$.
  $\overrightarrow{SC}=\begin{pmatrix}-10\\-10\sqrt{3}\\-30\end{pmatrix}$.
  Déterminant (en factorisant $10$ sur chaque colonne) :
  $\det = 10^3 \begin{vmatrix} 2 & -1 & -1 \\ 0 & \sqrt{3} & -\sqrt{3} \\ -3 & -3 & -3 \end{vmatrix} = 1000 \times (-3) \begin{vmatrix} 2 & -1 & -1 \\ 0 & \sqrt{3} & -\sqrt{3} \\ 1 & 1 & 1 \end{vmatrix}$.
  Calcul : $-3000 [ 2(\sqrt{3}+\sqrt{3}) - (-1)(0+\sqrt{3}) + (-1)(0-\sqrt{3}) ] = -3000 [ 4\sqrt{3} + \sqrt{3} + \sqrt{3} ] = -3000 \times 6\sqrt{3} = -18000\sqrt{3}$.
  Volume $= \frac{1}{6} \times 18000\sqrt{3} = 3000\sqrt{3} \approx 5196 \text{ m}^3$.
\end{enumerate}
\end{corrige}

\findecours

\end{document}
